%% toc: P15 169 169.1,169.2,169.3,169.4
\section{Proofs for likelihood estimation under stated regularity}
\label{sec:p15-w02b-p15-057-annex}

Throughout, \(E_\theta\) is expectation under \(P_\theta\), and
\(X\) denotes one observation with that law.

\begin{annexitem}{annex:p15-w02b-p15-057.halmos-savage}{thm:p15-w02b-p15-057-halmos-savage}{A countable mixture dominates a dominated family}
\begin{proof}
\emph{An equivalent probability.} Write \(\mathsf X\) as a disjoint union
of sets \(E_j\) with \(\mu(E_j)<\infty\), and let \(J\) be the set of
\(j\) with \(\mu(E_j)>0\). Put
\(\mu'(B)=\sum_{j\in J}2^{-j}\mu(B\cap E_j)/\mu(E_j)\), normalized by the
positive finite sum \(\sum_{j\in J}2^{-j}\). Then \(\mu'\) is a probability
measure and \(\mu'(B)=0\) exactly when \(\mu(B)=0\).

\emph{The supremum of supports.} Let \(\mathcal K\) be the set of
measures \(Q=\sum_ic_iP_{\theta_i}\) with countably many indices,
\(c_i>0\) and \(\sum_ic_i=1\). Each \(Q\in\mathcal K\) is a probability
with \(Q\ll\mu\); fix its density \(q=\sum_ic_if_{\theta_i}\) and put
\(S_Q=\{q>0\}\). Let \(s=\sup_{Q\in\mathcal K}\mu'(S_Q)\leq1\) and choose
\(Q_k\in\mathcal K\) with \(\mu'(S_{Q_k})\to s\). The measure
\(Q^*=\sum_k2^{-k}Q_k\) is again in \(\mathcal K\), being a countable
positive combination of the countably many \(P_{\theta_i}\) involved, with
total weight one. Its density \(\sum_k2^{-k}q_k\) is positive exactly on
\(\bigcup_kS_{Q_k}\), so \(\mu'(S_{Q^*})\geq\mu'(S_{Q_k})\) for each
\(k\) and therefore \(\mu'(S_{Q^*})=s\).

\emph{Domination.} Fix \(\theta\) and a set \(N\) with \(Q^*(N)=0\).
Since \(Q^*(N)=\int_Nq^*\,d\mu\) and \(q^*>0\) on \(S_{Q^*}\),
\(\mu(N\cap S_{Q^*})=0\), so \(P_\theta(N\cap S_{Q^*})=0\). Suppose
\(P_\theta(N\setminus S_{Q^*})>0\). Then
\(\mu(\{f_\theta>0\}\setminus S_{Q^*})>0\), hence
\(\mu'(\{f_\theta>0\}\setminus S_{Q^*})>0\). The measure
\(\tfrac12Q^*+\tfrac12P_\theta\) belongs to \(\mathcal K\) and its
density is positive on \(S_{Q^*}\cup\{f_\theta>0\}\), whose
\(\mu'\)-measure exceeds \(s\); this contradicts the definition of
\(s\). Hence \(P_\theta(N)=0\), and \(\lambda=Q^*\) has the required form.
\end{proof}
\end{annexitem}

\begin{annexitem}{annex:p15-w02b-p15-057.factorization}{thm:p15-w02b-p15-057-factorization}{The factorization criterion}
\begin{proof}
Let \(\lambda=\sum_ic_iP_{\theta_i}\) be the measure of
\cref{thm:p15-w02b-p15-057-halmos-savage}. Two facts are used. First,
a \(\sigma(T)\)-measurable \(W:\mathsf X\to[0,\infty]\) has the form
\(w\circ T\) with \(w\) \(\mathscr T\)-measurable: for an indicator
\(\mathbf1_{T^{-1}(D)}=\mathbf1_D\circ T\); hence for simple functions by
linearity; and in general \(W\) is the increasing limit of simple
\(\sigma(T)\)-measurable \(W_k=w_k\circ T\), so \(W=w\circ T\) with
\(w=\sup_kw_k\). Second, for a probability \(\nu\), a set
\(B\in\mathscr X\) and \(\phi=E_\nu[\mathbf1_B\mid\sigma(T)]\) as in
\cref{thm:p15-conditional-existence}, the identity
\(\int\mathbf1_BW\,d\nu=\int\phi W\,d\nu\) holds for every
\(\sigma(T)\)-measurable \(W\geq0\): it is the defining identity for
indicators \(W=\mathbf1_C\), extends to simple \(W\) by linearity and to
general \(W\) by monotone convergence.

\emph{Factorization implies sufficiency.} Suppose
\(f_\theta=g_\theta(T)h\) \(\mu\)-almost everywhere. The density of
\(\lambda\) is \(k(T)h\) with \(k=\sum_ic_ig_{\theta_i}\), by monotone
convergence; it is finite \(\mu\)-almost everywhere because \(\lambda\) is
a probability. Define \(r_\theta=g_\theta(T)/k(T)\) on
\(G=\{0<k(T)<\infty\}\) and \(r_\theta=0\) off \(G\); it is
\(\sigma(T)\)-measurable. For \(B\in\mathscr X\),
\(\int_Br_\theta\,d\lambda=\int_{B\cap G}g_\theta(T)h\,d\mu\). The
remaining part of \(P_\theta(B)=\int_Bg_\theta(T)h\,d\mu\) lies in
\(\{h=0\}\), where the integrand vanishes; in
\(\{k(T)=0,h>0\}\), which is \(\lambda\)-null and therefore
\(P_\theta\)-null; or in \(\{k(T)=\infty,h>0\}\), which is \(\mu\)-null.
Hence \(r_\theta=dP_\theta/d\lambda\). Let
\(\phi_B=E_\lambda[\mathbf1_B\mid\sigma(T)]\), chosen with values in
\([0,1]\). For \(C\in\sigma(T)\), applying the second fact with
\(W=\mathbf1_Cr_\theta\),
\[
 P_\theta(B\cap C)=\int\mathbf1_B\mathbf1_Cr_\theta\,d\lambda
 =\int\phi_B\mathbf1_Cr_\theta\,d\lambda=\int_C\phi_B\,dP_\theta .
\]
Thus \(\phi_B\) serves for every \(\theta\), and \(T\) is sufficient.

\emph{Sufficiency implies factorization.} Let \(\phi_B\) be the common
versions. Summing \(P_{\theta_i}(B\cap C)=\int_C\phi_B\,dP_{\theta_i}\)
with weights \(c_i\) gives \(\lambda(B\cap C)=\int_C\phi_B\,d\lambda\), so
\(\phi_B=E_\lambda[\mathbf1_B\mid\sigma(T)]\) \(\lambda\)-almost surely.
Let \(r_\theta=dP_\theta/d\lambda\geq0\), which is \(\lambda\)-integrable
with integral one, and \(\rho_\theta=E_\lambda[r_\theta\mid\sigma(T)]\),
chosen nonnegative. For \(B\in\mathscr X\), taking \(C=\mathsf X\),
\[
 P_\theta(B)=\int\phi_B\,dP_\theta=\int\phi_Br_\theta\,d\lambda
 =\int\phi_B\rho_\theta\,d\lambda=\int\mathbf1_B\rho_\theta\,d\lambda,
\]
the third equality because \(\phi_B\) is bounded and
\(\sigma(T)\)-measurable, the fourth by the second fact with
\(\nu=\lambda\) and \(W=\rho_\theta\). Hence
\(\rho_\theta=dP_\theta/d\lambda\), and by the first fact
\(\rho_\theta=g_\theta(T)\) with \(g_\theta\geq0\)
\(\mathscr T\)-measurable. Let \(h=d\lambda/d\mu\). For
\(B\in\mathscr X\), \(\int_B\rho_\theta\,d\lambda=\int_B\rho_\theta h\,d\mu\),
first for simple \(\rho_\theta\) and then by monotone convergence. So
\(dP_\theta/d\mu=g_\theta(T)h\) \(\mu\)-almost everywhere.
\end{proof}
\end{annexitem}

\begin{annexitem}{annex:p15-w02b-p15-057.wald}{thm:p15-w02b-p15-057-wald}{Strong consistency on a compact parameter space}
\begin{proof}
\emph{Existence.} Each \(\theta\mapsto\log f_\theta(X_i)\) is continuous
into \([-\infty,\infty)\) by (W2), so their finite sum \(\ell_n\) is
continuous into \([-\infty,\infty)\). Its image of the compact \(\Theta\)
is compact, so the supremum is attained. Since
\(P_{\theta_0}(f_{\theta_0}(X)>0)=1\), the event
\(\{\ell_n(\theta_0)>-\infty\text{ for all }n\}\) has probability one.

\emph{Negative expected log ratio.} On \(\{f_{\theta_0}>0\}\) put
\(r_\theta=\log f_\theta-\log f_{\theta_0}\in[-\infty,\infty)\). By (W4),
\(E_{\theta_0}r_\theta^+<\infty\), so \(E_{\theta_0}r_\theta\in[-\infty,
\infty)\) is defined. Since \(\log t\leq t-1\) for \(t\in[0,\infty)\),
with \(\log0=-\infty\) and equality only at \(t=1\),
\[
 E_{\theta_0}r_\theta\leq E_{\theta_0}\bigl(e^{r_\theta}-1\bigr)
 =\int_{\{f_{\theta_0}>0\}}f_\theta\,d\mu-1\leq0 .
\]
If \(E_{\theta_0}r_\theta=0\), both inequalities are equalities. The first
forces \(e^{r_\theta}=1\) \(P_{\theta_0}\)-almost surely, so
\(f_\theta=f_{\theta_0}\) \(\mu\)-almost everywhere on
\(\{f_{\theta_0}>0\}\). The second forces
\(P_\theta(f_{\theta_0}>0)=1\). Then for every \(B\),
\(P_\theta(B)=\int_{B\cap\{f_{\theta_0}>0\}}f_{\theta_0}\,d\mu
=P_{\theta_0}(B)\), contradicting (W3) when \(\theta\neq\theta_0\).
Hence \(E_{\theta_0}r_\theta<0\) for \(\theta\neq\theta_0\).

\emph{Local upper envelopes.} For \(\theta\neq\theta_0\) and \(\rho>0\)
let \(G_\rho=\sup\{r_{\theta'}:d(\theta',\theta)<\rho\}\). A compact
metric space is separable, and for a continuous
\([-\infty,\infty)\)-valued function the supremum over an open ball equals
the supremum over a countable dense subset of it; so \(G_\rho\) is
measurable. By (W4), \(G_\rho^+\) is integrable. As \(\rho\downarrow0\)
along a sequence, \(G_\rho\downarrow r_\theta\) pointwise by (W2). Then
\(G_\rho^+\downarrow r_\theta^+\) with an integrable bound, so
\(E G_\rho^+\to Er_\theta^+\) by dominated convergence, and
\(G_\rho^-\uparrow r_\theta^-\), so \(EG_\rho^-\to Er_\theta^-\) by monotone
convergence. Hence \(EG_\rho\to Er_\theta<0\), and some \(\rho_\theta\)
has \(EG_{\rho_\theta}<0\). The functions
\(\max(G_{\rho_\theta},-N)\) decrease to \(G_{\rho_\theta}\) as
\(N\to\infty\), with integrable positive parts, so by the same two
convergence theorems some \(N_\theta\) gives
\(H_\theta=\max(G_{\rho_\theta},-N_\theta)\) with
\(-N_\theta\leq H_\theta\leq G_{\rho_\theta}^+\) and \(EH_\theta<0\).

\emph{Cover and strong law.} Fix \(\varepsilon>0\). The set
\(K_\varepsilon=\{\theta:d(\theta,\theta_0)\geq\varepsilon\}\) is compact
and covered by the open balls of radius \(\rho_\theta\), so by finitely
many of them, with centres \(\vartheta_1,\ldots,\vartheta_J\). Each
\(H_{\vartheta_j}(X)\) is integrable, so by
\cref{thm:p15-iid-series-strong-law}
\(n^{-1}\sum_{i\leq n}H_{\vartheta_j}(X_i)\to EH_{\vartheta_j}<0\) almost
surely. On the intersection of these \(J\) events with
\(\{f_{\theta_0}(X_i)>0\ \forall i\}\) there is \(n_\varepsilon\) such that
for \(n\geq n_\varepsilon\) and every \(\theta\in K_\varepsilon\), lying in
the \(j\)-th ball,
\[
 \frac1n\bigl(\ell_n(\theta)-\ell_n(\theta_0)\bigr)
 =\frac1n\sum_{i\leq n}r_\theta(X_i)
 \leq\frac1n\sum_{i\leq n}H_{\vartheta_j}(X_i)<0 .
\]
A maximizer satisfies \(\ell_n(\widehat\theta_n)\geq\ell_n(\theta_0)\),
so \(\widehat\theta_n\notin K_\varepsilon\) for \(n\geq n_\varepsilon\),
whichever maximizer is chosen. Intersecting the probability-one events for
\(\varepsilon=1/k\), \(k=1,2,\ldots\), gives one probability-one event on
which every choice of maximizers converges to \(\theta_0\).
\end{proof}
\end{annexitem}

\begin{annexitem}{annex:p15-w02b-p15-057.wald-noncompact}{thm:p15-w02b-p15-057-wald-noncompact}{Strong consistency on a noncompact parameter space}
\begin{proof}
Write \(f_0=f_{\theta_0}\) and \(E=E_{\theta_0}\). The event
\(\Omega_0=\{f_0(X_i)>0\text{ for all }i\}\) has probability one, and on it
\(\ell_n(\theta_0)\) is finite for every \(n\). Under (C) extend \(\ell_n\)
to \(\overline\Theta\) by \(\ell_n(c)=\sum_{i\leq n}\log\overline f_c(X_i)\);
under (K) put \(\overline\Theta=\Theta\) and \(\overline f=f\). For
\(x\) with \(f_0(x)>0\) let
\(r_c(x)=\log\overline f_c(x)-\log f_0(x)\in[-\infty,\infty)\).

\emph{Step 1: negative expected log ratio.} Let \(c\) be a point with
\(\int\overline f_c\,d\mu\leq1\), with \(\overline f_c\neq f_0\) on a set of
positive \(\mu\)-measure, and with \(Er_c^+<\infty\). Then
\(Er_c\in[-\infty,\infty)\) is defined, and \(Er_c<0\). Indeed
\(\log t\leq t-1\) for \(t\in[0,\infty)\), with equality only at \(t=1\), so
\[
 Er_c\leq E\bigl(e^{r_c}-1\bigr)=\int_{\{f_0>0\}}\overline f_c\,d\mu-1\leq0 .
\]
If \(Er_c=0\), both inequalities are equalities. The first forces
\(e^{r_c}=1\) almost surely, that is \(\overline f_c=f_0\) \(\mu\)-almost
everywhere on \(\{f_0>0\}\); the second, with
\(\int\overline f_c\,d\mu\leq1\), forces \(\int_{\{f_0=0\}}\overline
f_c\,d\mu\leq0\), so \(\overline f_c=0=f_0\) \(\mu\)-almost everywhere on
\(\{f_0=0\}\). Then \(\overline f_c=f_0\) \(\mu\)-almost everywhere, a
contradiction. Every \(\theta\in\Theta\setminus\{\theta_0\}\) is such a
point when \(Er_\theta^+<\infty\): \(\int f_\theta\,d\mu=1\), and by (W3)
\(P_\theta\neq P_{\theta_0}\), so the two densities differ on a set of
positive \(\mu\)-measure. Under (C) the points of
\(\overline\Theta\setminus\Theta\) are such points by hypothesis.

\emph{Step 2: uniform negative drift on a compact set.} Let
\(C\subseteq\overline\Theta\) be compact with \(\theta_0\notin C\), every
\(c\in C\) as in Step 1, and \(E\sup_{c\in C}r_c^+<\infty\). We show that
there are \(\eta>0\) and a probability-one event on which
\(\sup_{c\in C}n^{-1}\sum_{i\leq n}r_c(X_i)\leq-\eta\) for all large \(n\).
If \(C\) is empty there is nothing to prove. For \(c\in C\) and
\(\rho>0\) let \(G_{c,\rho}=\sup\{r_{c'}:c'\in C,\ d(c',c)<\rho\}\). The
compact metric space \(C\) is separable; a relatively open subset of it
contains a countable dense subset, and since \(c'\mapsto r_{c'}(x)\) is
continuous into \([-\infty,\infty)\) by (W2) or by (C), the supremum over
the subset equals the supremum over that countable set, so \(G_{c,\rho}\)
is measurable. Its positive part is at most
\(\sup_{c'\in C}r_{c'}^+\), which is integrable. As \(\rho\downarrow0\)
along a sequence, \(G_{c,\rho}\downarrow r_c\) pointwise, by continuity.
Then \(G_{c,\rho}^+\downarrow r_c^+\) under an integrable bound, so
\(EG_{c,\rho}^+\to Er_c^+\) by dominated convergence, and
\(G_{c,\rho}^-\uparrow r_c^-\), so \(EG_{c,\rho}^-\to Er_c^-\) by monotone
convergence. Hence \(EG_{c,\rho}\to Er_c<0\), and some \(\rho_c\) has
\(EG_{c,\rho_c}<0\). By the same two convergence theorems applied to
\(\max(G_{c,\rho_c},-N)\downarrow G_{c,\rho_c}\) as \(N\to\infty\), some
\(N_c\) gives \(H_c=\max(G_{c,\rho_c},-N_c)\), integrable, with
\(EH_c<0\) and \(H_c\geq G_{c,\rho_c}\). Finitely many balls of radii
\(\rho_{c_1},\ldots,\rho_{c_J}\) around \(c_1,\ldots,c_J\) cover \(C\). By
\cref{thm:p15-iid-series-strong-law},
\(n^{-1}\sum_{i\leq n}H_{c_j}(X_i)\to EH_{c_j}\) almost surely for each
\(j\). Put \(\eta=-\tfrac12\max_jEH_{c_j}>0\). On the intersection of these
\(J\) events with \(\Omega_0\), for all large \(n\) and every \(c\in C\),
lying in the \(j\)-th ball,
\[
 \frac1n\sum_{i\leq n}r_c(X_i)\leq\frac1n\sum_{i\leq n}H_{c_j}(X_i)
 \leq-\eta .
\]

\emph{Step 3: the region outside \(K\) under (K).} The set
\(\Theta\setminus K\) is open in \(\Theta\), because a compact subset of a
metric space is closed. As \(\Theta\) is separable, \(\Theta\setminus K\)
has a countable dense subset, and by (W2) the supremum defining \(G\) may
be taken over it, so \(G\) is measurable. By hypothesis \(EG^+<\infty\) and
\(EG<0\), so, as in Step 2, some \(N\) gives an integrable
\(H=\max(G,-N)\geq G\) with \(EH<0\), and on a probability-one event
\(\sup_{\theta\in\Theta\setminus K}n^{-1}\sum_{i\leq n}r_\theta(X_i)\leq
n^{-1}\sum_{i\leq n}H(X_i)\leq\tfrac12EH\) for all large \(n\).

\emph{Step 4: consistency.} For \(k\geq1\) let
\(B_k=\{c\in\overline\Theta:d(c,\theta_0)<1/k\}\). Under (C) let
\(C_k=\overline\Theta\setminus B_k\), a closed subset of a compact space,
hence compact; \(E\sup_{C_k}r^+\leq E\sup_{\overline\Theta}r^+<\infty\),
and Step 1 applies to each of its points. Under (K) let
\(C_k=K\setminus B_k\), compact, with
\(E\sup_{C_k}r^+\leq E\sup_Kr^+<\infty\), and every point of \(C_k\) is a
point of \(\Theta\setminus\{\theta_0\}\) to which Step 1 applies. Step 2
gives \(\eta_k>0\) and a probability-one event \(\Omega_k\). Under (K),
\(\Theta\setminus B_k\subseteq C_k\cup(\Theta\setminus K)\); intersect
\(\Omega_k\) with the event of Step 3 and replace \(\eta_k\) by
\(\min(\eta_k,-\tfrac12EH)\). In either case, on \(\Omega_0\cap\Omega_k\)
there is \(n_k\) with
\[
 \ell_n(\theta)-\ell_n(\theta_0)\leq-n\eta_k\qquad
 (n\geq n_k,\ \theta\in\Theta\setminus B_k).
\]
Let \(\widehat\theta_n\) satisfy
\(\ell_n(\widehat\theta_n)\geq\ell_n(\theta_0)-n\varepsilon_n\) for all
large \(n\). Once \(\varepsilon_n<\eta_k\) and \(n\geq n_k\), the point
\(\widehat\theta_n\) cannot lie in \(\Theta\setminus B_k\), so
\(d(\widehat\theta_n,\theta_0)<1/k\). On the probability-one event
\(\Omega_0\cap\bigcap_k\Omega_k\), which does not depend on the sequence,
this holds for every \(k\), so \(\widehat\theta_n\to\theta_0\). A maximum
likelihood estimate satisfies the inequality with \(\varepsilon_n=0\).

\emph{Step 5: existence.} On \(\Omega_0\), \(\ell_n\) is a finite sum of
functions continuous into \([-\infty,\infty)\), hence continuous into
\([-\infty,\infty)\), and attains its supremum on every compact set. Under
(K) with \(\theta_0\in K\): by Step 3, for large \(n\),
\(\sup_{\Theta\setminus K}\ell_n\leq\ell_n(\theta_0)+\tfrac12nEH
<\ell_n(\theta_0)\). A maximizer \(\theta^*\) of \(\ell_n\) on \(K\) has
\(\ell_n(\theta^*)\geq\ell_n(\theta_0)\), so it maximizes \(\ell_n\) on
\(\Theta\). Under (C) with \(\Theta\) open in \(\overline\Theta\): choose
\(\rho>0\) with \(\{c:d(c,\theta_0)<2\rho\}\subseteq\Theta\), and let
\(F=\{c:d(c,\theta_0)\leq\rho\}\subseteq\Theta\), compact. Take \(k\) with
\(1/k\leq\rho\); then \(\overline\Theta\setminus F\subseteq C_k\), so for
\(n\geq n_k\), \(\sup_{\overline\Theta\setminus F}\ell_n<\ell_n(\theta_0)\).
A maximizer \(\theta^*\in F\) of \(\ell_n\) on \(F\) satisfies
\(\ell_n(\theta^*)\geq\ell_n(\theta_0)\), so it maximizes \(\ell_n\) on
\(\overline\Theta\), and in particular on \(\Theta\).
\end{proof}
\end{annexitem}

\begin{annexitem}{annex:p15-w02b-p15-057.cauchy-line}{thm:p15-w02b-p15-057-cauchy-line}{Cauchy location on the whole line}
\begin{proof}
\emph{The compactification.} The map \(\arctan\) is a homeomorphism of
\(\mathbb R\) onto \((-\pi/2,\pi/2)\) and, extended by
\(\arctan(\pm\infty)=\pm\pi/2\), an isometry of
\((\overline\Theta,d)\) onto \([-\pi/2,\pi/2]\). Hence \(\overline\Theta\)
is a compact metric space, \(d\) induces the usual topology on
\(\mathbb R\), and \(\mathbb R\) is open in \(\overline\Theta\). For fixed
\(x\), \(\theta\mapsto f_\theta(x)\) is continuous on \(\mathbb R\) and
tends to \(0\) as \(\theta\to\pm\infty\), so with
\(\overline f_{\pm\infty}=0\) it is continuous on \(\overline\Theta\). At the
two added points \(\int\overline f\,d\mu=0\leq1\), and
\(\overline f=0\neq f_{\theta_0}\) everywhere, since \(f_{\theta_0}>0\).
(W2) holds, and (W3) holds because \(f_\theta\) has its unique maximum at
\(x=\theta\), so distinct parameters give distinct continuous densities,
which differ on an interval.

\emph{The envelope.} For \(\theta\in\mathbb R\),
\(\log f_\theta(x)-\log f_{\theta_0}(x)=\log(1+(x-\theta_0)^2)
-\log(1+(x-\theta)^2)\), which is at most \(\log(1+(x-\theta_0)^2)\), with
equality at \(\theta=x\); at \(\pm\infty\) the difference is \(-\infty\). So
the supremum in (C) is \(\log(1+(x-\theta_0)^2)\geq0\), and under
\(P_{\theta_0}\) the variable \(Z=X-\theta_0\) is standard Cauchy. Since
\(1+z^2\leq(1+|z|)^2\), \(\log(1+Z^2)\leq2\log(1+|Z|)\). For \(t\geq0\),
\(\log(1+t)\leq\sqrt t\): both vanish at \(t=0\), and
\(1/(1+t)\leq1/(2\sqrt t)\) for \(t>0\) because \(2\sqrt t\leq1+t\). Hence
\[
 E\log(1+Z^2)\leq2E|Z|^{1/2}
 =\frac4\pi\int_0^\infty\frac{z^{1/2}}{1+z^2}\,dz
 \leq\frac4\pi\Bigl(\int_0^1dz+\int_1^\infty z^{-3/2}\,dz\Bigr)
 =\frac{12}\pi<\infty .
\]
So (C) holds.

\emph{Existence and consistency.} For every sample, \(\ell_n\) is
continuous on \(\mathbb R\) and tends to \(-\infty\) as
\(\theta\to\pm\infty\), because each term \(\log f_\theta(X_i)\) does; so
it attains its maximum. Strong consistency of every choice is
\cref{thm:p15-w02b-p15-057-wald-noncompact} under (C).

\emph{The limit law.} The verification of (R1)--(R5) with \(M\equiv2\) and
\(I(\theta_0)=\tfrac12\) in the proof of
\cref{thm:p15-w02b-p15-057-cauchy-check} uses only the densities near
\(\theta_0\), so the family is regular at \(\theta_0\). A maximizer over
the open set \(\mathbb R\) of the differentiable function \(\ell_n\) is a
root of \(\ell_n'\), so \(P(\ell_n'(\widehat\theta_n)=0)=1\), and almost
sure convergence gives convergence in probability. For a measurable choice,
\cref{thm:p15-w02b-p15-057-asymptotic} gives
\(\sqrt n(\widehat\theta_n-\theta_0)\Rightarrow N(0,2)\).
\end{proof}
\end{annexitem}

\begin{annexitem}{annex:p15-w02b-p15-057.asymptotic}{thm:p15-w02b-p15-057-asymptotic}{A consistent root of the likelihood equation}
\begin{proof}
\emph{Score moments.} Let \(A\) and \(N\) be as in (R1). Since
\(P_{\theta_0}(A)=1\), almost surely all \(X_i\in A\), and then by (R2)
\(\ell_n\) is twice continuously differentiable on \(N\). Measurability
of \(s_\theta(x)\) and \(\partial_\theta^2\log f_\theta(x)\) in \(x\)
follows because they are pointwise limits of difference quotients of
measurable functions. On \(A\), \(s_{\theta}f_{\theta}=\partial_\theta
f_\theta\) and \(\partial_\theta^2\log f_\theta=\partial_\theta^2f_\theta
/f_\theta-s_\theta^2\). By (R3),
\(E_{\theta_0}s_{\theta_0}=\int_A\partial_\theta f_{\theta_0}\,d\mu=0\) and
\(E_{\theta_0}[\partial_\theta^2f_{\theta_0}/f_{\theta_0}]=\int_A
\partial_\theta^2f_{\theta_0}\,d\mu=0\), whence
\(E_{\theta_0}\partial_\theta^2\log f_{\theta_0}=-I(\theta_0)\); the
expectation exists by (R5).

\emph{Uniform control of the curvature.} For \(0<\eta<\delta\) put
\[
 h_\eta(x)=\sup_{|\theta-\theta_0|\leq\eta}
 \bigl|\partial_\theta^2\log f_\theta(x)-\partial_\theta^2\log
 f_{\theta_0}(x)\bigr|\qquad(x\in A),
\]
measurable as a supremum over the rationals in the interval and its
endpoints, by continuity. By (R5), \(0\leq h_\eta\leq2M\), and by (R2),
\(h_\eta(x)\to0\) as \(\eta\to0\). Dominated convergence gives
\(E_{\theta_0}h_\eta\to0\).

\emph{Expansion.} Let \(\Delta_n=\widehat\theta_n-\theta_0\) and
\[
 B_n=\int_0^1\ell_n''(\theta_0+t\Delta_n)\,dt
\]
on the event \(\Omega_n\) where all \(X_i\in A\), \(|\Delta_n|<\eta\) and
\(\ell_n'(\widehat\theta_n)=0\), and \(B_n=\ell_n''(\theta_0)\) elsewhere;
\(B_n\) is measurable as a limit of Riemann sums of a continuous integrand.
By the fundamental theorem of calculus, on \(\Omega_n\),
\(0=\ell_n'(\widehat\theta_n)=\ell_n'(\theta_0)+\Delta_nB_n\). The
hypotheses give \(P_{\theta_0}(\Omega_n)\to1\) for each fixed \(\eta\).
On \(\Omega_n\),
\[
 \Bigl|\frac{B_n}{n}-\frac{\ell_n''(\theta_0)}{n}\Bigr|
 \leq\frac1n\sum_{i\leq n}h_\eta(X_i),
\]
and the right side converges almost surely to \(E_{\theta_0}h_\eta\) by
\cref{thm:p15-iid-series-strong-law}. Also
\(\ell_n''(\theta_0)/n\to-I(\theta_0)\) almost surely by the same strong
law applied to the integrable \(\partial_\theta^2\log f_{\theta_0}(X_i)\).
Given \(\varepsilon>0\), choose \(\eta\) with
\(E_{\theta_0}h_\eta<\varepsilon\); then
\(P_{\theta_0}(|B_n/n+I(\theta_0)|>2\varepsilon)\to0\). Hence
\(-B_n/n\to I(\theta_0)\) in probability. The same bound with
\(\widehat\theta_n\) in place of the intermediate points shows
\(-\ell_n''(\widehat\theta_n)/n\to I(\theta_0)\) in probability.

\emph{Limit law.} The summands \(s_{\theta_0}(X_i)\) are independent
identically distributed with mean zero and variance
\(I(\theta_0)\in(0,\infty)\), so
\(\ell_n'(\theta_0)/\sqrt n\Rightarrow N(0,I(\theta_0))\) by
\cref{thm:p15-clt-iid}. Put
\(Y_n=(\ell_n'(\theta_0)/\sqrt n)/(-B_n/n)\), assigning the value zero on
\(\{B_n=0\}\). By \cref{thm:p15-clt-perturbation}, with the random
denominator tending in probability to \(I(\theta_0)\neq0\),
\(Y_n\Rightarrow N(0,I(\theta_0))/I(\theta_0)=N(0,I(\theta_0)^{-1})\). On
\(\Omega_n\cap\{B_n\neq0\}\), \(\sqrt n\,\Delta_n=Y_n\), and the
probability of this event tends to one because \(-B_n/n\to
I(\theta_0)>0\). Thus \(\sqrt n\,\Delta_n-Y_n\to0\) in probability, and the
same lemma gives \(\sqrt n\,\Delta_n\Rightarrow N(0,I(\theta_0)^{-1})\).
The asymptotic variance is \(I(\theta_0)^{-1}\), so the efficiency of
\cref{def:p15-w02b-p15-057-efficiency} equals one.
\end{proof}
\end{annexitem}

\begin{annexitem}{annex:p15-w02b-p15-057.mle}{thm:p15-w02b-p15-057-mle}{Maximum likelihood on a compact interval}
\begin{proof}
Hypothesis (W1) holds for \([c,d]\), and (W2)--(W4) are assumed, so
\cref{thm:p15-w02b-p15-057-wald} applies. Hence, almost surely, every maximizer converges to \(\theta_0\);
in particular a measurable choice converges in probability. Let \(N\)
and \(A\) be as in (R1), and shrink \(N\) to an open interval inside
\((c,d)\). Almost surely all \(X_i\in A\) and, for all large \(n\),
\(\widehat\theta_n\in N\). There \(\ell_n\) is differentiable and finite,
and \(\widehat\theta_n\) is an interior maximum of \(\ell_n\) on
\([c,d]\), so \(\ell_n'(\widehat\theta_n)=0\). Almost sure eventual
validity implies \(P_{\theta_0}(\ell_n'(\widehat\theta_n)=0)\to1\), and
\cref{thm:p15-w02b-p15-057-asymptotic} applies.
\end{proof}
\end{annexitem}

\begin{annexitem}{annex:p15-w02b-p15-057.information}{thm:p15-w02b-p15-057-information}{The information bound for a fixed sample size}
\begin{proof}
As in the previous proof, \(E_\theta s_\theta(X)=\int_A\partial_\theta
f_\theta\,d\mu=0\), and \(E_\theta s_\theta(X)^2=I(\theta)\). The score of
the sample is \(\ell_n'(\theta)=\sum_{i\leq n}s_\theta(X_i)\) almost surely;
its summands are independent with mean zero, so
\(E_\theta\ell_n'(\theta)=0\) and
\(\operatorname{Var}_\theta\ell_n'(\theta)=nI(\theta)\). Since
\(E_\theta T^2<\infty\),
\(\operatorname{Cov}_\theta(T,\ell_n'(\theta))=E_\theta[T\ell_n'(\theta)]
-E_\theta T\cdot0=m'(\theta)\). The Cauchy--Schwarz inequality gives
\(m'(\theta)^2\leq\operatorname{Var}_\theta(T)\,nI(\theta)\); divide by
\(nI(\theta)>0\). For an unbiased \(T\), \(m(\vartheta)=\vartheta\), so
\(m'=1\) and \(1/(nI(\theta)\operatorname{Var}_\theta T)\leq1\).
\end{proof}
\end{annexitem}

\begin{annexitem}{annex:p15-w02b-p15-057.hodges}{thm:p15-w02b-p15-057-hodges}{Efficiency above one at a single point}
\begin{proof}
Since \(\operatorname{Var}_\theta\overline X_n=1/n\), Chebyshev's inequality
gives \(P_\theta(|\overline X_n-\theta|\geq t)\leq1/(nt^2)\).

At \(\theta=0\), \(T_n\neq0\) only if \(|\overline X_n|\geq n^{-1/4}\), an
event of probability at most \(n^{-1/2}\). So
\(P_0(\sqrt n\,|T_n|>\varepsilon)\leq n^{-1/2}\to0\) for every
\(\varepsilon>0\).

At \(\theta\neq0\), once \(n^{-1/4}<|\theta|/2\), the event
\(T_n\neq\overline X_n\) requires \(|\overline X_n|<n^{-1/4}\) and hence
\(|\overline X_n-\theta|>|\theta|/2\), of probability at most
\(4/(n\theta^2)\to0\). Thus \(\sqrt n\,(T_n-\overline X_n)\to0\) in
probability. By \cref{thm:p15-clt-iid},
\(\sqrt n\,(\overline X_n-\theta)\Rightarrow N(0,1)\), and
\cref{thm:p15-clt-perturbation} transfers this limit to
\(\sqrt n\,(T_n-\theta)\).

For the normal location family the density is
\((2\pi)^{-1/2}e^{-(x-\theta)^2/2}\), the score is \(x-\theta\) and
\(I\equiv E(X-\theta)^2=1\). The asymptotic variances \(0\) and \(1\)
therefore give efficiencies \(\infty\) and \(1\).
\end{proof}
\end{annexitem}

\begin{annexitem}{annex:p15-w02b-p15-057.cauchy-check}{thm:p15-w02b-p15-057-cauchy-check}{The Cauchy location family meets every hypothesis}
\begin{proof}
Let \(g(u)=\pi^{-1}(1+u^2)^{-1}\), so \(f_\theta(x)=g(x-\theta)\).

(W1) and (W2) are immediate. (W3): \(f_\theta\) has its unique maximum at
\(x=\theta\), so distinct \(\theta\) give distinct continuous densities,
which differ on an interval and so define distinct laws. (W4): for
\(\theta,\theta_0\in[c,d]\),
\[
 \frac{f_\theta(x)}{f_{\theta_0}(x)}=\frac{1+(x-\theta_0)^2}{1+(x-\theta)^2}
 \leq\frac{1+2(x-\theta)^2+2(\theta-\theta_0)^2}{1+(x-\theta)^2}
 \leq2+2(d-c)^2,
\]
using \((u+v)^2\leq2u^2+2v^2\). So the positive part in (W4) is bounded.

Regularity. (R1) holds with \(A=\mathbb R\), and (R2) because \(\log g\)
is smooth. With \(\phi=\log g+\log\pi=-\log(1+u^2)\),
\(\phi'(u)=-2u/(1+u^2)\) and \(\phi''(u)=-2(1-u^2)/(1+u^2)^2\). Since
\(|1-u^2|\leq1+u^2\), \(|\phi''|\leq2\), which is (R5) with \(M\equiv2\),
because \(\partial_\theta^2\log f_\theta(x)=\phi''(x-\theta)\). For (R3),
\(\partial_\theta f_\theta(x)=-g'(x-\theta)\) and
\(\partial_\theta^2f_\theta(x)=g''(x-\theta)\); here
\(|g'|\leq2|u|/(\pi(1+u^2)^2)\) and \(|g''|\leq C(1+u^2)^{-2}\) for a
constant \(C\), both integrable, and \(g\), \(g'\) tend to zero at
\(\pm\infty\), so \(\int g'=\int g''=0\).

Information. By \cref{thm:p15-w02b-p15-057-location-scale} with \(a=1\)
(whose hypotheses for this \(\phi\) on \(D=\mathbb R\) are verified in its
proof), \(I(\theta_0)=-E\phi''=\tfrac12\), in \((0,\infty)\), giving (R4).

Conclusion. \cref{thm:p15-w02b-p15-057-mle} gives
\(\sqrt n(\widehat\theta_n-\theta_0)\Rightarrow N(0,2)\). The statement on
the sample mean is \cref{thm:p15-fe-cauchy}.
\end{proof}
\end{annexitem}

\begin{annexitem}{annex:p15-w02b-p15-057.minor}{thm:p15-w02b-p15-057-minor}{Fisher's minor-over-determinant rule}
\begin{proof}
A positive definite matrix has positive eigenvalues, so \(\det J>0\) and
\(J\) is invertible with \(J^{-1}=\operatorname{adj}(J)/\det J\). The
\((j,j)\) entry of the adjugate is the cofactor \((-1)^{2j}\det J_{(jj)}\),
which gives the first claim. For \(k=2\), \(\Sigma=J^{-1}\) is positive
definite, so \(\sigma_1,\sigma_2>0\) and
\(\det\Sigma=\sigma_1^2\sigma_2^2(1-r^2)>0\), whence \(|r|<1\). Inverting
the \(2\times2\) matrix \(\Sigma\),
\[
 J=\frac1{\sigma_1^2\sigma_2^2(1-r^2)}
 \begin{pmatrix}\sigma_2^2&-r\sigma_1\sigma_2\\
 -r\sigma_1\sigma_2&\sigma_1^2\end{pmatrix},
\]
whose entries are the three displayed values.
\end{proof}
\end{annexitem}

\begin{annexitem}{annex:p15-w02b-p15-057.multi}{thm:p15-w02b-p15-057-multi}{Joint limit of a local root of the likelihood equations}
\begin{proof}
Let \(\lambda>0\) be the least eigenvalue of \(J\), so that
\(|Jv|\geq\lambda|v|\), \(v^{\top}Jv\geq\lambda|v|^2\) and
\(\|J^{-1}\|=1/\lambda\). Let \(G_n\) be the events of (S) and
\(S_n=\nabla\ell_n(\theta_0)\), equal to \(\sum_{i\leq n}s(X_i)\) on
\(G_n\).

\emph{A matrix bound.} If \(A\) is a \(k\times k\) matrix with
\(\|A-J\|\leq\varepsilon\leq\lambda/2\), then for \(|v|=1\),
\(|Av|\geq|Jv|-\varepsilon\geq\lambda/2\), so \(A\) is invertible with
\(\|A^{-1}\|\leq2/\lambda\), and
\(A^{-1}-J^{-1}=A^{-1}(J-A)J^{-1}\) gives
\(\|A^{-1}-J^{-1}\|\leq2\varepsilon/\lambda^2\).

\emph{The normalized score is tight.} The vectors \(s(X_i)\) are
independent with mean zero and covariance \(J\), so
\(E\bigl|n^{-1/2}\sum_{i\leq n}s(X_i)\bigr|^2=\operatorname{tr}J\), and by
Markov's inequality
\(P\bigl(|n^{-1/2}\sum_{i\leq n}s(X_i)|>c\bigr)\leq\operatorname{tr}J/c^2\).

\emph{Item 1.} Given \(\eta>0\), choose \(c\) with
\(\operatorname{tr}J/c^2<\eta/2\) and \(K>4c/\lambda\). Let \(E_n\) be the
intersection of \(\Omega_n(K,\lambda/2)\), \(G_n\) and
\(\{|n^{-1/2}S_n|\leq c\}\); then \(P(E_n)\geq1-\eta\) for all large
\(n\). Fix a sample in \(E_n\). For \(\theta\) in the ball
\(B=B(\theta_0,Kn^{-1/2})\) and \(v\in\mathbb R^k\),
\[
 v^{\top}\bigl(-n^{-1}\nabla^2\ell_n(\theta)\bigr)v
 =v^{\top}Jv-v^{\top}\bigl(J+n^{-1}\nabla^2\ell_n(\theta)\bigr)v
 \geq\tfrac{\lambda}{2}|v|^2 .
\]
Taylor's formula with integral remainder along the segment from
\(\theta_0\) to \(\theta_0+u\in B\), valid because \(\ell_n\) is twice
continuously differentiable on an open set containing the convex set
\(B\), gives
\[
 \ell_n(\theta_0+u)-\ell_n(\theta_0)
 =u^{\top}S_n+\int_0^1(1-t)\,u^{\top}\nabla^2\ell_n(\theta_0+tu)\,u\,dt
 \leq|u|\,|S_n|-\tfrac{n\lambda}{4}|u|^2 .
\]
For \(|u|=Kn^{-1/2}\) the right side is at most
\(Kc-\lambda K^2/4<0\). The continuous function \(\ell_n\) attains its
maximum over the compact \(B\), and by the last inequality not on the
boundary sphere; at an interior maximizer the gradient vanishes. If
\(\theta_1\neq\theta_2\) were stationary points in the open ball, the
function \(\varphi(t)=\ell_n(\theta_1+t(\theta_2-\theta_1))\) on
\([0,1]\) would have \(\varphi'(0)=\varphi'(1)=0\) and
\(\varphi''<0\) by the displayed quadratic bound, which is impossible
since \(\varphi'\) is strictly decreasing. The same strict concavity makes
the maximizer over \(B\) unique, and it is the stationary point.

\emph{The expansion of a root.} Fix \(0<\varepsilon\leq\lambda/2\) and
\(\eta>0\). By tightness choose \(K\) with
\(P(\sqrt n\,|\widehat\theta_n-\theta_0|\leq K)\geq1-\eta\) for every
\(n\). Let \(F_n\) be the intersection of \(\Omega_n(K,\varepsilon)\),
\(G_n\), \(\{\sqrt n\,|\widehat\theta_n-\theta_0|\leq K\}\) and
\(\{\nabla\ell_n(\widehat\theta_n)=0\}\); then
\(P(F_n)\geq1-2\eta\) for all large \(n\). On \(F_n\) put
\(\Delta_n=\widehat\theta_n-\theta_0\) and
\(B_n=\int_0^1\nabla^2\ell_n(\theta_0+t\Delta_n)\,dt\), the integral taken
entry by entry. The fundamental theorem of calculus applied to
\(t\mapsto\nabla\ell_n(\theta_0+t\Delta_n)\) gives
\(0=\nabla\ell_n(\widehat\theta_n)=S_n+B_n\Delta_n\). The matrix
\(A_n=-n^{-1}B_n\) satisfies \(\|A_n-J\|\leq\varepsilon\), the norm of an
integral being at most the supremum of the norms, so \(A_n\) is invertible
and \(\sqrt n\,\Delta_n=A_n^{-1}n^{-1/2}S_n\). Hence on \(F_n\)
\[
 \Bigl|\sqrt n\,\Delta_n-J^{-1}n^{-1/2}\sum_{i\leq n}s(X_i)\Bigr|
 \leq\frac{2\varepsilon}{\lambda^2}\,\Bigl|n^{-1/2}\sum_{i\leq n}s(X_i)\Bigr| .
\]
For \(\rho>0\) the probability that the left side exceeds \(\rho\) is at
most \(2\eta+(2\varepsilon/\lambda^2)^2\operatorname{tr}J/\rho^2\) for
large \(n\). Since \(\eta\) and \(\varepsilon\) are arbitrary, the
difference tends to zero in probability. On \(F_n\) also
\(\|n^{-1}\nabla^2\ell_n(\widehat\theta_n)+J\|\leq\varepsilon\), which
gives the last assertion of item~2.

\emph{Linear combinations.} For \(t\in\mathbb R^k\),
\(t^{\top}J^{-1}n^{-1/2}\sum_is(X_i)=n^{-1/2}\sum_it^{\top}J^{-1}s(X_i)\)
is a normalized sum of independent identically distributed variables with
mean zero and variance \(t^{\top}J^{-1}JJ^{-1}t=t^{\top}J^{-1}t\). By
\cref{thm:p15-clt-iid} it converges in law to \(N(0,t^{\top}J^{-1}t)\),
and by \cref{thm:p15-clt-perturbation} so does
\(\sqrt n\,t^{\top}\Delta_n\), which differs from it by a quantity tending
to zero in probability.

\emph{Functions.} Let \(\tau\) be Borel and differentiable at
\(\theta_0\) with gradient \(g\), and write
\(\tau(\theta_0+u)=\tau(\theta_0)+g^{\top}u+|u|\,r(u)\) with
\(\sup_{|u|\leq\rho}|r(u)|\to0\) as \(\rho\to0\). Then
\(\sqrt n\,(\tau(\widehat\theta_n)-\tau(\theta_0))=\sqrt n\,g^{\top}\Delta_n
+\sqrt n\,|\Delta_n|\,r(\Delta_n)\). For \(\rho>0\),
\(P(\sqrt n|\Delta_n||r(\Delta_n)|>\rho)\leq P(\sqrt n|\Delta_n|>K)
+\mathbf 1\{\sup_{|u|\leq Kn^{-1/2}}|r(u)|>\rho/K\}\), which is at most
\(\eta\) for large \(n\). The first term converges in law to
\(N(0,g^{\top}J^{-1}g)\) by the previous paragraph with \(t=g\), and
\cref{thm:p15-clt-perturbation} completes item~2.

\emph{Item 3.} Take \(\varepsilon=\lambda/2\) and \(\delta\) from
(C\(^{+}\)). On the intersection of its events with \(G_n\),
\(\{|\widehat\theta_n-\theta_0|<\delta\}\) and
\(\{\nabla\ell_n(\widehat\theta_n)=0\}\), an event of probability tending
to one, the expansion above holds with \(B(\theta_0,\delta)\) in place of
\(B(\theta_0,Kn^{-1/2})\), so
\(\sqrt n\,|\Delta_n|\leq(2/\lambda)|n^{-1/2}S_n|\). The right side is
tight by the score bound, so \(\sqrt n\,\Delta_n\) is tight.
\end{proof}
\end{annexitem}

\begin{annexitem}{annex:p15-w02b-p15-057.multi-bound}{thm:p15-w02b-p15-057-multi-bound}{The information bound for a function of several parameters}
\begin{proof}
Write \(S=\nabla\ell_n(\theta_0)=\sum_{i\leq n}s(X_i)\). By independence
\(ES=0\) and \(E\,SS^{\top}=nJ\). By the Cauchy--Schwarz inequality
\(TS\) is integrable, so \(g=E[TS]=\operatorname{Cov}(T,S)\) is defined.
Put \(c=(nJ)^{-1}g\) and \(U=T-ET-c^{\top}S\). Then
\[
 EU^2=\operatorname{Var}T-2c^{\top}g+c^{\top}(nJ)c
 =\operatorname{Var}T-\tfrac1n\,g^{\top}J^{-1}g\geq0,
\]
with equality exactly when \(U=0\) almost surely, which is the stated
condition. For the last sentence, differentiation under the integral sign
means
\(\nabla_\vartheta\int T\prod_if_\vartheta(x_i)\,d\mu^n=\int
T\,\nabla_\vartheta\prod_if_\vartheta(x_i)\,d\mu^n\) at \(\theta_0\). Where
\(\prod_if_{\theta_0}(x_i)>0\) the integrand gradient equals
\(T\prod_if_{\theta_0}(x_i)\,\nabla\ell_n(\theta_0)\); where the product
vanishes, \(\theta_0\) minimizes the nonnegative differentiable function
\(\vartheta\mapsto\prod_if_\vartheta(x_i)\), so its gradient is zero.
The right side is therefore \(E_{\theta_0}[T\nabla\ell_n(\theta_0)]=g\),
and the left side is the gradient of \(\tau\) when \(T\) is unbiased.
\end{proof}
\end{annexitem}

\begin{annexitem}{annex:p15-w02b-p15-057.multi-regular}{thm:p15-w02b-p15-057-multi-regular}{The regular case}
\begin{proof}
Almost surely every \(X_i\in A\), and then \(\ell_n\) is finite and twice
continuously differentiable on \(N\) with
\(\nabla\ell_n(\theta_0)=\sum_is_{\theta_0}(X_i)\). On \(A\),
\(s_{\theta_0}f_{\theta_0}=\nabla f_{\theta_0}\), so
\(E_{\theta_0}s_{\theta_0}=\int_A\nabla f_{\theta_0}\,d\mu=0\); with the
assumption on \(J\) this is (S). Also
\(\nabla^2\log f_{\theta_0}=\nabla^2f_{\theta_0}/f_{\theta_0}
-s_{\theta_0}s_{\theta_0}^{\top}\) on \(A\); the left side is integrable by
the domination and \(s_{\theta_0}s_{\theta_0}^{\top}\) by finiteness of
\(J\), so \(E_{\theta_0}\nabla^2\log f_{\theta_0}
=\int_A\nabla^2f_{\theta_0}\,d\mu-J=-J\).

For (C\(^{+}\)), let
\(h_\eta(x)=\sup_{|\theta-\theta_0|\leq\eta}\|\nabla^2\log f_\theta(x)
-\nabla^2\log f_{\theta_0}(x)\|\) for \(x\in A\) and \(\eta<\delta\) with
\(B(\theta_0,\eta)\subseteq N\). By continuity the supremum may be taken
over a countable dense subset of the ball, so \(h_\eta\) is measurable;
\(0\leq h_\eta\leq2M\) and \(h_\eta\to0\) pointwise as \(\eta\to0\), so
\(E_{\theta_0}h_\eta\to0\) by dominated convergence. Given
\(\varepsilon>0\), choose \(\eta\) with \(E_{\theta_0}h_\eta<\varepsilon/2\).
For \(|\theta-\theta_0|\leq\eta\),
\[
 \bigl\|n^{-1}\nabla^2\ell_n(\theta)+J\bigr\|
 \leq\frac1n\sum_{i\leq n}h_\eta(X_i)
 +\bigl\|n^{-1}\nabla^2\ell_n(\theta_0)+J\bigr\| .
\]
By \cref{thm:p15-iid-series-strong-law} the first term converges almost
surely to \(E_{\theta_0}h_\eta<\varepsilon/2\), and, entry by entry, the
second converges almost surely to zero, the operator norm being at most
the sum of the absolute values of the entries. The events on which every
\(X_i\in A\), the first term is below \(3\varepsilon/4\) and the second
below \(\varepsilon/4\) have probability tending to one, which is
(C\(^{+}\)) with \(\delta=\eta\). The final sentence is item~3 followed by
item~2 of \cref{thm:p15-w02b-p15-057-multi}.
\end{proof}
\end{annexitem}

\begin{annexitem}{annex:p15-w02b-p15-057.normal-scale}{thm:p15-w02b-p15-057-normal-scale}{Mean error and mean square error for a normal scale}
\begin{proof}
Write \(Z=(X-\mu)/\sigma\), a standard normal variable, and
\(\varphi(z)=(2\pi)^{-1/2}e^{-z^2/2}\). Since
\(\int_0^\infty z\varphi(z)\,dz=(2\pi)^{-1/2}\) and
\(\int_0^\infty z^3\varphi(z)\,dz=2(2\pi)^{-1/2}\) (substitute
\(w=z^2/2\)), symmetry gives \(E|Z|=\sqrt{2/\pi}\) and
\(E|Z|^3=2\sqrt{2/\pi}\); also \(EZ^2=1\) and \(EZ^4=3\).

\emph{Item 1.} \(\sqrt{\pi/2}\,E|X-\mu|=\sigma\sqrt{\pi/2}\,E|Z|=\sigma\)
and \(E(X-\mu)^2=\sigma^2\). By \cref{thm:p15-iid-series-strong-law},
\(A_n=n^{-1}\sum_{i\leq n}|X_i-\mu|\to\sigma\sqrt{2/\pi}\),
\(\bar X_n\to\mu\) and \(V_n=n^{-1}\sum_{i\leq n}(X_i-\mu)^2\to\sigma^2\)
almost surely. Since \(\bigl||X_i-\bar X_n|-|X_i-\mu|\bigr|\leq|\bar
X_n-\mu|\), the mean of the \(|X_i-\bar X_n|\) differs from \(A_n\) by at
most \(|\bar X_n-\mu|\to0\), so \(\sigma_1\to\sigma\) almost surely. Also
\(\sigma_2^2=V_n-(\bar X_n-\mu)^2\to\sigma^2\), so \(\sigma_2\to\sigma\)
almost surely.

\emph{The centring of the mean error.} Put \(h=\bar X_n-\mu\) and
\(D_n=n^{-1}\sum_{i\leq n}(|X_i-\mu-h|-|X_i-\mu|)\). For real \(y\), if
\(|y|>|h|\) then \(y-h\) and \(y\) have the same sign and
\(|y-h|-|y|=-h\operatorname{sgn}y\); if \(|y|\leq|h|\) then
\(\bigl||y-h|-|y|+h\operatorname{sgn}y\bigr|\leq2|h|\). Hence
\[
 |D_n+hG_n|\leq2|h|\,H_n(|h|),\qquad
 G_n=\frac1n\sum_{i\leq n}\operatorname{sgn}(X_i-\mu),\quad
 H_n(\eta)=\frac1n\#\{i\leq n:|X_i-\mu|\leq\eta\}.
\]
Now \(\sqrt n\,h\) has the law \(N(0,\sigma^2)\) for every \(n\), so for
\(\varepsilon>0\) there is \(M\) with \(P(\sqrt n|h|>M)<\varepsilon\) for
all \(n\). The variables \(\operatorname{sgn}(X_i-\mu)\) have mean zero, so
\(G_n\to0\) almost surely by \cref{thm:p15-iid-series-strong-law}; hence
\(\sqrt n|hG_n|\to0\) in probability. Given \(\gamma>0\) choose
\(\eta>0\) with \(P(|X-\mu|\leq\eta)<\gamma\); on \(\{|h|\leq\eta\}\),
whose probability tends to one, \(H_n(|h|)\leq H_n(\eta)\), and
\(H_n(\eta)\to P(|X-\mu|\leq\eta)<\gamma\) almost surely. So
\(\sqrt n|D_n+hG_n|\leq2\sqrt n|h|\,H_n(|h|)\) is at most \(2M\gamma\)
with probability at least \(1-2\varepsilon\) for all large \(n\). As
\(\gamma\) and \(\varepsilon\) are arbitrary, \(\sqrt n D_n\to0\) in
probability.

\emph{Linear representations.} By the last paragraph,
\[
 \sqrt n(\sigma_1-\sigma)=\frac1{\sqrt n}\sum_{i\leq n}Y_{1i}+R_{1n},\qquad
 Y_{1i}=\sqrt{\tfrac\pi2}\bigl(|X_i-\mu|-\sigma\sqrt{2/\pi}\bigr),
\]
with \(R_{1n}=\sqrt{\pi/2}\,\sqrt n D_n\to0\) in probability. For the
second statistic put \(B_n=n^{-1/2}\sum_{i\leq n}((X_i-\mu)^2-\sigma^2)\);
then \(\sqrt n(\sigma_2^2-\sigma^2)=B_n-\sqrt n h^2\), and
\(E\sqrt n h^2=\sigma^2/\sqrt n\to0\), so \(\sqrt nh^2\to0\) in
probability by Markov's inequality. Hence
\[
 \sqrt n(\sigma_2-\sigma)=\frac{\sqrt n(\sigma_2^2-\sigma^2)}{\sigma_2+\sigma}
 =\frac{B_n}{2\sigma}+R_{2n},\qquad
 R_{2n}=B_n\Bigl(\frac1{\sigma_2+\sigma}-\frac1{2\sigma}\Bigr)
 -\frac{\sqrt nh^2}{\sigma_2+\sigma}.
\]
Since \(EB_n^2=E(X-\mu)^4-\sigma^4=2\sigma^4\), Chebyshev's inequality
makes \(B_n\) tight; the bracket tends to zero almost surely by item~1, so
\(R_{2n}\to0\) in probability. Put \(Y_{2i}=((X_i-\mu)^2-\sigma^2)/(2\sigma)\).

\emph{Item 2.} The pairs \((Y_{1i},Y_{2i})\) are independent and
identically distributed with mean zero, and
\[
 EY_1^2=\tfrac\pi2\sigma^2\bigl(1-\tfrac2\pi\bigr)=v_1,\quad
 EY_2^2=\frac{\sigma^4(EZ^4-1)}{4\sigma^2}=v_2,\quad
 EY_1Y_2=\frac{\sqrt{\pi/2}}{2\sigma}\,\sigma^3\bigl(E|Z|^3-E|Z|\bigr)=c,
\]
the last because \(E[|Z|(Z^2-1)]=\sqrt{2/\pi}\) and \(\sqrt{\pi/2}\sqrt{2/\pi}=1\).
The determinant \(v_1v_2-c^2=\tfrac14\sigma^4(\pi-3)\) is positive, so for
\((s,t)\neq(0,0)\) the variable \(sY_1+tY_2\) has positive variance
\(s^2v_1+2stc+t^2v_2\). By \cref{thm:p15-clt-iid},
\(n^{-1/2}\sum_{i\leq n}(sY_{1i}+tY_{2i})\) converges in law to the normal
law with that variance, and adding \(sR_{1n}+tR_{2n}\to0\) does not change
the limit by \cref{thm:p15-clt-perturbation}. For \(s=t=0\) both sides are
the point mass at zero.

\emph{Item 3.} With \(\theta=(\mu,\sigma)\) and support \(\mathbb R\),
\(\log f_\theta(x)=-\log\sigma-\tfrac12\log(2\pi)-(x-\mu)^2/(2\sigma^2)\)
is infinitely differentiable, with score
\(s_\theta=\bigl((x-\mu)/\sigma^2,\,-1/\sigma+(x-\mu)^2/\sigma^3\bigr)\)
and Hessian entries \(-1/\sigma^2\), \(-2(x-\mu)/\sigma^3\) and
\(1/\sigma^2-3(x-\mu)^2/\sigma^4\). On the ball of radius
\(\delta=\min(1,\sigma_0/2)\) about \((\mu_0,\sigma_0)\) these entries are
bounded in absolute value by \(M(x)=C(1+(x-\mu_0)^2)\) for a constant \(C\)
depending on \(\sigma_0\) only, and \(E_{\theta_0}M(X)<\infty\). The
derivatives of \(f_\theta=\exp(\log f_\theta)\) are bounded on the same
ball by \(C'(1+(x-\mu_0)^4)\sup_\theta f_\theta(x)\), and
\(\sup_\theta f_\theta(x)\leq C''e^{-(|x-\mu_0|-1)_+^2/(2(3\sigma_0/2)^2)}\),
so differentiating \(\int f_\theta\,dx=1\) twice under the integral sign
gives \(\int\nabla f_{\theta_0}=0\) and \(\int\nabla^2f_{\theta_0}=0\). So
\cref{def:p15-w02b-p15-057-multi-regular} holds, and
\(E s_\theta s_\theta^\top=\operatorname{diag}(EZ^2/\sigma^2,
E(Z^2-1)^2/\sigma^2)=\operatorname{diag}(1/\sigma^2,2/\sigma^2)\), the cross
term being \(E[Z(Z^2-1)]/\sigma^2=0\). The gradient of \(\tau\) is
\(g=(0,1)\), and \(g^\top J^{-1}g=\sigma^2/2\). The efficiencies are
\((\sigma^2/2)/v_2=1\) and \((\sigma^2/2)/v_1=1/(\pi-2)\); numerically
\(\pi-2=1.14159\) and \(1/(\pi-2)=0.87597\).

\emph{Item 4.} \(\sqrt{v_1v_2}=\tfrac12\sigma^2\sqrt{\pi-2}\), so
\(r=1/\sqrt{\pi-2}\); then
\(r\sqrt{v_1}=\sigma\sqrt{(\pi-2)/2}/\sqrt{\pi-2}=\sigma/\sqrt2=\sqrt{v_2}\)
and \(r^2=1/(\pi-2)=v_2/v_1\).

\emph{Item 5.} Since
\(\sum_i(x_i-\mu)^2=S(x-\bar x)^2+n(\bar x-\mu)^2\), the joint density of
the sample with respect to Lebesgue measure on \(\mathbb R^n\) is
\((2\pi\sigma^2)^{-n/2}\exp\{-(S(x-\bar x)^2+n(\bar x-\mu)^2)/(2\sigma^2)\}\),
a function of \((\mu,\sigma)\) and of the pair alone, so
\cref{thm:p15-w02b-p15-057-factorization} applies with \(h=1\); and
\(\sigma_2=(S(x-\bar x)^2/n)^{1/2}\).
\end{proof}
\end{annexitem}

\begin{annexitem}{annex:p15-w02b-p15-057.suff-eff}{thm:p15-w02b-p15-057-suff-eff}{Fisher's sufficiency-to-efficiency argument for a bivariate normal pair}
\begin{proof}
\emph{Item 1.} Put \(\beta=r\sigma_2/\sigma_1\) and
\(W=V-\theta-\beta(U-\theta)\). The pair \((U,W)\) is an invertible linear
image of \((U,V)\), so it has a bivariate normal density, and
\(\operatorname{Cov}(U,W)=r\sigma_1\sigma_2-\beta\sigma_1^2=0\). A
bivariate normal density with diagonal covariance matrix is the product
of its marginal densities, so \(U\) and \(W\) are independent, and \(W\)
has the law
\(N(0,\sigma_2^2-2\beta r\sigma_1\sigma_2+\beta^2\sigma_1^2)
=N(0,\sigma_2^2(1-r^2))\). For Borel sets \(A,B\), independence and
Fubini's theorem give
\(P_\theta(U\in A,V\in B)=\int_A P(\theta+\beta(u-\theta)+W\in B)\,
P_\theta^U(du)\), which is the defining property of the stated kernel as a
version of the conditional law.

\emph{Item 2.} If \(r\sigma_2=\sigma_1\), then \(\beta=1\) and the kernel of
item~1 is \(N(u,\sigma_2^2(1-r^2))\), free of \(\theta\). Conversely,
suppose one kernel \(\kappa\) is a version under \(\theta\) and under some
\(\theta^*\neq\theta\) in \(\Theta\). For each rational \(q\), the function
\(u\mapsto\kappa(u,(-\infty,q])\) and the corresponding function computed
from item~1 are both versions of the conditional probability of
\(\{V\leq q\}\) given \(U\) under \(\theta\), so they agree for
\(P_\theta^U\)-almost every \(u\) by the almost sure uniqueness of
conditional expectation (\cref{def:p15-general-conditional-expectation});
the same holds under \(\theta^*\). The law of \(U\) under either parameter
has an everywhere positive density, so its null sets are the Lebesgue null
sets. Off one Lebesgue null set, therefore, the normal laws
\(N(\theta+\beta(u-\theta),\sigma_2^2(1-r^2))\) and
\(N(\theta^*+\beta(u-\theta^*),\sigma_2^2(1-r^2))\) have equal
distribution functions at every rational point, hence everywhere by right
continuity, hence are equal and have equal means:
\((\theta-\theta^*)(1-\beta)=0\). So \(\beta=1\), that is,
\(r\sigma_2=\sigma_1\).

\emph{Item 3.} From \(r\sigma_2=\sigma_1>0\) and \(\sigma_2>0\), \(r>0\);
with \(r<1\), \(\sigma_1=r\sigma_2<\sigma_2\) and
\(\sigma_1^2/\sigma_2^2=r^2\).
\end{proof}
\end{annexitem}

\begin{annexitem}{annex:p15-w02b-p15-057.typeiii}{thm:p15-w02b-p15-057-typeiii}{The Type III information matrix and the moment efficiencies}
\begin{proof}
Put \(k=p+1>0\). The variable \(\xi=(X-m)/a\) has density
\(\xi^{p}e^{-\xi}/\Gamma(k)\) on \((0,\infty)\), so
\(E\xi^r=\Gamma(k+r)/\Gamma(k)\) whenever \(k+r>0\) and
\(E\xi^r=\infty\) for \(r\leq-k\). In particular \(E\xi=k\),
\(E\xi^2=k(k+1)\), \(E\xi^{-1}=1/p\) for \(p>0\) and
\(E\xi^{-2}=1/(p(p-1))\) for \(p>1\), while \(E\xi^{-2}=\infty\) for
\(0<p\leq1\). Write \(\psi\) for the derivative of \(\log\Gamma\); then
\(\frac{d}{dp}\log(p!)=\psi(k)\) and \(F(p)=\psi'(k)\). Differentiation of
\(\int_0^\infty\xi^{s}e^{-\xi}\,d\xi=\Gamma(s+1)\) under the integral, valid
for \(s>-1\) by domination of \(|\log\xi|^j\xi^{s}e^{-\xi}\) locally
uniformly in \(s\), gives \(E\log\xi=\psi(k)\),
\(\operatorname{Var}\log\xi=\psi'(k)\), \(E[\xi\log\xi]=k\psi(k+1)\) and,
for \(p>0\), \(E[\xi^{-1}\log\xi]=\psi(p)/p\). The recurrence
\(\Gamma(s+1)=s\Gamma(s)\) gives \(\psi(s+1)=\psi(s)+1/s\).

\emph{Derivatives.} For \(x>m\),
\(\log f=-\log a-\log\Gamma(k)+p\log\xi-\xi\) with
\(\partial_m\xi=-1/a\), \(\partial_a\xi=-\xi/a\). Hence
\[
 \partial_m\log f=\frac{1-p/\xi}{a},\quad
 \partial_a\log f=\frac{\xi-k}{a},\quad
 \partial_p\log f=\log\xi-\psi(k),
\]
\[
 \partial_m^2\log f=-\frac{p}{a^2\xi^2},\quad
 \partial_m\partial_a\log f=-\frac1{a^2},\quad
 \partial_m\partial_p\log f=-\frac1{a\xi},
\]
\[
 \partial_a^2\log f=\frac{k}{a^2}-\frac{2\xi}{a^2},\quad
 \partial_a\partial_p\log f=-\frac1a,\quad
 \partial_p^2\log f=-\psi'(k).
\]
These are Fisher's displays (1) and (2) and the derivatives on pp.~333--335
for one observation.

\emph{Item 1.} The support statement is immediate. The value of
\(E[-\partial_m^2\log f]=pE\xi^{-2}/a^2\) follows from the moments above:
for \(p>1\) it is \(1/(a^2(p-1))\); for \(0<p\leq1\) the integrand is
positive with infinite mean; for \(p=0\) it vanishes identically; for
\(-1<p<0\) it is negative with infinite mean. Likewise
\(E[-\partial_m\partial_p\log f]=E\xi^{-1}/a\), finite exactly when
\(p>0\), and then equal to \(1/(ap)\).

\emph{Item 2.} For \(p>1\) every entry is finite. The remaining expectations
are \(E[-\partial_a^2\log f]=(2k-k)/a^2=k/a^2\),
\(E[-\partial_m\partial_a\log f]=1/a^2\),
\(E[-\partial_a\partial_p\log f]=1/a\) and \(E[-\partial_p^2\log f]=F\),
giving \(J\). The scores have covariances equal to these entries:
\(E(1-p/\xi)^2/a^2=(1-2pE\xi^{-1}+p^2E\xi^{-2})/a^2
=(1-2+p/(p-1))/a^2=1/(a^2(p-1))\);
\(E[(1-p/\xi)(\xi-k)]/a^2=(E\xi-k-p+pkE\xi^{-1})/a^2=(k-p)/a^2=1/a^2\);
\(E(\xi-k)^2/a^2=k/a^2\);
\(E[(1-p/\xi)(\log\xi-\psi(k))]/a
=(E\log\xi-\psi(k)-pE[\xi^{-1}\log\xi]+p\psi(k)E\xi^{-1})/a
=(\psi(k)-\psi(p))/a=1/(ap)\);
\(\operatorname{Cov}(\xi,\log\xi)/a=(k\psi(k+1)-k\psi(k))/a=1/a\); and
\(\operatorname{Var}\log\xi=F\). The scores have mean zero by the same
moments. Writing \(J=DMD\) with \(D=\operatorname{diag}(1/a,1/a,1)\) and
\(M\) the matrix with rows \((1/(p-1),1,1/p)\), \((1,p+1,1)\),
\((1/p,1,F)\), cofactor expansion along the first row gives
\[
 \det M=\frac{(p+1)F-1}{p-1}-\Bigl(F-\frac1p\Bigr)
 +\frac1p\Bigl(1-\frac{p+1}{p}\Bigr)
 =\frac{1}{p-1}\Bigl(2F-\frac2p+\frac1{p^2}\Bigr),
\]
after placing the last two terms over \(p-1\). Hence
\(\det J=\det M/a^4\). As a covariance matrix \(J\) is positive
semidefinite. It is definite: if
\(\alpha(1-p/\xi)+\beta(\xi-k)+\gamma(\log\xi-\psi(k))\) were almost surely
constant, then, the gamma law charging every interval of \((0,\infty)\)
and the expression being continuous in \(\xi\), a relation
\(\alpha'\xi^{-1}+\beta\xi+\gamma\log\xi=\text{const}\) would hold on
\((0,\infty)\); letting \(\xi\to0\) forces \(\alpha'=0\), then
\(\gamma=0\), then \(\beta=0\), with \(\alpha'=-\alpha p\) and \(p>1\).
So \(\det J>0\). By \cref{thm:p15-w02b-p15-057-minor}, the \((p,p)\) entry
of \(J^{-1}\) is the determinant of the \((m,a)\) block,
\((1/a^4)\bigl((p+1)/(p-1)-1\bigr)=2/(a^4(p-1))\), divided by \(\det J\),
which is \(2/(2F-2/p+1/p^2)\). The \((m,a)\) block
\(a^{-2}\bigl(\begin{smallmatrix}1/(p-1)&1\\1&p+1\end{smallmatrix}\bigr)\)
has determinant \(2/(a^4(p-1))\) and inverse
\(\frac{a^2(p-1)}2\bigl(\begin{smallmatrix}p+1&-1\\-1&1/(p-1)
\end{smallmatrix}\bigr)\), whose diagonal is as stated. The
\((m,m)\) entry alone inverts to \(a^2(p-1)\).

\emph{Item 3.} \(E\overline x=m+aE\xi=m+ak\) and
\(\operatorname{Var}\overline x=a^2\operatorname{Var}\xi/n=a^2k/n\), since
\(\operatorname{Var}\xi=k(k+1)-k^2=k\) and the observations are independent.
So \(m_\mu\) is unbiased with variance \(a^2(p+1)/n\). The ratio of
\(a^2(p-1)/n\) to it is \((p-1)/(p+1)\).

\emph{Item 4.} Let \(\mu_2=a^2k\) and \(\mu_4=E(X-EX)^4=a^4E(\xi-k)^4\).
Expanding with the moments \(E\xi^r=k(k+1)\cdots(k+r-1)\) for
\(r=1,2,3,4\),
\[
 E(\xi-k)^4=E\xi^4-4kE\xi^3+6k^2E\xi^2-4k^3E\xi+k^4=3k^2+6k .
\]
So \(\mu_4-\mu_2^2=a^4(2k^2+6k)=2a^4k(k+3)\), finite for every \(p>-1\).
Write \(m_2=n^{-1}\sum(X_i-\overline X)^2=V_n-(\overline X-EX)^2\) with
\(V_n=n^{-1}\sum(X_i-EX)^2\). By \cref{thm:p15-clt-iid} applied to the
variables \((X_i-EX)^2\), of mean \(\mu_2\) and variance
\(\mu_4-\mu_2^2>0\),
\(\sqrt n\,(V_n-\mu_2)\Rightarrow N(0,\mu_4-\mu_2^2)\). Also
\(E[\sqrt n(\overline X-EX)^2]=\mu_2/\sqrt n\to0\), so that term tends to
zero in probability by Markov's inequality, and
\cref{thm:p15-clt-perturbation} gives
\(\sqrt n\,(m_2-\mu_2)\Rightarrow N(0,\mu_4-\mu_2^2)\). By
\cref{thm:p15-iid-series-strong-law}, \(m_2\to\mu_2\) almost surely. Now
\[
 \sqrt n\,(a_\mu-a)=\frac{\sqrt n\,(m_2-\mu_2)}
 {\sqrt{k}\,(\sqrt{m_2}+\sqrt{\mu_2})},
\]
whose denominator tends in probability to \(2\sqrt{k\mu_2}=2a k\). By
\cref{thm:p15-clt-perturbation} the limit law is normal with variance
\(2a^4k(k+3)/(4a^2k^2)=a^2(k+3)/(2k)=a^2(p+4)/(2(p+1))\). For \(p>1\),
dividing the inverse-information value \(a^2/2\) of item~2 by it gives
\((p+1)/(p+4)\).

\emph{Item 5.} Summing the \(a\)-score over the sample, the likelihood
equation \(\partial_aL=-nk/a+a^{-2}\sum(x_i-m)=0\) is equivalent, for
\(a>0\), to \(\overline x-m=a(p+1)\). The moment estimates are by
definition solutions of the moment equations, the first of which is
\(\overline x=m_\mu+a_\mu(p_\mu+1)\), since \(EX=m+a(p+1)\).
\end{proof}
\end{annexitem}

\begin{annexitem}{annex:p15-w02b-p15-057.typeiii-joint}{thm:p15-w02b-p15-057-typeiii-joint}{The joint Type III limit for \(p>1\)}
\begin{proof}
Write \(\theta_0=(m,a,p)\), \(k=p+1>2\), \(Y_i=X_i-m\) and
\(\xi_i=Y_i/a\), so that \(\xi_i\) has the gamma law with shape \(k\) and
unit scale, and \(Y_i>0\) almost surely. For a parameter point
\(\theta'=(m',a',p')\) and \(x>m'\),
\(\log f_{\theta'}(x)=p'\log(x-m')-(x-m')/a'-(p'+1)\log a'-\log\Gamma(p'+1)\);
for \(x\leq m'\) the density is zero. The moments used below are those of
the proof of \cref{thm:p15-w02b-p15-057-typeiii}:
\(EY^{-2}=1/(a^2p(p-1))\), \(EY^{-1}=1/(ap)\), \(EY=ak\).

\emph{Condition (S).} On the event that every \(X_i>m\), which has
probability one, \(\ell_n\) is finite and infinitely differentiable on
the open set \(\{m'<\min_iX_i,\ a'>0,\ p'>-1\}\), which contains
\(\theta_0\), and its gradient at \(\theta_0\) is \(\sum_is(X_i)\) with
\(s\) the vector of first derivatives computed in the proof of
\cref{thm:p15-w02b-p15-057-typeiii}. That proof shows that \(s(X)\) has mean
zero and covariance \(J\), and that \(J\) is positive definite for
\(p>1\).

\emph{The Hessian.} At \(\theta'\) in that open set, with
\(u=m'-m\), the entries of \(n^{-1}\nabla^2\ell_n(\theta')\) for the
pairs \(mm,ma,mp,aa,ap,pp\) are
\[
 -p'Q_{-2}(u),\quad -\frac1{a'^2},\quad -Q_{-1}(u),\quad
 \frac{p'+1}{a'^2}-\frac{2}{a'^3}Q_1(u),\quad -\frac1{a'},\quad
 -\psi'(p'+1),
\]
where \(Q_r(u)=n^{-1}\sum_{i\leq n}(Y_i-u)^r\). At \(\theta_0\) the
expectations of the one-observation versions are the negatives of the
entries of \(J\). The coefficients \(p'\), \(1/a'^2\), \((p'+1)/a'^2\),
\(2/a'^3\), \(1/a'\) and \(\psi'(p'+1)\) are continuous at \(\theta_0\);
for the last, \(\log\Gamma\) is twice continuously differentiable on
\((0,\infty)\) by the differentiation under the integral recorded in that
proof.

\emph{Condition (C).} Fix \(K\) and \(\varepsilon>0\) and put
\(r_n=Kn^{-1/2}\). For large \(n\), \(r_n<\min(a/2,(p+1)/2)\), so every
\(\theta'\in B(\theta_0,r_n)\) has \(a'>0\) and \(p'>-1\). Let
\(D_n=\{\min_iY_i>2r_n\}\). On \(D_n\) the open set of the previous
paragraph contains \(B(\theta_0,r_n)\), since \(|m'-m|\leq r_n<\min_iY_i\)
there, and \(Y_i-u\geq Y_i/2\) whenever \(|u|\leq r_n\). The probability
of \(D_n\) tends to one:
\[
 P(D_n^c)\leq nP(\xi\leq2r_n/a)
 =\frac{n}{\Gamma(k)}\int_0^{2r_n/a}\xi^{k-1}e^{-\xi}\,d\xi
 \leq\frac{n\,(2r_n/a)^k}{k\,\Gamma(k)}
 =O\bigl(n^{1-k/2}\bigr)\to0,
\]
because \(k>2\). On \(D_n\) and for \(|u|\leq r_n\):
\(Q_1(u)=Q_1(0)-u\);
\(|(Y_i-u)^{-1}-Y_i^{-1}|=|u|/(Y_i(Y_i-u))\leq2r_nY_i^{-2}\); and
\(|(Y_i-u)^{-2}-Y_i^{-2}|=|u|\,|2Y_i-u|/(Y_i^2(Y_i-u)^2)
\leq10\,r_nY_i^{-3}\), using \(|2Y_i-u|\leq5Y_i/2\) and
\((Y_i-u)^2\geq Y_i^2/4\). Hence on \(D_n\)
\[
 \sup_{|u|\leq r_n}|Q_{-1}(u)-EY^{-1}|\leq|Q_{-1}(0)-EY^{-1}|
 +2r_nQ_{-2}(0),
\]
\[
 \sup_{|u|\leq r_n}|Q_{-2}(u)-EY^{-2}|\leq|Q_{-2}(0)-EY^{-2}|+10W_n,
 \qquad W_n=\frac{r_n}{n}\sum_{i\leq n}Y_i^{-3}\mathbf 1\{Y_i>2r_n\}.
\]
By \cref{thm:p15-iid-series-strong-law}, \(Q_r(0)\to EY^r\) almost surely
for \(r=-2,-1,1\), the three moments being finite because \(p>1\). For
\(W_n\), with \(c_n=2r_n/a\),
\[
 EW_n=\frac{r_n}{a^3\Gamma(k)}\int_{c_n}^\infty\xi^{k-4}e^{-\xi}\,d\xi
 \leq\frac{r_n}{a^3\Gamma(k)}\Bigl(\int_{c_n}^1\xi^{k-4}\,d\xi
 +\Gamma(k+1)\Bigr),
\]
using \(\xi^{k-4}\leq\xi^k\) for \(\xi\geq1\). The remaining integral is at
most \(1/(k-3)\) when \(k>3\), equals \(\log(1/c_n)\) when \(k=3\), and is
at most \(c_n^{k-3}/(3-k)\) when \(2<k<3\). In each case \(EW_n\to0\):
\(r_n\to0\), \(r_n\log(1/c_n)\to0\), and
\(r_nc_n^{k-3}=O(n^{-1/2}n^{(3-k)/2})=O(n^{(2-k)/2})\to0\). By Markov's
inequality \(W_n\to0\) in probability. Combining with the continuity of
the coefficients, each entry of \(n^{-1}\nabla^2\ell_n(\theta')+J\) is, on
\(D_n\), bounded uniformly over \(B(\theta_0,r_n)\) by a quantity tending to
zero in probability; for instance the \(mm\) entry is at most
\(|p'|\sup_u|Q_{-2}(u)-EY^{-2}|+|p'-p|EY^{-2}\). The operator norm is at
most the sum of the nine absolute entries, so the events on which
\(D_n\) holds and that sum is at most \(\varepsilon\) have probability
tending to one. This is (C), and item~1 follows from item~1 of
\cref{thm:p15-w02b-p15-057-multi}.

\emph{Item 2.} Item~2 of \cref{thm:p15-w02b-p15-057-multi} with \(t\) the
third and the first coordinate vectors gives the limit variances
\((J^{-1})_{pp}\) and \((J^{-1})_{mm}\). The first is
\(2/(2F-2/p+1/p^2)\) by \cref{thm:p15-w02b-p15-057-typeiii}. By
\cref{thm:p15-w02b-p15-057-minor} the second is the determinant of the
\((a,p)\) block, \(((p+1)F-1)/a^2\), divided by
\(\det J=(2F-2/p+1/p^2)/(a^4(p-1))\), which is the stated value. For
\(\tau(m,a,p)=m+a(p+1)\), continuous and differentiable with gradient
\(g=(1,p+1,a)\), observe that \(g=a^2Je_a\), \(e_a\) being the second
coordinate vector, because the second column of \(J\) is
\((1/a^2,(p+1)/a^2,1/a)\). Therefore
\(g^{\top}J^{-1}g=a^4e_a^{\top}Je_a=a^4(p+1)/a^2=a^2(p+1)\).

\emph{Item 3.} By \cref{thm:p15-w02b-p15-057-typeiii}, items~3 and~5,
\(\overline x=m_\mu+a_\mu(p_\mu+1)\), \(E\overline x=m+a(p+1)\) and
\(\operatorname{Var}\overline x=a^2(p+1)/n\). The score has
\(\nabla\ell_n(\theta_0)=\sum_is(X_i)\) almost surely. With
\(T=\overline x\), independence and \(Es(X)=0\) give
\(E[T\nabla\ell_n(\theta_0)]=E[Xs(X)]=E[(X-EX)s(X)]\). Since
\(X-EX=a(\xi-k)=a^2s_a(X)\), with \(s_a=(\xi-k)/a\) the \(a\)-component,
this equals \(a^2E[s\,s_a]=a^2Je_a=g\). The bound of
\cref{thm:p15-w02b-p15-057-multi-bound} is therefore
\(n^{-1}g^{\top}J^{-1}g=a^2(p+1)/n=\operatorname{Var}\overline x\). By
\cref{thm:p15-clt-iid},
\(\sqrt n\,(\overline x-m-a(p+1))\Rightarrow N(0,a^2(p+1))\), so the
efficiency of \cref{def:p15-w02b-p15-057-multi-efficiency} is
\(a^2(p+1)/(a^2(p+1))=1\).

\emph{Item 4.} Divide the value of \((J^{-1})_{mm}\) from item~2 by
\(a^2(p-1)\).

\emph{Item 5.} Let \(x_{(1)}\) be the least observation. Take \(a'=1\),
\(p'=-1/2\) and \(m'<x_{(1)}\). Then
\(\ell_n(m',1,-\tfrac12)=\sum_i\bigl(-\tfrac12\log(x_i-m')-(x_i-m')
-\log\Gamma(\tfrac12)\bigr)\). As \(m'\uparrow x_{(1)}\) each term with
\(x_i=x_{(1)}\) tends to \(+\infty\), and the remaining terms stay bounded
for \(m'\in[x_{(1)}-1,x_{(1)})\). So \(\ell_n\) is unbounded above on
\(\Theta\).
\end{proof}
\end{annexitem}

\begin{annexitem}{annex:p15-w02b-p15-057.moment-shape}{thm:p15-w02b-p15-057-moment-shape}{The moment estimate of the Type III shape}
\begin{proof}
Put \(k=p+1\), \(Z_i=X_i-EX=a\zeta_i\) with \(\zeta_i=\xi_i-k\), and
\(\nu_r=E\zeta^r\). Expanding \(E(\xi-k)^r\) binomially with
\(E\xi^j=k(k+1)\cdots(k+j-1)\) gives
\(\nu_2=k\), \(\nu_3=2k\), \(\nu_4=3k^2+6k\), \(\nu_5=20k^2+24k\) and
\(\nu_6=15k^3+130k^2+120k\); all are finite for every \(k>0\). The
central moments of \(X\) are \(\mu_r=a^r\nu_r\).

\emph{Linearization of the sample moments.} With
\(\overline Z=n^{-1}\sum_iZ_i\),
\(m_2=n^{-1}\sum_iZ_i^2-\overline Z^2\) and, expanding
\((Z_i-\overline Z)^3\) and summing,
\(m_3=n^{-1}\sum_iZ_i^3-3\overline Zn^{-1}\sum_iZ_i^2+2\overline Z^3\).
Let \(U_i=Z_i^2-\mu_2\) and \(V_i=Z_i^3-\mu_3-3\mu_2Z_i\), independent
identically distributed with mean zero and finite variance. Then
\[
 \sqrt n(m_2-\mu_2)=n^{-1/2}\sum_iU_i-\sqrt n\,\overline Z\cdot\overline Z,
\]
\[
 \sqrt n(m_3-\mu_3)=n^{-1/2}\sum_iV_i
 -3\sqrt n\,\overline Z\Bigl(n^{-1}\sum_iZ_i^2-\mu_2\Bigr)
 +2\sqrt n\,\overline Z\cdot\overline Z^2 .
\]
Here \(\sqrt n\,\overline Z\) is tight by Chebyshev's inequality, while
\(\overline Z\to0\) and \(n^{-1}\sum_iZ_i^2\to\mu_2\) almost surely by
\cref{thm:p15-iid-series-strong-law}. A product of a tight sequence and one
tending to zero in probability tends to zero in probability, since
\(P(|A_nB_n|>\rho)\leq P(|A_n|>L)+P(|B_n|>\rho/L)\). So both differences
equal the normalized sums of \(U_i\) and \(V_i\) up to terms tending to
zero in probability, and in particular are tight.

\emph{Delta step.} The function \(G(x,y)=4x^3/y^2-1\) is differentiable at
\((\mu_2,\mu_3)=(a^2k,2a^3k)\), where \(\mu_3\neq0\), with
\(G(\mu_2,\mu_3)=k-1=p\), \(\partial_xG=12\mu_2^2/\mu_3^2=3/a^2\) and
\(\partial_yG=-8\mu_2^3/\mu_3^3=-1/a^3\). Write
\(w_n=(m_2-\mu_2,m_3-\mu_3)\) and
\(G(\mu+w)=p+\nabla G^{\top}w+|w|\,r(w)\) with \(r(w)\to0\) as \(w\to0\).
On the event \(|w_n|<|\mu_3|/2\), whose probability tends to one,
\(m_3\neq0\) and \(p_\mu=G(m_2,m_3)\), so
\[
 \sqrt n(p_\mu-p)=\frac{3}{a^2}\sqrt n(m_2-\mu_2)
 -\frac1{a^3}\sqrt n(m_3-\mu_3)+\sqrt n\,|w_n|\,r(w_n),
\]
and the last term tends to zero in probability as in the proof of
\cref{thm:p15-w02b-p15-057-multi}. The first two terms equal
\(n^{-1/2}\sum_iR_i\) up to terms tending to zero in probability, where
\[
 R_i=\frac{3U_i}{a^2}-\frac{V_i}{a^3}
 =3(\zeta_i^2-k)-(\zeta_i^3-2k-3k\zeta_i)
 =-\zeta_i^3+3\zeta_i^2+3k\zeta_i-k .
\]

\emph{The variance.} \(ER_i=-2k+3k-k=0\). Squaring and using
\(E\zeta=0\),
\[
 ER_i^2=\nu_6-6\nu_5+(9-6k)\nu_4+20k\,\nu_3+(9k^2-6k)\nu_2+k^2 .
\]
Substituting the values of \(\nu_r\) gives
\(6k^3+36k^2+30k=6k(k+1)(k+5)=6(p+1)(p+2)(p+6)\). By
\cref{thm:p15-clt-iid} and \cref{thm:p15-clt-perturbation},
\(\sqrt n(p_\mu-p)\Rightarrow N(0,6(p+1)(p+2)(p+6))\).

\emph{Efficiency.} For \(p>1\), item~2 of
\cref{thm:p15-w02b-p15-057-typeiii-joint} gives the limit variance
\(2/(2F-2/p+1/p^2)\) for the shape coordinate, the inverse-information
value of \cref{def:p15-w02b-p15-057-multi-efficiency} for
\(\tau(m,a,p)=p\); its ratio to \(6(p+1)(p+2)(p+6)\) is the stated
efficiency.
\end{proof}
\end{annexitem}

\begin{annexitem}{annex:p15-w02b-p15-057.trigamma}{thm:p15-w02b-p15-057-trigamma}{Series, integral and two-sided bounds for the trigamma function}
\begin{proof}
\emph{The series.} For \(y>0\) let \(\xi_y\) have the gamma law with shape
\(y\) and unit scale. In the proof of
\cref{thm:p15-w02b-p15-057-typeiii} it is shown that
\(\psi'(y)=\operatorname{Var}\log\xi_y\) and that
\(\psi(s+1)=\psi(s)+1/s\); differentiating the latter,
\(\psi'(s+1)=\psi'(s)-1/s^2\) for \(s>0\). Iterating,
\(\psi'(x)=\sum_{j=0}^{J-1}(x+j)^{-2}+\psi'(x+J)\) for every \(J\). For
\(v>0\) put \(v=e^w\); then \((\log v)^2=w^2\leq(e^{w/2}-e^{-w/2})^2=
(v-1)^2/v\), because \(|e^{w/2}-e^{-w/2}|=2|\sinh(w/2)|\geq|w|\). Hence for
\(y>1\),
\[
 \psi'(y)\leq E\bigl(\log(\xi_y/y)\bigr)^2\leq
 E\frac{(\xi_y-y)^2}{y\,\xi_y}
 =\frac1y\Bigl(E\xi_y-2y+y^2E\xi_y^{-1}\Bigr)
 =\frac1y\Bigl(-y+\frac{y^2}{y-1}\Bigr)=\frac1{y-1},
\]
using \(E\xi_y=y\) and \(E\xi_y^{-1}=1/(y-1)\). So
\(\psi'(x+J)\to0\) as \(J\to\infty\), which gives the series.

\emph{The integral.} For \(u>0\), \(u^{-2}=\int_0^\infty te^{-ut}\,dt\).
Summing over \(u=x+j\) and exchanging sum and integral by Tonelli's
theorem, with \(\sum_{j\geq0}e^{-jt}=1/(1-e^{-t})\) for \(t>0\), gives the
integral.

\emph{Bounds for \(u\coth u\).} For \(u>0\) compare the power series of
\(P(u)\sinh u\) and \(u\cosh u\), where
\(P(u)=1+u^2/3-u^4/45+\lambda\,2u^6/945\) with \(\lambda\in\{0,1\}\). The
coefficient of \(u^{2k+1}\) in \(P(u)\sinh u-u\cosh u\), multiplied by
\((2k+1)!\), is, writing \((m)_j=m(m-1)\cdots(m-j+1)\),
\[
 c_\lambda(k)=1+\frac{(2k+1)_2}3-\frac{(2k+1)_4}{45}
 +\lambda\frac{2(2k+1)_6}{945}-(2k+1),
\]
with \((m)_j=0\) when \(j>m\). Direct expansion gives
\[
 c_0(k)=-\frac{16}{45}k(k-1)(k-2)(k+2),\qquad
 c_1(k)=\frac{16}{945}k(k-1)(k-2)(k-3)(8k^2+12k+13),
\]
for every integer \(k\geq0\). (For \(c_1\), factor
\(16k(k-1)(k-2)/945\) from both terms to get
\((4k^2-1)(2k-3)-21(k+2)=8k^3-12k^2-23k-39=(k-3)(8k^2+12k+13)\).) Thus
\(c_0(k)\leq0\) for all \(k\), with strict inequality for \(k\geq3\), and
\(c_1(k)\geq0\) for all \(k\), with strict inequality for \(k\geq4\). Since
\(u^{2k+1}>0\), dividing by \(\sinh u>0\),
\[
 1+\frac{u^2}3-\frac{u^4}{45}<u\coth u<1+\frac{u^2}3-\frac{u^4}{45}
 +\frac{2u^6}{945}\qquad(u>0).
\]

\emph{The bounds for \(\psi'\).} For \(t>0\),
\(t/(1-e^{-t})=t/2+(t/2)\coth(t/2)\), because
\(te^{t/2}/(e^{t/2}-e^{-t/2})=(t/2)(\cosh(t/2)+\sinh(t/2))/\sinh(t/2)\).
With \(u=t/2\) the bounds read
\[
 1+\frac t2+\frac{t^2}{12}-\frac{t^4}{720}<\frac t{1-e^{-t}}<
 1+\frac t2+\frac{t^2}{12}-\frac{t^4}{720}+\frac{t^6}{30240}.
\]
Multiply by \(e^{-xt}\) and integrate, using
\(\int_0^\infty t^je^{-xt}\,dt=j!/x^{j+1}\): the terms give \(1/x\),
\(1/(2x^2)\), \(2/(12x^3)=1/(6x^3)\), \(24/(720x^5)=1/(30x^5)\) and
\(720/(30240x^7)=1/(42x^7)\). The inequalities between the integrands are
strict on \((0,\infty)\), so those between the integrals are strict.

\emph{Consequences.} For \(p>0\),
\(F(p)=\psi'(p+1)=\psi'(p)-1/p^2\). Subtracting \(1/p^2\) from the bounds at
\(x=p\) gives
\(F(p)-\bigl(\frac1p-\frac1{2p^2}+\frac1{6p^3}-\frac1{30p^5}\bigr)\in
(0,\frac1{42p^7})\), which is the stated form with
\(\vartheta\in(0,1)\). Finally
\(F(p)=\sum_{j\geq1}(p+j)^{-2}>\sum_{j\geq1}\int_j^{j+1}(p+t)^{-2}\,dt
=1/(p+1)\).
\end{proof}
\end{annexitem}

\begin{annexitem}{annex:p15-w02b-p15-057.typeiii-largep}{thm:p15-w02b-p15-057-typeiii-largep}{Fisher's large-\(p\) expansions with remainders, and the 80 per cent threshold}
\begin{proof}
Write \(F=F(p)=\frac1p-\frac1{2p^2}+\frac1{6p^3}-\frac1{30p^5}+\frac{\vartheta}{42p^7}\)
by \cref{thm:p15-w02b-p15-057-trigamma}, with \(p\geq1\). The
interpretations of \(2/D\) and \(n'\) for \(p>1\) are items~2 and~4 of
\cref{thm:p15-w02b-p15-057-typeiii-joint}.

\emph{Item 1.} \(2F-\frac2p+\frac1{p^2}=\frac1{3p^3}-\frac1{15p^5}
+\frac{\vartheta}{21p^7}\), which is the first display. Multiplying out,
\[
 (p+1)F-1=\frac1{2p}-\frac1{3p^2}+\frac1{6p^3}-\frac1{30p^4}-\frac1{30p^5}
 +\frac{\vartheta(p+1)}{42p^7},
\]
the constant terms \(1\) and \(-1\) cancelling and the \(p^{-1}\) terms
being \(\frac1p-\frac1{2p}\); factoring \(\frac1{2p}\) gives the second
display.

\emph{Item 2.} Write \(D=\frac1{3p^3}(1-\varepsilon)\) with
\(\varepsilon=\frac1{5p^2}-\eta\), \(\eta=\frac{\vartheta}{7p^4}\in
(0,\frac1{7p^4})\). For \(p\geq1\), \(0<\varepsilon<\frac1{5p^2}\leq\frac15\).
Then \(2/D-6p^3-\frac65p=6p^3\frac{\varepsilon}{1-\varepsilon}-\frac65p\).
As \(\varepsilon/(1-\varepsilon)\) increases with \(\varepsilon\), it is
less than \(\frac1{5p^2-1}\), and the difference is less than
\(\frac{6p^3}{5p^2-1}-\frac{6p}5=\frac{6p}{5(5p^2-1)}\), which is less than
\(\frac6{7p}\) because \(42p^2<150p^2-30\) for \(p\geq1\). Also
\(\varepsilon/(1-\varepsilon)>\varepsilon>\frac1{5p^2}-\frac1{7p^4}\), so
the difference exceeds \(\frac{6p}5-\frac6{7p}-\frac{6p}5=-\frac6{7p}\).

\emph{Item 3.} By \cref{thm:p15-w02b-p15-057-moment-shape},
\(E(p)=(2/D)/\bigl(6(p+1)(p+2)(p+6)\bigr)\), and the approximation is
\(6(p^3+\frac p5)/\bigl(6(p+1)(p+2)(p+6)\bigr)\). Divide item~2 by
\(6(p+1)(p+2)(p+6)\).

\emph{Item 4.} Put \(E_F(p)=(p^3+\frac p5)/((p+1)(p+2)(p+6))\) and
\(g(p)=p^3-36p^2-79p-48\). Since
\((p+1)(p+2)(p+6)=p^3+9p^2+20p+12\),
\(p^3+\frac p5-0.8(p^3+9p^2+20p+12)=0.2\,g(p)\), so by item~3
\[
 (E(p)-0.8)(p+1)(p+2)(p+6)\in\Bigl(0.2\,g(p)-\frac1{7p},\;
 0.2\,g(p)+\frac1{7p}\Bigr).
\]
Hence \(E(p)<0.8\) when \(g(p)<-\frac5{7p}\), and \(E(p)>0.8\) when
\(g(p)>\frac5{7p}\). Now \(g'(p)=3p^2-72p-79\) has its positive zero at
\(p_0=(72+\sqrt{6132})/6\), which lies between \(25\) and \(26\) because
\(g'(25)=-4<0<g'(26)=77\); so \(g\) decreases on \([1,p_0]\) and increases
on \([p_0,\infty)\). On \([1,38.1]\), therefore,
\(g\leq\max(g(1),g(38.1))\), and \(g(1)=-162\),
\(g(38.1)=55306.341-52257.96-3009.9-48=-9.519\); both are below
\(-\frac57\leq-\frac5{7p}\). For \(p\geq38.11\),
\(g(p)\geq g(38.11)=55349.900731-52285.3956-3010.69-48=5.815131\), which
exceeds \(\frac5{7p}\). (The products are exact decimal arithmetic:
\(38.1^2=1451.61\), \(38.1^3=55306.341\), \(38.11^2=1452.3721\),
\(38.11^3=55349.900731\).) Finally \(4/39.1=0.10230\ldots\).

\emph{Item 5.} \(D>0\) for \(p\geq1\) by item~1, since
\(\frac1{3p^3}-\frac1{15p^5}>0\). Then
\[
 n'-1=\frac{\mathcal N-D}D=\frac{(p-1)F-(p-1)^2/p^2}{D}
 =\frac{(p-1)\bigl(F-(p-1)/p^2\bigr)}{D},
\]
because \(\mathcal N-D=(p-1)F-1+\frac2p-\frac1{p^2}\). This is zero at \(p=1\). For
\(p>1\), \(F>\frac1{p+1}>\frac{p-1}{p^2}\), the last since
\(p^2>(p-1)(p+1)\); so \(n'>1\). For the expansion put
\(A_0=\frac32p^2-p+\frac45>0\) and note \(n'=3p^3\mathcal N/(1-\varepsilon)\) with
\(\varepsilon\) as in item~2. By item~1,
\(3p^3\mathcal N=\frac32p^2-p+\frac12-\frac1{10p}-\frac1{10p^2}
+\frac{\vartheta(p+1)}{14p^4}\), and
\(\varepsilon A_0=\frac3{10}-\frac1{5p}+\frac4{25p^2}-\eta A_0\). Hence
\[
 (1-\varepsilon)(n'-A_0)=3p^3\mathcal N-A_0+\varepsilon A_0
 =-\frac3{10p}+\frac3{50p^2}+\frac{\vartheta(p+1)}{14p^4}-\eta A_0 .
\]
For \(p\geq1\): \(0<\frac{\vartheta(p+1)}{14p^4}<\frac{2p}{14p^4}\leq
\frac1{7p}\) and \(0<\eta A_0<\frac{1}{7p^4}\bigl(\frac32p^2+\frac45\bigr)
\leq\frac{23}{70p}\). So the right side lies between
\(-\frac3{10p}-\frac{23}{70p}=-\frac{22}{35p}\) and
\(\bigl(-\frac3{10}+\frac3{50}+\frac17\bigr)\frac1p=-\frac{34}{350p}<0\).
Since \(\frac45\leq1-\varepsilon<1\), \(n'-A_0\) is negative and greater
than \(-\frac{22}{35p}\cdot\frac54=-\frac{11}{14p}\).

\emph{Item 6.} Put \(q=\frac23n'\) and \(b=p-\frac13>0\). Then
\(\frac23A_0=p^2-\frac23p+\frac8{15}\) and \(b^2=p^2-\frac23p+\frac19\), so
by item~5
\(q-b^2=\frac{19}{45}+\frac23(n'-A_0)\in\bigl(\frac{19}{45}-\frac{11}{21p},
\frac{19}{45}\bigr)\). For \(p\geq2\) the lower end exceeds
\(\frac{19}{45}-\frac{11}{42}>0\). So \(\sqrt q>b\) and
\(0<\sqrt q-b=\frac{q-b^2}{\sqrt q+b}<\frac{19}{45\cdot2b}\). Multiply by
\(\sqrt{3/2}\), noting \(\sqrt{n'}=\sqrt{3/2}\sqrt q\).
\end{proof}
\end{annexitem}

\begin{annexitem}{annex:p15-w02b-p15-057.location-scale}{thm:p15-w02b-p15-057-location-scale}{Location--scale information and the centre of location}
\begin{proof}
\emph{Moment identities.} Since \(g'=\phi'g\) and
\(g''=(\phi''+\phi'^2)g\), integration on \(D\) with the stated limits and
integrability gives
\(E\phi'=\int g'=0\);
\(E\xi\phi'=\int\xi g'=[\xi g]-\int g=-1\);
\(E\phi''+E\phi'^2=\int g''=[g']=0\);
\(E\xi\phi''+E\xi\phi'^2=\int\xi g''=[\xi g']-\int g'=0\); and
\(E\xi^2\phi''+E\xi^2\phi'^2=\int\xi^2g''=[\xi^2g']-2\int\xi g'=2\).
Here \([\cdot]\) is the difference of the end limits, all zero, and
integration by parts on compact subintervals followed by the limits is
justified by the stated integrability.

\emph{Information matrix.} With \(\xi=(x-m)/a\), \(\partial_m\xi=-1/a\)
and \(\partial_a\xi=-\xi/a\), so
\(\partial_m\log f=-\phi'(\xi)/a\) and
\(\partial_a\log f=-(1+\xi\phi'(\xi))/a\). The scores are square integrable
by the finiteness of \(E\phi'^2\) and \(E\xi^2\phi'^2\). Their second
moments are \(E\phi'^2/a^2=-E\phi''/a^2\);
\(E[\phi'(1+\xi\phi')]/a^2=E\xi\phi'^2/a^2=-E\xi\phi''/a^2\); and
\(E(1+\xi\phi')^2/a^2=(1-2+E\xi^2\phi'^2)/a^2=(1-E\xi^2\phi'')/a^2\).
The second derivatives are
\(\partial_m\partial_a\log f=a^{-2}(\phi'+\xi\phi'')\) and
\(\partial_a^2\log f=a^{-2}(1+2\xi\phi'+\xi^2\phi'')\); their sample sums
are Fisher's expressions \(a^{-2}S(\phi'+\xi\phi'')\) and
\(a^{-2}S(2\xi\phi'+\xi^2\phi'')+n/a^2\). Replacing \(\phi'\) by its mean
zero and \(\xi\phi'\) by its mean \(-1\) gives his shortened forms; this is
valid for expectations and not for the sums themselves.

\emph{Centre of location.} The density \(g_c(\eta)=g(\eta+c)\) on
\(D-c\) has \(\phi_c(\eta)=\phi(\eta+c)\), and each hypothesis for
\(g_c\) is a linear combination of the hypotheses for \(g\) with powers at
most two, so the matrix formula applies to \(g_c\). Its off-diagonal entry
is \(-E[(\xi-c)\phi''(\xi)]/a^2=-(E\xi\phi''-cE\phi'')/a^2\), zero for
\(c=E\xi\phi''/E\phi''\). The family \(a^{-1}g_c((x-m')/a)\) is the
original one with \(m'=m+ac\).

\emph{Cauchy form.} With \(\xi=\tan t\), \(t\in(-\pi/2,\pi/2)\),
\(d\xi=\sec^2t\,dt\) and \(1+\xi^2=\sec^2t\):
\[
 \frac1\pi\int_{\mathbb R}\frac{1-\xi^2}{(1+\xi^2)^3}\,d\xi
 =\frac1\pi\int_{-\pi/2}^{\pi/2}\cos^2t\cos2t\,dt
 =\frac1\pi\int_{-\pi/2}^{\pi/2}\frac{\cos2t+\cos^22t}{2}\,dt
 =\frac14,
\]
as \((1-\tan^2t)\cos^4t=\cos^2t\cos2t\). Hence
\(E\phi''=-2\cdot\tfrac14=-\tfrac12\), the value printed on p.~339.
The hypotheses hold: \(g,\xi g,g',\xi g',\xi^2g'\) are
\(O(|\xi|^{-1})\) at \(\pm\infty\); \(g',\xi g',g'',\xi g'',\xi^2g''\) are
\(O((1+\xi^2)^{-1})\); and \(\phi'^2\), \(\xi^2\phi'^2\), \(|\phi''|\),
\(\xi^2|\phi''|\) are bounded, because \(|\phi'|\leq2|\xi|/(1+\xi^2)\) and
\(|\phi''|\leq2/(1+\xi^2)\).

\emph{Type III form.} Here \(g=\xi^pe^{-\xi}/\Gamma(p+1)\) on
\((0,\infty)\), \(\phi''=-p/\xi^2\). For \(p>1\),
\(g'=(p/\xi-1)g\) behaves as a constant times \(\xi^{p-1}\) at zero and
\(g''=((p/\xi-1)^2-p/\xi^2)g\) as a constant times \(\xi^{p-2}\), which is
integrable at zero exactly because \(p>1\); at infinity all the listed
functions decay exponentially. The moments in the hypotheses are finite
because \(E\xi^{-2}<\infty\) for \(p>1\). Using the moments of the
previous proof, \(E\phi''=-pE\xi^{-2}=-1/(p-1)\), \(E\xi\phi''=-pE\xi^{-1}
=-1\), and \(E\xi^2\phi''=-p\). Then \(c=p-1\), and after the shift
\(E(\xi-c)^2\phi''=-p+2(p-1)-(p-1)=-1\), so \(1-E(\xi-c)^2\phi''=2\).
\end{proof}
\end{annexitem}

\begin{annexitem}{annex:p15-w02b-p15-057.typeiv}{thm:p15-w02b-p15-057-typeiv}{Type IV information and the efficiency of moments}
\begin{proof}
\emph{Angular coordinates.} Put \(\xi=\tan\vartheta\) with
\(\vartheta\in(-\pi/2,\pi/2)\). Then \(1+\xi^2=\sec^2\vartheta\),
\(d\xi=\sec^2\vartheta\,d\vartheta\) and
\(g(\xi)\,d\xi=C\,e^{-\nu\vartheta}\cos^r\vartheta\,d\vartheta\). For
\(k>-1\) put
\(I_k=\int_{-\pi/2}^{\pi/2}e^{-\nu\vartheta}\cos^k\vartheta\,d\vartheta\),
finite because \(\cos^k\vartheta\) is of order
\((\pi/2-|\vartheta|)^k\) at the ends, and
\(K_k=\int_{-\pi/2}^{\pi/2}e^{-\nu\vartheta}\cos^k\vartheta\sin\vartheta
\,d\vartheta\). Then \(C=1/I_r\) and \(Eh(\xi)=I_r^{-1}\int
h(\tan\vartheta)e^{-\nu\vartheta}\cos^r\vartheta\,d\vartheta\) whenever
either side is absolutely convergent.

\emph{Two recursions.} For \(k>-1\), integrating by parts on
\([-\pi/2+\delta,\pi/2-\delta]\) with
\(d(\cos^{k+1}\vartheta)=-(k+1)\cos^k\vartheta\sin\vartheta\,d\vartheta\)
and letting \(\delta\to0\), the boundary term vanishing because
\(k+1>0\), gives \(-(k+1)K_k=\nu I_{k+1}\), that is
\(K_k=-\nu I_{k+1}/(k+1)\). For \(k>1\),
\(\frac{d}{d\vartheta}(\cos^{k-1}\vartheta\sin\vartheta)
=k\cos^k\vartheta-(k-1)\cos^{k-2}\vartheta\), since
\(\sin^2=1-\cos^2\). Integrating this against \(e^{-\nu\vartheta}\) and
by parts, the boundary term vanishing because \(k-1>0\),
\(kI_k-(k-1)I_{k-2}=\nu K_{k-1}=-\nu^2I_k/k\). Hence
\[
 I_k=\frac{k(k-1)}{k^2+\nu^2}\,I_{k-2}\qquad(k>1).
\]

\emph{Item 1: hypotheses.} Here
\(\phi'=-(\nu+(r+2)\xi)/(1+\xi^2)\) and
\(\phi''=\bigl((r+2)(\xi^2-1)+2\nu\xi\bigr)/(1+\xi^2)^2\), so there is a
constant \(c_0\) with \(|\phi'|\leq c_0/(1+|\xi|)\) and
\(|\phi''|\leq c_0/(1+\xi^2)\). Also \(g\leq Ce^{\pi|\nu|/2}
(1+\xi^2)^{-(r+2)/2}\), \(g'=\phi'g\) and \(g''=(\phi''+\phi'^2)g\). For
\(j\leq2\), \(\xi^jg'\) is \(O(|\xi|^{j-r-3})\) and \(\xi^jg''\) is
\(O(|\xi|^{j-r-4})\) at \(\pm\infty\); since \(r>-1\), the first tends to
zero and the second is integrable. Likewise \(g\) and \(\xi g\) tend to
zero, and \(g'\), \(\xi g'\) are integrable. The functions \(\phi'^2\),
\(\xi^2\phi'^2\), \(\phi''\) and \(\xi^2\phi''\) are bounded, so the moments
in the hypotheses are finite.

\emph{Item 1: values.} In angular coordinates
\(1/(1+\xi^2)=\cos^2\vartheta\), \(\xi/(1+\xi^2)=\sin\vartheta\cos\vartheta\),
\(1/(1+\xi^2)^2=\cos^4\vartheta\),
\(\xi/(1+\xi^2)^2=\sin\vartheta\cos^3\vartheta\) and
\(\xi^2/(1+\xi^2)^2=\cos^2\vartheta-\cos^4\vartheta\). Hence
\[
\begin{aligned}
 \phi''&=(r+2)(\cos^2\vartheta-2\cos^4\vartheta)
  +2\nu\sin\vartheta\cos^3\vartheta,\\
 \xi\phi''&=(r+2)(\sin\vartheta\cos\vartheta
  -2\sin\vartheta\cos^3\vartheta)
  +2\nu(\cos^2\vartheta-\cos^4\vartheta),\\
 \xi^2\phi''&=(r+2)(1-3\cos^2\vartheta+2\cos^4\vartheta)
  +2\nu(\sin\vartheta\cos\vartheta-\sin\vartheta\cos^3\vartheta),
\end{aligned}
\]
the last because, with \(w=\cos^2\vartheta\),
\(\xi^2(\xi^2-1)w^2=(1-w)(1-2w)\) and \(\xi^3w^2=\xi w(1-w)\). Put
\(D_2=(r+2)^2+\nu^2\) and \(D=(r+4)^2+\nu^2\). The recursions give
\[
\begin{aligned}
 E\cos^2\vartheta&=\frac{I_{r+2}}{I_r}=\frac{(r+1)(r+2)}{D_2},&
 E\cos^4\vartheta&=\frac{I_{r+4}}{I_r}=\frac{(r+1)(r+2)(r+3)(r+4)}{D_2D},\\
 E\sin\vartheta\cos\vartheta&=\frac{K_{r+1}}{I_r}=-\frac{\nu(r+1)}{D_2},&
 E\sin\vartheta\cos^3\vartheta&=\frac{K_{r+3}}{I_r}
 =-\frac{\nu(r+1)(r+2)(r+3)}{D_2D}.
\end{aligned}
\]
Substituting,
\(E\phi''=\frac{(r+1)(r+2)}{D_2D}\bigl[(r+2)D-2(r+2)(r+3)(r+4)
-2\nu^2(r+3)\bigr]\), and the bracket equals
\((r+2)(r+4)(-r-2)-\nu^2(r+4)=-(r+4)D_2\); so
\(E\phi''=-(r+1)(r+2)(r+4)/D\). Next
\(E\xi\phi''=\frac{\nu(r+1)(r+2)}{D_2D}\bigl[-D+2(r+2)(r+3)+2D
-2(r+3)(r+4)\bigr]=\frac{\nu(r+1)(r+2)}{D_2D}\bigl[D-4(r+3)\bigr]\), and
\(D-4(r+3)=D_2\). Finally
\(E\xi^2\phi''-1=(r+1)+\frac{(r+1)}{D_2D}\,Q\) with
\[
 Q=(r+2)^2\bigl[2(r+3)(r+4)-3D\bigr]+2\nu^2\bigl[(r+2)(r+3)-D\bigr]
 =-(r+2)^2\bigl((r+4)(r+6)+3\nu^2\bigr)-2\nu^2(3r+10+\nu^2),
\]
and \(D_2D+Q=-2(r+2)^2(r+4)-\nu^2(r^2+6r+12)-\nu^4
=-\bigl(2(r+4)+\nu^2\bigr)D_2\), using
\(D_2D=(r+2)^2(r+4)^2+\nu^2\bigl((r+2)^2+(r+4)^2\bigr)+\nu^4\) and
\((r+2)^2+2(r+4)=r^2+6r+12\). So
\(E\xi^2\phi''-1=-(r+1)(2(r+4)+\nu^2)/D\). The centre of location of
\cref{thm:p15-w02b-p15-057-location-scale} is
\(c=E\xi\phi''/E\phi''=-\nu/(r+4)\). The derivative \(\phi'\) has the
single zero \(\xi=-\nu/(r+2)\) and is positive to its left and negative
to its right, so that point is the mode.

\emph{Item 2.} By \cref{thm:p15-w02b-p15-057-location-scale} the
information for \(m\) is \(-E\phi''/a^2\). For a symmetric \(2\times2\)
matrix with diagonal entries \(A,A'\), off-diagonal entry \(B\) and
\(A>0\), the \((2,2)\) entry of the inverse is \(1/(A'-B^2/A)\), and the
matrix is positive definite when also \(A'-B^2/A>0\). Here
\[
 a^2\Bigl(A'-\frac{B^2}A\Bigr)
 =1-E\xi^2\phi''+\frac{(E\xi\phi'')^2}{E\phi''}
 =\frac{r+1}{(r+4)D}\bigl[(2(r+4)+\nu^2)(r+4)-\nu^2(r+2)\bigr]
 =\frac{2(r+1)}{r+4},
\]
since the bracket is \(2(r+4)^2+2\nu^2=2D\). For
\cref{def:p15-w02b-p15-057-multi-regular}: the support is \(\mathbb R\) for
every parameter; \(\log f_{m,a}(x)=\log C-\log a+\phi((x-m)/a)\) is
infinitely differentiable in \((m,a)\) for \(a>0\); with
\(\eta=(x-m)/a\) the scores are \(-\phi'(\eta)/a\) and
\(-(1+\eta\phi'(\eta))/a\), and the second derivatives are
\(a^{-2}\phi''(\eta)\), \(a^{-2}(\phi'+\eta\phi'')(\eta)\) and
\(a^{-2}(1+2\eta\phi'+\eta^2\phi'')(\eta)\). The functions \(\phi''\),
\(\phi'+\eta\phi''\) and \(1+2\eta\phi'+\eta^2\phi''\) are bounded on
\(\mathbb R\) by the bounds above, so for parameters with \(a'\geq a/2\)
every entry of the Hessian of \(\log f\) is bounded by a constant, which
serves as \(M\). Because \(\nabla f=f\,s\) and
\(\nabla^2f=f(\nabla^2\log f+ss^{\top})\) with bounded factors, the two
integrals \(\int\nabla f\) and \(\int\nabla^2f\) equal the expectation of
the score, zero by \cref{thm:p15-w02b-p15-057-location-scale}, and the
expected Hessian plus the information matrix, which is zero by the
expectations computed in the proof of that theorem. Hence
\cref{thm:p15-w02b-p15-057-multi-regular} applies. Holding \(a\) fixed,
the same bounds give (R1)--(R5) for \(m\).

\emph{Item 3.} Since \(\tan\vartheta\cos^r\vartheta=\sin\vartheta
\cos^{r-1}\vartheta\), \(E|\xi|<\infty\) exactly when \(r>0\), and then
\(E\xi=K_{r-1}/I_r=-\nu/r\). Since
\(\tan^2\vartheta=\sec^2\vartheta-1\), \(E\xi^2<\infty\) exactly when
\(r>1\), and then \(E\xi^2=I_{r-2}/I_r-1=(r^2+\nu^2)/(r(r-1))-1
=(\nu^2+r)/(r(r-1))\), so
\(\operatorname{Var}\xi=(\nu^2+r)/(r(r-1))-\nu^2/r^2=(r^2+\nu^2)/(r^2(r-1))\).
Thus \(E\overline x=m-a\nu/r\), \(m_\mu\) is unbiased, and
\(\operatorname{Var}m_\mu=a^2\operatorname{Var}\xi/n\). Its finite-sample
information efficiency is
\(1/\bigl(n\cdot\frac{(r+1)(r+2)(r+4)}{a^2D}\cdot\frac{a^2(r^2+\nu^2)}
{nr^2(r-1)}\bigr)\), the displayed expression. By
\cref{thm:p15-clt-iid}, \(\sqrt n\,(m_\mu-m)\) converges in law to
\(N(0,a^2\operatorname{Var}\xi)\), and the model is regular in \(m\) by
item~2, so the efficiency of
\cref{def:p15-w02b-p15-057-efficiency} is the same number. At \(\nu=0\) it
is \((r-1)(r+4)/((r+1)(r+2))=1-6/((r+1)(r+2))\). At \(r=7\) it is
\(\frac{49\cdot6\,(121+\nu^2)}{8\cdot9\cdot11\,(49+\nu^2)}
=\frac{49}{132}\cdot\frac{121+\nu^2}{49+\nu^2}\), and
\((121+x)/(49+x)=1+72/(49+x)\) decreases from \(121/49\) to \(1\) as
\(x=\nu^2\) increases.

\emph{Item 4.} Let \(r>3\). Two steps of the recursion give
\(E\sec^4\vartheta=I_{r-4}/I_r=\frac{(r^2+\nu^2)((r-2)^2+\nu^2)}
{r(r-1)(r-2)(r-3)}\), and
\(E\tan\vartheta\sec^2\vartheta=K_{r-3}/I_r=-\nu I_{r-2}/((r-2)I_r)
=-\frac{\nu(r^2+\nu^2)}{r(r-1)(r-2)}\). Since
\(\xi^3=\tan\vartheta\sec^2\vartheta-\tan\vartheta\) and
\(\xi^4=\sec^4\vartheta-2\sec^2\vartheta+1\), all four raw moments of
\(\xi\) are finite and explicit, and the central fourth moment
\(E\xi^4-4E\xi E\xi^3+6(E\xi)^2E\xi^2-3(E\xi)^4\) equals
\[
 \frac{3(r^2+\nu^2)\bigl(r^3-2r^2+\nu^2r+6\nu^2\bigr)}
 {r^4(r-1)(r-2)(r-3)} .
\]
Dividing by \((\operatorname{Var}\xi)^2\) gives
\(\beta_2=\frac{3(r-1)(r^3-2r^2+\nu^2r+6\nu^2)}{(r-2)(r-3)(r^2+\nu^2)}\),
which is the displayed form because
\((r+6)(r^2+\nu^2)-8r^2=r^3-2r^2+\nu^2r+6\nu^2\). Then
\((\beta_2-1)(r-2)(r-3)(r^2+\nu^2)
=3(r-1)(r^3-2r^2+\nu^2r+6\nu^2)-(r-2)(r-3)(r^2+\nu^2)\); its
\(\nu\)-free part is \(r^2(r-2)\bigl(3(r-1)-(r-3)\bigr)=2r^3(r-2)\) and its
\(\nu^2\) coefficient is \(3(r-1)(r+6)-(r-2)(r-3)=2(r^2+10r-12)\), which
gives the displayed \((\beta_2-1)/4\). Write
\(\kappa=\operatorname{Var}\xi\), \(\mu_2=a^2\kappa\) and
\(\mu_4=\beta_2\mu_2^2\) for the central moments of one observation. As
in the proof of item~4 of \cref{thm:p15-w02b-p15-057-typeiii},
\(\sqrt n\,(m_2-\mu_2)\Rightarrow N(0,\mu_4-\mu_2^2)\), with
\(\mu_4-\mu_2^2>0\) because \((X-EX)^2\) is not almost surely constant, and
\(m_2\to\mu_2\) almost surely; and
\(\sqrt n\,(a_\mu-a)=\sqrt n\,(m_2-\mu_2)/\bigl(\sqrt\kappa
(\sqrt{m_2}+\sqrt{\mu_2})\bigr)\), whose denominator tends to
\(2\sqrt{\kappa\mu_2}=2a\kappa\). By \cref{thm:p15-clt-perturbation} the
limit is normal with variance
\((\beta_2-1)\mu_2^2/(4a^2\kappa^2)=a^2(\beta_2-1)/4\). By item~2 the
inverse-information value for \(\tau(m,a)=a\) is \(a^2(r+4)/(2(r+1))\),
and (S) holds by \cref{thm:p15-w02b-p15-057-multi-regular}; the ratio is
the displayed efficiency. At \(\nu=0\) it is
\((r-3)(r+4)/(r(r+1))=1-12/(r(r+1))\). At \(r=7\) it is
\(\frac{5\cdot4\cdot11\,(49+\nu^2)}{8\,(1715+107\nu^2)}
=\frac{55}2\cdot\frac{49+\nu^2}{1715+107\nu^2}\), whose derivative in
\(x=\nu^2\) has the sign of \(1715-49\cdot107<0\); the values at
\(x=0\) and as \(x\to\infty\) are \(2695/3430\) and \(55/214\).
\end{proof}
\end{annexitem}

\begin{annexitem}{annex:p15-w02b-p15-057.pearson-moments}{thm:p15-w02b-p15-057-pearson-moments}{Moments of the Type IV curve and the boundaries they impose}
\begin{proof}
\emph{Item 1.} The factor \(e^{-\nu\tan^{-1}\xi}\) lies between
\(e^{-|\nu|\pi/2}\) and \(e^{|\nu|\pi/2}\), and
\(2^{-(r+2)/2}|\xi|^{-(r+2)}\leq(1+\xi^2)^{-(r+2)/2}\leq|\xi|^{-(r+2)}\)
for \(|\xi|\geq1\). So \(|\xi|^kg(\xi)\) lies between positive multiples of
\(|\xi|^{k-r-2}\) for \(|\xi|\geq1\) and is bounded on \([-1,1]\). Since
\(\int_1^\infty t^{k-r-2}dt<\infty\) if and only if \(k-r-2<-1\), that is
\(k<r+1\), item~1 follows.

\emph{Item 2.} Differentiating,
\((1+\xi^2)g'(\xi)=-(\nu+(r+2)\xi)g(\xi)\). Let \(0\leq k<r\). For
\(R>0\), integration by parts gives
\[
 \int_{-R}^R\xi^k(1+\xi^2)g'\,d\xi
 =\Bigl[\xi^k(1+\xi^2)g\Bigr]_{-R}^R
 -\int_{-R}^R\bigl(k\xi^{k-1}+(k+2)\xi^{k+1}\bigr)g\,d\xi .
\]
By the bounds of item~1, \(|\xi|^k(1+\xi^2)g(\xi)\) is at most a constant
times \(|\xi|^{k-r}\) for \(|\xi|\geq1\), which tends to \(0\) since
\(k<r\). The integrands \(\xi^{k-1}g\), \(\xi^kg\), \(\xi^{k+1}g\) are
integrable by item~1 since \(k+1<r+1\). Letting \(R\to\infty\) and
substituting the differential equation,
\(-\nu\mu_k'-(r+2)\mu_{k+1}'=-k\mu_{k-1}'-(k+2)\mu_{k+1}'\), which is
\((r-k)\mu_{k+1}'=k\mu_{k-1}'-\nu\mu_k'\).

\emph{Item 3.} The variance of a mean of \(n\) observations is
\(\operatorname{Var}\xi/n\), finite exactly when \(E\xi^2<\infty\), that is
\(r>1\) by item~1. The central fourth moment
\(E(\xi-E\xi)^4\) is finite exactly when \(E\xi^4<\infty\), that is
\(r>3\); for \(1<r\leq3\) the variance is finite and positive while the
fourth central moment is infinite, so \(\beta_2\) is infinite. The
variables \(\xi_i^j\) are independent with the law of \(\xi^j\), so
\(n^{-1}\sum\xi_i^j\) has variance \(\operatorname{Var}(\xi^j)/n\) when
\(E\xi^{2j}<\infty\) and infinite variance when \(E|\xi|^j<\infty\) but
\(E\xi^{2j}=\infty\); by item~1 the first case is \(2j<r+1\). For \(j=2\)
this is \(r>3\) and for \(j=4\) it is \(r>7\).

\emph{Item 4.} Let \(r>3\). Item~2 with \(k=0,1,2,3\) gives
\[
 \mu_1'=-\frac\nu r,\quad
 \mu_2'=\frac{\nu^2+r}{r(r-1)},\quad
 \mu_3'=-\frac{\nu(\nu^2+3r-2)}{r(r-1)(r-2)},\quad
 \mu_4'=\frac{\nu^4+6\nu^2r-8\nu^2+3r^2-6r}{r(r-1)(r-2)(r-3)} .
\]
The central moments \(c_j=E(\xi-\mu_1')^j\) follow by expansion:
\[
 c_2=\frac{r^2+\nu^2}{r^2(r-1)},\quad
 c_3=-\frac{4\nu(r^2+\nu^2)}{r^3(r-1)(r-2)},\quad
 c_4=\frac{3(r^2+\nu^2)\bigl(r^3-2r^2+\nu^2(r+6)\bigr)}{r^4(r-1)(r-2)(r-3)} .
\]
Then \(\beta_1=c_3^2/c_2^3=16\nu^2(r-1)/((r-2)^2(r^2+\nu^2))\), and
\(\beta_2=c_4/c_2^2\) is the expression of item~4 of
\cref{thm:p15-w02b-p15-057-typeiv}, since
\((r+6)(r^2+\nu^2)-8r^2=r^3-2r^2+\nu^2(r+6)\). Put
\(W=\nu^2+(r-2)^2>0\) and \(S=r^2+\nu^2\). Direct reduction to a common
denominator gives
\[
 \beta_2+3=\frac{6r^2W}{S(r-2)(r-3)},\quad
 4\beta_2-3\beta_1=\frac{12r^2(r-1)W}{S(r-2)^2(r-3)},\quad
 2\beta_2-3\beta_1-6=\frac{12r^2W}{S(r-2)^2(r-3)},
\]
all nonzero for \(r>3\). Hence
\[
 \frac{\beta_1(\beta_2+3)^2}{4(4\beta_2-3\beta_1)(2\beta_2-3\beta_1-6)}
 =\frac{16\nu^2(r-1)}{(r-2)^2S}\cdot\frac{36r^4W^2}{S^2(r-2)^2(r-3)^2}
 \cdot\frac{S^2(r-2)^4(r-3)^2}{576\,r^4(r-1)W^2}
 =\frac{\nu^2}{S} .
\]
\end{proof}
\end{annexitem}

\begin{annexitem}{annex:p15-w02b-p15-057.contours}{thm:p15-w02b-p15-057-contours}{The efficiency contours near the normal point}
\begin{proof}
\emph{The forms in \(r\) and \(\nu\).} Write \(S=r^2+\nu^2\). From item~3
of \cref{thm:p15-w02b-p15-057-typeiv},
\(1-E_1=\bigl[(r+1)(r+2)(r+4)S-r^2(r-1)((r+4)^2+\nu^2)\bigr]/
\bigl[(r+1)(r+2)(r+4)S\bigr]\). Since
\((r+1)(r+2)(r+4)=r^3+7r^2+14r+8\) and
\(r^2(r-1)(r+4)^2=r^5+7r^4+8r^3-16r^2\), the numerator is
\(6r^3+24r^2+\nu^2\bigl(r^3+7r^2+14r+8-r^3+r^2\bigr)
=2\bigl(3r^3+12r^2+\nu^2(4r^2+7r+4)\bigr)\). From item~4 there,
\(1-E_2=\bigl[(r+1)P-(r-2)(r-3)(r+4)S\bigr]/\bigl[(r+1)P\bigr]\) with
\(P=r^3(r-2)+\nu^2(r^2+10r-12)\); since \((r-2)(r-3)(r+4)=r^3-r^2-14r+24\)
and \((r+1)r^3(r-2)=r^5-r^4-2r^3\), the numerator is
\(12r^3-24r^2+\nu^2\bigl(r^3+11r^2-2r-12-r^3+r^2+14r-24\bigr)
=12\bigl(r^3-2r^2+\nu^2(r^2+r-3)\bigr)\).

\emph{Substitution.} In the region \(y<0\), \(x^2<-y\), put
\(r=(y^2-1)/y\) and \(\nu^2=-x^2(x^2y+(1+y^2)^2)/(y^2(x^2+y))\), which is
the solution recorded in \cref{def:p15-w02b-p15-057-pearson-system}: indeed
\(r^2+\nu^2\) then equals
\(\bigl((1+y)^2+x^2\bigr)\bigl((1-y)^2-x^2\bigr)/(y(x^2+y))\), as one
checks by writing both sides over \(y^2(x^2+y)\): with
\(P=x^2y+(1+y^2)^2\), so that \(\nu^2=-x^2P/(y^2(x^2+y))\), the
numerator of \(r^2+\nu^2\) is \((y^2-1)^2(x^2+y)-x^2P
=x^2\bigl((y^2-1)^2-(1+y^2)^2\bigr)+y(y^2-1)^2-x^4y
=y\bigl((1-y^2)^2-4x^2y-x^4\bigr)\), and
\(\bigl((1+y)^2+x^2\bigr)\bigl((1-y)^2-x^2\bigr)
=(1-y^2)^2-x^2\bigl((1+y)^2-(1-y)^2\bigr)-x^4\) is the same bracket. Note \(r+1=(y^2+y-1)/y\),
\(r+2=(y^2+2y-1)/y\), \(r+4=(y^2+4y-1)/y\) and \(r-2=(y^2-2y-1)/y\).

For \(1-E_1\), multiply numerator and denominator of the form in \(r,\nu\)
by \(y^4(x^2+y)\). Since \(y^2(4r^2+7r+4)=4(y^2-1)^2+7y(y^2-1)+4y^2
=q_1\) with \(q_1=4y^4+7y^3-4y^2-7y+4\), the numerator becomes
\(2\bigl[3y(y^2-1)^3(x^2+y)+12y^2(y^2-1)^2(x^2+y)-x^2Pq_1\bigr]\), and the
denominator becomes
\(\bigl((1+y)^2+x^2\bigr)\bigl((1-y)^2-x^2\bigr)(y^2+y-1)(y^2+2y-1)
(y^2+4y-1)\). Changing the sign of both and collecting powers of \(x^2\)
(\(P\) contributes \(x^2y\) and \((1+y^2)^2\)), the numerator is
\[
 -6y^2(y^2-1)^2(y^2+4y-1)
 +2x^2\bigl[(1+y^2)^2q_1-3y(y^2-1)^3-12y^2(y^2-1)^2\bigr]+2x^4yq_1 .
\]
For \(1-E_2\), multiply by \(y^4(x^2+y)\) likewise. With
\(y^2(r^2+r-3)=q_2=y^4+y^3-5y^2-y+1\), \(y^2(r^2+10r-12)=q_3
=y^4+10y^3-14y^2-10y+1\) and \(y^4r^3(r-2)=(y^2-1)^3(y^2-2y-1)\), the
numerator becomes \(12Q_2\) with
\(Q_2=y(y^2-1)^3(x^2+y)-2y^2(y^2-1)^2(x^2+y)-x^2Pq_2\), and the
denominator becomes \(\frac{y^2+y-1}y\,R_2\) with
\(R_2=(y^2-1)^3(y^2-2y-1)(x^2+y)-x^2Pq_3\). So
\(1-E_2=12yQ_2/((y^2+y-1)R_2)\). The constant term of \(R_2\) in \(x\) is
\(y(y^2-1)^3(y^2-2y-1)\), the coefficient of \(x^2\) is
\((y^2-1)^3(y^2-2y-1)-(1+y^2)^2q_3\), which vanishes at \(y=0\), and that
of \(x^4\) is \(-yq_3\); so \(R_2=-y\widetilde R_2\) with
\[
 \widetilde R_2=-(y^2-1)^3(y^2-2y-1)
 -x^2\,\frac{(y^2-1)^3(y^2-2y-1)-(1+y^2)^2q_3}{y}+x^4q_3 ,
\]
a polynomial, and \(1-E_2=-12Q_2/((y^2+y-1)\widetilde R_2)\). Collecting
powers of \(x^2\), \(-12Q_2\) is
\[
 -12y^2(y^2-1)^2(y^2-2y-1)
 -12x^2\bigl[y(y^2-1)^3-2y^2(y^2-1)^2-(1+y^2)^2q_2\bigr]+12x^4yq_2 .
\]

The result is \(1-E_1=N_1/D_1\) with
\begin{align*}
 N_1&=6y^2-24y^3-18y^4+48y^5+18y^6-24y^7-6y^8\\
 &\quad+x^2\bigl(8-8y-16y^2-32y^3+48y^4+32y^5-16y^6+8y^7+8y^8\bigr)\\
 &\quad+x^4\bigl(8y-14y^2-8y^3+14y^4+8y^5\bigr),\\
 D_1&=\bigl(x^2-(1-y)^2\bigr)(y^2+y-1)(y^2+2y-1)(y^2+4y-1)
 \bigl((1+y)^2+x^2\bigr),
\end{align*}
and \(1-E_2=N_2/D_2\) with
\begin{align*}
 N_2&=12y^2+24y^3-36y^4-48y^5+36y^6+24y^7-12y^8\\
 &\quad+x^2\bigl(12-12y^2-48y^3-144y^4+48y^5-12y^6+12y^8\bigr)\\
 &\quad+x^4\bigl(12y-12y^2-60y^3+12y^4+12y^5\bigr),\\
 D_2&=(y^2+y-1)\bigl(-1-2y+4y^2+6y^3-6y^4-6y^5+4y^6+2y^7-y^8\\
 &\quad+x^2(-12-8y-4y^2-32y^3+4y^4-8y^5+12y^6)
 +x^4(1-10y-14y^2+10y^3+y^4)\bigr).
\end{align*}
Expanding the products in the last two displays and in
\(\widetilde R_2\) gives the polynomials \(N_1\), \(N_2\) and the second
factor of \(D_2\) shown above, with \(D_1\) the factored denominator of
\(1-E_1\) after the change of sign and \(D_2=(y^2+y-1)\widetilde R_2\).

\emph{Expansion.} \(D_1(0,0)=(-1)(-1)(-1)(-1)\cdot1=1\) and
\(D_2(0,0)=(-1)(-1)=1\). The terms of \(N_1\) of total degree two are
\(8x^2+6y^2\), and those of \(N_2\) are \(12x^2+12y^2\); all other terms of
\(N_1,N_2\) have total degree at least three. Let \(Q_1=8x^2+6y^2\),
\(Q_2=12(x^2+y^2)\). Then \(N_i-Q_iD_i\) is a polynomial whose terms all
have total degree at least three, because \(D_i-1\) has no constant term.
Choose \(\rho\in(0,1)\) with \(|D_i|\geq\frac12\) on the closed disc of
radius \(\rho\), possible by continuity. On that disc each monomial of
degree \(d\geq3\) is at most \((x^2+y^2)^{d/2}\leq(x^2+y^2)^{3/2}\) in
absolute value, so \(|N_i-Q_iD_i|\leq c_i(x^2+y^2)^{3/2}\) with \(c_i\) the
sum of the absolute coefficients, and
\(|1-E_i-Q_i|=|N_i-Q_iD_i|/|D_i|\leq2c_i(x^2+y^2)^{3/2}\). Take
\(C=2\max(c_1,c_2)\). On the Type~IV side of this disc, \(r=y-1/y=1/|y|-|y|>1/\rho-\rho\)
is large, so the efficiencies of \cref{thm:p15-w02b-p15-057-typeiv} are
defined (\(r>3\) once \(\rho<\frac14\)) and equal these rational functions.
\end{proof}
\end{annexitem}

\begin{annexitem}{annex:p15-w02b-p15-057.terminus}{thm:p15-w02b-p15-057-terminus}{Curves that end at a finite point}
\begin{proof}
\emph{Item 1.} By independence,
\(P(X_{(1)}-m>s)=(1-F(m+s))^n\) for \(s\geq0\). Fix \(t>0\) and put
\(s_n=tn^{-1/(\alpha+1)}\to0\). Then
\(u_n=F(m+s_n)=k't^{\alpha+1}n^{-1}(1+o(1))\), so \(nu_n\to k't^{\alpha+1}\)
and \(u_n\to0\). For \(0\leq u\leq\frac12\),
\(|\log(1-u)+u|\leq u^2\), so \(|n\log(1-u_n)+nu_n|\leq nu_n^2\to0\), and
\((1-u_n)^n\to\exp(-k't^{\alpha+1})\). The case \(t=0\) is trivial
(\(F(m)=0\)). If \(\alpha<1\), then for \(\varepsilon>0\),
\(P(\sqrt n(X_{(1)}-m)>\varepsilon)=P\bigl(n^{1/(\alpha+1)}(X_{(1)}-m)
>\varepsilon n^{1/(\alpha+1)-1/2}\bigr)\); the threshold tends to
infinity because \(1/(\alpha+1)>1/2\), and for any fixed \(T\) the
probability is eventually at most
\(P(n^{1/(\alpha+1)}(X_{(1)}-m)>T)\to\exp(-k'T^{\alpha+1})\), which is
arbitrarily small for large \(T\). Since \(X_{(1)}\geq m\) almost surely,
\(\sqrt n(X_{(1)}-m)\to0\) in probability.

\emph{Item 2.} On \((-a_1,a_2)\),
\(f'/f=m_1/(a_1+x)-m_2/(a_2-x)\). On \([-a_1/2,a_2/2]\) the integrand
\((f'/f)^2f\) is bounded. Near \(a_2\), with \(h=a_2-x\in(0,a_2/2)\), the
term \(m_1/(a_1+x)\) is bounded by \(|m_1|/a_1\) and \(f\) lies between
positive multiples of \(h^{m_2}\). If \(m_2>1\), then
\((f'/f)^2\leq2m_1^2/a_1^2+2m_2^2h^{-2}\) and the integrand is at most a
constant times \(h^{m_2}+h^{m_2-2}\), integrable at \(h=0\) since
\(m_2-2>-1\). If \(m_2\leq1\) (and \(m_2\neq0\)), then for \(h\) small
enough that \(|m_2|/h\geq2|m_1|/a_1\),
\((f'/f)^2\geq m_2^2/(4h^2)\), and the integrand is at least a positive
multiple of \(h^{m_2-2}\), which is not integrable at \(0\) since
\(m_2-2\leq-1\). The end \(-a_1\) is treated in the same way with
\(h=a_1+x\) and the exponent \(m_1\). The integral is finite exactly when
both ends give finite contributions.
\end{proof}
\end{annexitem}

\begin{annexitem}{annex:p15-w02b-p15-057.rectangle}{thm:p15-w02b-p15-057-rectangle}{The rectangular law: mean against the mean of the extremes}
\begin{proof}
Write \(X_i=m-\frac a2+aU_i\) with \(U_i\) independent uniform on
\((0,1)\); then \(\overline x-m=a(\overline U-\frac12)\) and
\(\hat m-m=a(M-\frac12)\) with \(M=\frac12(U_{(1)}+U_{(n)})\).

\emph{Item 1.} \(E\overline U=\frac12\) and
\(\operatorname{Var}\overline U=1/(12n)\). The pair
\((U_{(1)},U_{(n)})\) has density \(n(n-1)(v-u)^{n-2}\) on \(0<u<v<1\),
since the event that the least lies in \(du\), the greatest in \(dv\) and
the other \(n-2\) between them has probability
\(n(n-1)(v-u)^{n-2}du\,dv\). For \(0\leq t\leq\frac12\), \(M\leq t\)
forces \(U_{(1)}\leq t\) and \(U_{(n)}\leq2t-U_{(1)}\), so
\[
 P(M\leq t)=\int_0^t\!\int_u^{2t-u}n(n-1)(v-u)^{n-2}\,dv\,du
 =\int_0^tn(2t-2u)^{n-1}\,du=2^{n-1}t^n .
\]
The map \(U_i\mapsto1-U_i\) preserves the joint law and sends \(M\) to
\(1-M\), so \(M\) is symmetric about \(\frac12\), \(EM=\frac12\), and its
density is \(n(2t)^{n-1}\) on \((0,\frac12]\) and \(n(2-2t)^{n-1}\) on
\([\frac12,1)\), that is \(n(1-2|t-\frac12|)^{n-1}\). Hence
\(\hat m-m=a(M-\frac12)\) has density
\(\frac na(1-\frac{2|x|}a)^{n-1}\) on \(|x|<\frac a2\). Its variance is
\[
 2\int_0^{a/2}x^2\frac na\Bigl(1-\frac{2x}a\Bigr)^{n-1}dx
 =\frac{na^2}4\int_0^1w^2(1-w)^{n-1}\,dw
 =\frac{na^2}4\cdot\frac{2\,(n-1)!}{(n+2)!}
 =\frac{a^2}{2(n+1)(n+2)},
\]
with \(w=2x/a\) and the beta integral
\(\int_0^1w^2(1-w)^{n-1}dw=2!\,(n-1)!/(n+2)!\).

\emph{Item 2.} The ratio is
\(\frac{a^2}{2(n+1)(n+2)}\cdot\frac{12n}{a^2}=\frac{6n}{(n+1)(n+2)}\), and
\(\frac{6n}{(n+1)(n+2)}-\frac6n=\frac{6n^2-6(n+1)(n+2)}{n(n+1)(n+2)}
=-\frac{6(3n+2)}{n(n+1)(n+2)}\), which is \(O(n^{-2})\).

\emph{Item 3.} Let \(W_n=\frac{2n}a(\hat m-m)\). For \(0\leq t<n\), by
item~1 and symmetry,
\(P(W_n>t)=\int_{at/(2n)}^{a/2}\frac na(1-\frac{2x}a)^{n-1}dx
=\frac12(1-\frac tn)^n\to\frac12e^{-t}\), and \(P(W_n<-t)\) is the same.
These limits are the tail probabilities of the density \(\frac12e^{-|t|}\),
so the distribution functions converge at every point. For the density
\(\frac na e^{-2n|x|/a}\), the logarithmic derivative, which exists except
at \(0\), is \(-\frac{2n}a\operatorname{sign}x\); its mean square is
\(4n^2/a^2\). For a normal density of variance \(v\), the logarithmic
derivative is \(-x/v\), with mean square \(1/v=12n/a^2\) when
\(v=a^2/(12n)\). The ratio is \((12n/a^2)/(4n^2/a^2)=3/n\).
\end{proof}
\end{annexitem}

\begin{annexitem}{annex:p15-w02b-p15-057.form-vii}{thm:p15-w02b-p15-057-form-vii}{The form parameters of the Type IV and Type VII curves}
\begin{proof}
\emph{Normalizer.} Use the angular coordinate of the proof of
\cref{thm:p15-w02b-p15-057-typeiv}: \(\xi=\tan\vartheta\),
\(\vartheta\in(-\pi/2,\pi/2)\), so that
\((1+\xi^2)^{-(r+2)/2}e^{-\nu\tan^{-1}\xi}\,d\xi
=e^{-\nu\vartheta}\cos^r\vartheta\,d\vartheta\) and the kernel integrates
to \(I_r=\int_{-\pi/2}^{\pi/2}e^{-\nu\vartheta}\cos^r\vartheta\,d\vartheta\).
Substituting \(\vartheta=\pi/2-\theta\), \(\cos\vartheta=\sin\theta\) and
\(e^{-\nu\vartheta}=e^{-\nu\pi/2}e^{\nu\theta}\), so
\(I_r=\mathrm F(\nu,r)\). Under \(f\), \(\vartheta=\tan^{-1}\xi\) has the
density \(e^{-\nu\vartheta}\cos^r\vartheta/\mathrm F\) on
\((-\pi/2,\pi/2)\). For fixed \((\nu,r)\) with \(r>-1\) and any
\(j,k\geq0\), \(|\vartheta|^j|\log\cos\vartheta|^k\cos^{r'}\vartheta\) is
integrable for every \(r'\) in a neighbourhood of \(r\) inside
\((-1,\infty)\), and \(e^{-\nu'\vartheta}\leq e^{\pi|\nu'|/2}\); so
\(\mathrm F\) may be differentiated under the integral sign any number of
times in \((\nu,r)\), and
\(\partial_\nu\log\mathrm F=-E\vartheta\),
\(\partial_r\log\mathrm F=E\log\cos\vartheta\),
\(\partial_\nu^2\log\mathrm F=\operatorname{Var}\vartheta\),
\(\partial_\nu\partial_r\log\mathrm F=-\operatorname{Cov}(\vartheta,\log
\cos\vartheta)\), \(\partial_r^2\log\mathrm F=\operatorname{Var}\log\cos
\vartheta\).

\emph{Scores and second derivatives.} With \(\phi\) as in
\cref{thm:p15-w02b-p15-057-typeiv},
\(\log f=-\log a-\log\mathrm F(\nu,r)+\phi(\xi)\), where
\(\phi(\xi)=-\nu\tan^{-1}\xi-\frac{r+2}2\log(1+\xi^2)\), and
\(\partial_m\xi=-1/a\), \(\partial_a\xi=-\xi/a\). Hence
\(\partial_\nu\log f=-\partial_\nu\log\mathrm F-\tan^{-1}\xi\),
\(\partial_r\log f=-\partial_r\log\mathrm F-\frac12\log(1+\xi^2)\), and
\[
 \partial_m\partial_\nu\log f=\frac{1}{a(1+\xi^2)},\quad
 \partial_a\partial_\nu\log f=\frac{\xi}{a(1+\xi^2)},\quad
 \partial_m\partial_r\log f=\frac{\xi}{a(1+\xi^2)},\quad
 \partial_a\partial_r\log f=\frac{\xi^2}{a(1+\xi^2)},
\]
while the \((\nu,r)\) second derivatives are the nonrandom
\(-\nabla^2\log\mathrm F\) and the \((m,a)\) ones are those listed in the
proof of \cref{thm:p15-w02b-p15-057-typeiv}, item~2.

\emph{Regularity.} The support is \(\mathbb R\) for every parameter, and
\(\log f\) is infinitely differentiable in \((m,a,\nu,r)\) on the stated
open set. Fix a point \(\theta_0=(m_0,a_0,\nu_0,r_0)\) and a closed ball
\(N\) about it on which \(a\geq a_0/2\), \(a\leq2a_0\), \(|m-m_0|\leq1\),
\(|\nu|\leq|\nu_0|+1\) and \(r\geq r_-\), where \(-1<r_-<r_0\). On \(N\),
\(|\phi'|\), \(|\xi\phi'|\), \(|\phi''|\), \(|\xi\phi''|\) and
\(|\xi^2\phi''|\) are bounded by a constant depending on the bounds for
\(|\nu|\) and \(r\) (the formulas for \(\phi',\phi''\) in the proof of
\cref{thm:p15-w02b-p15-057-typeiv}), the four mixed derivatives above are
at most \(2/a_0\) in absolute value, and \(\nabla^2\log\mathrm F\) is
continuous, hence bounded; this gives the constant envelope \(M\). Next,
\(|x|\leq|m|+a|\xi|\) gives \(1+x^2\leq c(1+\xi^2)\) and
\(|\xi|\leq(|x|+|m|)/a\) gives \(1+\xi^2\leq c(1+x^2)\) on \(N\), for a
constant \(c\). So \(f\leq c'(1+x^2)^{-(r_-+2)/2}\) and
\(|s|\leq c''(1+\log(1+x^2))\) on \(N\), and \(|\nabla f|=f|s|\) and
\(\|\nabla^2f\|\leq f(\|\nabla^2\log f\|+|s|^2)\) are bounded by
\(G(x)=c'''(1+\log(1+x^2))^2(1+x^2)^{-(r_-+2)/2}\), which is integrable
because \(r_-+2>1\). By the mean value theorem and dominated convergence,
\(\int f\,dx=1\) may be differentiated twice under the integral sign, so
\(\int\nabla f\,dx=0\) and \(\int\nabla^2f\,dx=0\). The score has
\(E|s|^2<\infty\) by the same bound, so \(J\) is finite.

\emph{Positive definiteness.} If \(t^{\top}s=0\) almost surely for some
\(t\in\mathbb R^4\), then, the scores being continuous in \(\xi\), the
identity holds for every \(\xi\in\mathbb R\):
\[
 c_0+t_1\frac{\nu+(r+2)\xi}{a(1+\xi^2)}
 +t_2\frac{\xi(\nu+(r+2)\xi)}{a(1+\xi^2)}
 -t_3\tan^{-1}\xi-\frac{t_4}2\log(1+\xi^2)=0,
\]
\(c_0\) collecting the constants. As \(\xi\to\pm\infty\) the rational terms
are bounded and \(\log(1+\xi^2)\to\infty\), so \(t_4=0\). The rational
terms have the same limit at \(+\infty\) and \(-\infty\), whereas
\(\tan^{-1}\xi\) tends to \(\pm\pi/2\), so \(t_3=0\). Multiplying by
\(a(1+\xi^2)\), the polynomial \(ac_0(1+\xi^2)+(t_1+t_2\xi)(\nu+(r+2)\xi)\)
vanishes: its coefficients give \(ac_0+t_2(r+2)=0\),
\(t_1(r+2)+t_2\nu=0\), \(ac_0+t_1\nu=0\). Eliminating, \(t_1=-t_2\nu/(r+2)\)
and \(ac_0=t_2\nu^2/(r+2)=-t_2(r+2)\), so \(t_2D_2=0\), \(t_2=0\), and then
\(t_1=0\). So \(J=E\,ss^{\top}\) is positive definite. By
\cref{thm:p15-w02b-p15-057-multi-regular}, \(J=-E\nabla^2\log f\).

\emph{Entries.} The \((m,a)\) block is item~1 of
\cref{thm:p15-w02b-p15-057-typeiv}. In angular coordinates
\(1/(1+\xi^2)=\cos^2\vartheta\), \(\xi/(1+\xi^2)=\sin\vartheta\cos
\vartheta\) and \(\xi^2/(1+\xi^2)=1-\cos^2\vartheta\); by the values
\(E\cos^2\vartheta=(r+1)(r+2)/D_2\) and
\(E\sin\vartheta\cos\vartheta=-\nu(r+1)/D_2\) in that proof,
\(J_{m\nu}=-(r+1)(r+2)/(aD_2)\), \(J_{a\nu}=J_{mr}=\nu(r+1)/(aD_2)\) and
\(J_{ar}=-(1-(r+1)(r+2)/D_2)/a=-(r+2+\nu^2)/(aD_2)\), since
\(D_2-(r+1)(r+2)=r+2+\nu^2\). The \((\nu,r)\) block is
\(\nabla^2\log\mathrm F\). This is item~1.

\emph{Item 2.} At \(\nu=0\) the law of \(\vartheta\) is symmetric about
\(0\). Then \(J_{ma}\), \(J_{a\nu}\), \(J_{mr}\) vanish (each carries the
factor \(\nu\)), and \(\operatorname{Cov}(\vartheta,\log\cos\vartheta)=0\)
because the product is odd; so \(J\) is block diagonal. For the
\((\nu,\nu)\) entry, let \(k>1\), \(I_k=\int\cos^k\vartheta\,d\vartheta\)
and \(L_k=\int\vartheta^2\cos^k\vartheta\,d\vartheta\) over
\((-\pi/2,\pi/2)\). Differentiating,
\(\frac d{d\vartheta}\bigl(\vartheta^2\cos^{k-1}\vartheta\sin\vartheta\bigr)
=2\vartheta\cos^{k-1}\vartheta\sin\vartheta+\vartheta^2\bigl(k\cos^k
\vartheta-(k-1)\cos^{k-2}\vartheta\bigr)\); integrating, the boundary term
being zero, and using \(\int\vartheta\cos^{k-1}\vartheta\sin\vartheta\,
d\vartheta=I_k/k\) (by parts against \(d\cos^k\vartheta\)),
\(kL_k-(k-1)L_{k-2}=-2I_k/k\). With \(I_k=\frac{k-1}kI_{k-2}\) (the
recursion of that proof at \(\nu=0\)), \(V_k=L_k/I_k\) satisfies
\(V_k=V_{k-2}-2/k^2\). Since \(|\vartheta|\leq\frac\pi2|\sin\vartheta|\)
on \((-\pi/2,\pi/2)\) and \(E\sin^2\vartheta=1-\frac{k+1}{k+2}\) under the
density proportional to \(\cos^k\vartheta\), \(V_k\leq\pi^2/(4(k+2))\to0\)
as \(k\to\infty\). Iterating, \(V_r=\sum_{j\geq1}2/(r+2j)^2
=\frac12\sum_{j\geq0}(\frac r2+1+j)^{-2}=\frac12\psi'(\frac r2+1)\) by
\cref{thm:p15-w02b-p15-057-trigamma}; the iteration starts from any
\(r>-1\) because \(V_r=V_{r+2}+2/(r+2)^2\) uses only \(r+2>1\). For the
\((r,r)\) entry, at \(\nu=0\) the substitution \(t=\sin^2\vartheta\) gives
\(\mathrm F(0,r)=\int_0^1t^{-1/2}(1-t)^{(r-1)/2}\,dt
=\Gamma(\frac12)\Gamma(\frac{r+1}2)/\Gamma(\frac r2+1)\). Since
\(\nabla^2\log\mathrm F\) is continuous and \(\partial_r^2\log\mathrm F(0,r)\)
is the second derivative of \(r\mapsto\log\mathrm F(0,r)\), it equals
\(\frac14\psi'(\frac{r+1}2)-\frac14\psi'(\frac r2+1)\).

For a positive definite \(2\times2\) block with diagonal \(A,A'\) and
off-diagonal \(B\), the diagonal entries of the inverse are \(A'/\Delta\)
and \(A/\Delta\), \(\Delta=AA'-B^2>0\), and the block inverse is the
corresponding block of \(J^{-1}\) when \(J\) is block diagonal. For
\((m,\nu)\): \(a^2\Delta=\frac{(r+1)(r+2)}{r+4}\cdot\frac12\psi'(\frac r2+1)
-\frac{(r+1)^2}{(r+2)^2}=\frac{(r+1)G_\nu(r)}{2(r+2)^2(r+4)}\), so
\(G_\nu>0\) and \((J^{-1})_{\nu\nu}=\frac{(r+1)(r+2)}{r+4}\cdot
\frac{2(r+2)^2(r+4)}{(r+1)G_\nu}=\frac{2(r+2)^3}{G_\nu}\). For \((a,r)\):
\(a^2\Delta=\frac{2(r+1)}{r+4}\cdot\frac14\bigl(\psi'(\frac{r+1}2)
-\psi'(\frac r2+1)\bigr)-\frac1{(r+2)^2}=\frac{G_r(r)}{2(r+1)(r+2)^2(r+4)}\),
so \(G_r>0\) and \((J^{-1})_{rr}=\frac{2(r+1)}{r+4}\cdot
\frac{2(r+1)(r+2)^2(r+4)}{G_r}=\frac{4(r+1)^2(r+2)^2}{G_r}\). The limit
statement is \cref{thm:p15-w02b-p15-057-multi-regular} with
\(\tau(m,a,\nu,r)=\nu\) and \(=r\).

\emph{Item 3: the moment map.} For \(r>3\) the moments of \(\xi\) up to
order four exist (item~1 of \cref{thm:p15-w02b-p15-057-pearson-moments}),
and item~2 there gives \(\mu_1'=-\nu/r\),
\(\mu_2'=(1-\nu\mu_1')/(r-1)=(r+\nu^2)/(r(r-1))\) and
\(\mu_3'=(2\mu_1'-\nu\mu_2')/(r-2)=-\nu(\nu^2+3r-2)/(r(r-1)(r-2))\). The
central moments are \(c_2=\mu_2'-\mu_1'^2=(r^2+\nu^2)/(r^2(r-1))\) and
\(c_3=\mu_3'-3\mu_2'\mu_1'+2\mu_1'^3=-4\nu(r^2+\nu^2)/(r^3(r-1)(r-2))\),
the last because
\(-(\nu^2+3r-2)r^2+3r(r+\nu^2)(r-2)-2\nu^2(r-1)(r-2)=-4(r^2+\nu^2)\)
after division by \(\nu\) and collection over
\(r^3(r-1)(r-2)\). Hence
\(\gamma_1=c_3/c_2^{3/2}=-4\nu\sqrt{r-1}/((r-2)\sqrt{r^2+\nu^2})\); location
and scale do not change \(\gamma_1,\beta_2\). Both are continuously
differentiable on \(\{r>3\}\), \(\gamma_1\) is odd and \(\beta_2\) even in
\(\nu\) (the reflection \(\xi\mapsto-\xi\) maps \(\nu\) to \(-\nu\)), so at
\((0,r_0)\): \(\partial_\nu\gamma_1=-4\sqrt{r_0-1}/(r_0(r_0-2))\),
\(\partial_r\gamma_1=0\), \(\partial_\nu\beta_2=0\) and
\(\partial_r\beta_2=-6/(r_0-3)^2\), from \(\beta_2(0,r)=3(r-1)/(r-3)\). The
derivative is diagonal and nonsingular, so the inverse function theorem
gives \(\Psi\) as in \cref{def:p15-w02b-p15-057-form-moments}, with
\(\partial_1\Psi_1=1/\partial_\nu\gamma_1\), \(\partial_2\Psi_1=0\),
\(\partial_1\Psi_2=0\), \(\partial_2\Psi_2=1/\partial_r\beta_2\) at
\(\Gamma(0,r_0)=(0,\beta_2^0)\). Shrinking \(W\), the reflection symmetry
and local uniqueness of the inverse give \(\Psi_1(-g,b)=-\Psi_1(g,b)\) and
\(\Psi_2(-g,b)=\Psi_2(g,b)\); in particular \(\Psi_1(0,b)=0\), and
\(\Psi_2(0,b)=3(b-1)/(b-3)\), the inverse of \(r\mapsto3(r-1)/(r-3)\).

\emph{Item 3: sample moments.} At \(\nu=0\) and \(a=1\) (the ratios
\(g_1,b_2\) do not depend on \(m,a\)), the substitution
\(t=\sin^2\vartheta\) gives, for \(2j<r+1\),
\(E\xi^{2j}=B(j+\frac12,\frac{r+1}2-j)/B(\frac12,\frac{r+1}2)
=(2j-1)(2j-3)\cdots1\big/\bigl((r-1)(r-3)\cdots(r+1-2j)\bigr)\), and the
odd moments vanish. So \(\mu_2=1/(r-1)\), \(\mu_4=3/((r-1)(r-3))\),
\(\mu_6=15/((r-1)(r-3)(r-5))\), \(\mu_8=105/((r-1)(r-3)(r-5)(r-7))\), each
when finite. Write \(Z_i=x_i-m\) and \(\overline Z\) for their mean. As in
the proof of \cref{thm:p15-w02b-p15-057-moment-shape}, with
\(\mu_3=0\), \(\sqrt n\,m_3=n^{-1/2}\sum_iV_i+o_P(1)\),
\(V_i=Z_i^3-3\mu_2Z_i\), when \(\mu_6<\infty\); \(\sqrt n(m_2-\mu_2)\) is
tight when \(\mu_4<\infty\); and expanding
\(m_4=n^{-1}\sum Z_i^4-4\overline Z\,n^{-1}\sum Z_i^3
+6\overline Z^2n^{-1}\sum Z_i^2-3\overline Z^4\), each term after the first
is \(\sqrt n\,\overline Z\) (tight) times a sequence tending to zero
almost surely (\(n^{-1}\sum Z_i^3\to\mu_3=0\)), so
\(\sqrt n(m_4-\mu_4)=n^{-1/2}\sum_i(Z_i^4-\mu_4)+o_P(1)\) when
\(\mu_8<\infty\). Then
\(b_2-\beta_2=(m_4-\mu_4)/m_2^2-\mu_4(m_2^2-\mu_2^2)/(m_2^2\mu_2^2)\), and
since \(m_2\to\mu_2\) almost surely,
\(\sqrt n(b_2-\beta_2)=n^{-1/2}\sum_iU_i+o_P(1)\) with
\(U_i=(Z_i^4-\mu_4)/\mu_2^2-2\mu_4(Z_i^2-\mu_2)/\mu_2^3\). By
\cref{thm:p15-clt-iid} and \cref{thm:p15-clt-perturbation},
\(\sqrt n\,g_1\Rightarrow N(0,EV^2/\mu_2^3)\) when \(r>5\) and
\(\sqrt n(b_2-\beta_2)\Rightarrow N(0,EU^2)\) when \(r>7\). Expanding the
squares, \(EV^2/\mu_2^3=\mu_6/\mu_2^3-6\mu_4/\mu_2^2+9\) and
\(EU^2=\mu_8/\mu_2^4-\beta_2^2-4\beta_2\mu_6/\mu_2^3+4\beta_2^3\). With the
moments above, \(\mu_6/\mu_2^3=15(r-1)^2/((r-3)(r-5))\) and
\(15(r-1)^2-18(r-1)(r-5)+9(r-3)(r-5)=6r^2+6r+60\), so
\(EV^2/\mu_2^3=6(r^2+r+10)/((r-3)(r-5))\). Over the denominator
\((r-3)^3(r-5)(r-7)\), \(EU^2\) has numerator
\(105(r-1)^3(r-3)^2-9(r-1)^2(r-3)(r-5)(r-7)-180(r-1)^3(r-3)(r-7)
+108(r-1)^3(r-5)(r-7)\), which equals \(24r(r-1)^2(r^2-r+18)\) (both sides
are polynomials of degree five; dividing out \((r-1)^2\), the cubic
\(105(r-1)(r-3)^2-9(r-3)(r-5)(r-7)-180(r-1)(r-3)(r-7)+108(r-1)(r-5)(r-7)\)
has coefficients \(24,-24,432,0\), that is \(24r(r^2-r+18)\)).

\emph{Item 3: delta step.} Let \(r_0>5\). Since \(g_1\to0\) and
\(b_2\to\beta_2^0\) almost surely (\(\mu_4<\infty\)), the event
\(\{m_2>0,(g_1,b_2)\in W\}\) has probability tending to one, and on it
\(\nu_\mu=\Psi_1(g_1,b_2)=g_1\int_0^1\partial_1\Psi_1(tg_1,b_2)\,dt\), the
integral tending in probability to \(\partial_1\Psi_1(0,\beta_2^0)
=-r_0(r_0-2)/(4\sqrt{r_0-1})\) by continuity. So \(\sqrt n\,\nu_\mu\)
converges in law to the normal law with variance
\(\frac{r_0^2(r_0-2)^2}{16(r_0-1)}\cdot\frac{6(r_0^2+r_0+10)}{(r_0-3)(r_0-5)}\),
which is the stated value. Let \(r_0>7\). Then
\(r_\mu-r_0=\bigl(\Psi_2(g_1,b_2)-\Psi_2(0,b_2)\bigr)
+\bigl(\Psi_2(0,b_2)-\Psi_2(0,\beta_2^0)\bigr)\). The first bracket is
\(g_1\int_0^1\partial_1\Psi_2(tg_1,b_2)\,dt\), with \(\sqrt n\,g_1\) tight
and the integral tending to \(\partial_1\Psi_2(0,\beta_2^0)=0\), so it is
\(o_P(n^{-1/2})\). The second is the scalar delta method for the
differentiable \(b\mapsto3(b-1)/(b-3)\), whose derivative at \(\beta_2^0\)
is \(-(r_0-3)^2/6\). So \(\sqrt n(r_\mu-r_0)\) has the normal limit with
variance \(\frac{(r_0-3)^4}{36}\cdot\frac{24r_0(r_0-1)^2(r_0^2-r_0+18)}
{(r_0-3)^3(r_0-5)(r_0-7)}\), the stated value.

\emph{Item 4.} The efficiencies are the quotients of the item~2 values by
the item~3 values; the factors simplify to the displayed forms. By
\cref{thm:p15-w02b-p15-057-trigamma} at \(x=s/2\), \(s=r+2\),
\(\psi'(\frac s2)=\frac2s+\frac2{s^2}+\frac4{3s^3}-\frac{16}{15s^5}
+\frac{128\vartheta}{42s^7}\) with \(\vartheta\in(0,1)\); multiplying by
\(s^3\) and subtracting \(2(s-1)(s+2)=2s^2+2s-4\) gives
\(G_\nu=\frac{16}3-\frac{16}{15s^2}+\frac{64\vartheta}{21s^4}\). For \(G_r\),
apply the same bounds at \(x_1=\frac{r+1}2\) and \(x_2=\frac{r+2}2\): with
\(T(x)=\frac1x+\frac1{2x^2}+\frac1{6x^3}-\frac1{30x^5}\),
\(G_r=(r+1)^2(r+2)^2\bigl(T(x_1)-T(x_2)\bigr)-2(r+1)(r+4)+\rho\), where
\(|\rho|\leq(r+1)^2(r+2)^2\cdot\frac{128}{42}(r+1)^{-7}=O(r^{-3})\). The
first two terms are a rational function of \(r\); its expansion at
infinity is \(6-4r^{-2}+12r^{-3}+O(r^{-4})\), and
\(6-4(r+2)^{-2}=6-4r^{-2}+16r^{-3}+O(r^{-4})\), so
\(G_r=6-4(r+2)^{-2}+O(r^{-3})\). Substituting \(G_\nu=\frac{16}3
(1+O(r^{-2}))\) refined as above, and \(G_r\), into the formulas for
\(E_\nu,E_r\) and expanding the rational factors in powers of \(1/r\)
gives \(E_\nu=1-\frac{144}5r^{-2}+O(r^{-3})\) and
\(E_r=1-\frac{160}3r^{-2}+O(r^{-3})\): for \(E_\nu\),
\(\frac{16}{3G_\nu}=1+\frac1{5r^2}+O(r^{-3})\) and
\(\frac{(r+2)^3(r-1)(r-3)(r-5)}{r^2(r-2)^2(r^2+r+10)}
=1-\frac{29}{r^2}+O(r^{-3})\) (the \(1/r\) terms cancel: \(6-9+4-1=0\) from
\(+6\), \(-9\), \(+4\), \(-1\) for the factors \((1+2/r)^3\),
\((1-1/r)(1-3/r)(1-5/r)\), \((1-2/r)^{-2}\), \((1+1/r+10/r^2)^{-1}\)), and
\(\frac15-29=-\frac{144}5\); for \(E_r\), \(\frac6{G_r}=1+\frac2{3r^2}
+O(r^{-3})\) and \(\frac{(r+1)^2(r+2)^2(r-5)(r-7)}{r(r-1)^2(r-3)
(r^2-r+18)}=1-\frac{54}{r^2}+O(r^{-3})\), and \(\frac23-54
=-\frac{160}3\).
\end{proof}
\end{annexitem}

\begin{annexitem}{annex:p15-w02b-p15-057.form-ii}{thm:p15-w02b-p15-057-form-ii}{The form parameters of the Type II curves}
\begin{proof}
\emph{Normalizer.} With \(\xi=2u-1\),
\(\int_{-1}^1(1+\xi)^{\alpha}(1-\xi)^{\beta}d\xi
=2^{\alpha+\beta+1}B(\alpha+1,\beta+1)\) for \(\alpha,\beta>-1\); with
\(\alpha=(r-2-\nu)/2\), \(\beta=(r-2+\nu)/2\) this is \(Z(\nu,r)\), and
\(\alpha,\beta>-1\) is \(|\nu|<r\). In \((\nu,r)\) the family is
exponential with statistics \(-\tanh^{-1}\xi\) and
\(\frac12\log(1-\xi^2)\), since
\((1+\xi)^{\alpha}(1-\xi)^{\beta}=(1-\xi^2)^{(r-2)/2}e^{-\nu\tanh^{-1}\xi}\).
So \(\log Z\) is infinitely differentiable on \(\{|\nu|<r\}\) and the
\((\nu,r)\) block of the information is \(\nabla^2\log Z\). From
\(\log Z=(r-1)\log2+\log\Gamma(\frac{r-\nu}2)+\log\Gamma(\frac{r+\nu}2)
-\log\Gamma(r)\), at \(\nu=0\): \(\partial_\nu^2\log Z=\frac12\psi'(\frac r2)\),
\(\partial_\nu\partial_r\log Z=0\) and
\(\partial_r^2\log Z=\frac12\psi'(\frac r2)-\psi'(r)\). The duplication
formula \(\Gamma(r)=2^{r-1}\pi^{-1/2}\Gamma(\frac r2)\Gamma(\frac{r+1}2)\),
differentiated twice in logarithmic form, gives
\(\psi'(r)=\frac14\psi'(\frac r2)+\frac14\psi'(\frac{r+1}2)\), so
\(\partial_r^2\log Z=\frac14\bigl(\psi'(\frac r2)-\psi'(\frac{r+1}2)\bigr)\).

\emph{Item 1.} At \(\nu=0\), \(\xi^2\) has the beta law with parameters
\(\frac12,\frac r2\), so for real \(c\) with \(\frac r2+c>0\),
\(E(1-\xi^2)^{c}=B(\frac12,\frac r2+c)/B(\frac12,\frac r2)\); in particular
\(E(1-\xi^2)^{-1}=\frac{r-1}{r-2}\) for \(r>2\) and
\(E(1-\xi^2)^{-2}=\frac{(r-1)(r-3)}{(r-2)(r-4)}\) for \(r>4\). Let \(r>4\).
For every sample all \(|\xi_i|<1\) almost surely, and \(\log f\) is then
infinitely differentiable in \((m,a,\nu,r)\) near the true value, with
gradient the sum of the scores; this is the differentiability part of
(S). The scores at \(\nu=0\) are
\(s_m=(r-2)\xi/(a(1-\xi^2))\),
\(s_a=-\bigl(1-(r-2)\xi^2/(1-\xi^2)\bigr)/a\),
\(s_\nu=-\tanh^{-1}\xi\) and
\(s_r=\frac12\log(1-\xi^2)-\partial_r\log Z\). They have mean zero:
\(s_m\) and \(s_\nu\) are odd; \(E\frac{\xi^2}{1-\xi^2}=\frac1{r-2}\) gives
\(Es_a=0\); and \(E\frac12\log(1-\xi^2)=\partial_r\log Z\) by differentiating
\(\int(1-\xi^2)^{(r-2)/2}d\xi=Z(0,r)\) under the integral sign. Their
second moments are finite: \(E\xi^2(1-\xi^2)^{-2}
=E(1-\xi^2)^{-2}-E(1-\xi^2)^{-1}=\frac{r-1}{(r-2)(r-4)}\), and
\(\tanh^{-1}\xi\), \(\log(1-\xi^2)\) have all moments since
\((1-\xi^2)^{(r-2)/2}\) is bounded near \(\pm1\) times a logarithm power.
Hence \(J=Ess^{\top}\) has
\(J_{mm}=(r-2)^2\frac{r-1}{a^2(r-2)(r-4)}=\frac{(r-1)(r-2)}{a^2(r-4)}\)
and, since \(E\xi^4(1-\xi^2)^{-2}=E(1-\xi^2)^{-2}-2E(1-\xi^2)^{-1}+1
=\frac3{(r-2)(r-4)}\),
\(J_{aa}=\bigl((r-2)^2E\xi^4(1-\xi^2)^{-2}-1\bigr)/a^2=\frac{2(r-1)}{a^2(r-4)}\).
The cross entries follow by integration by parts in \(\xi\), the boundary
terms vanishing because \((1-\xi^2)^{(r-2)/2}\) vanishes to order greater
than \(1\) at \(\pm1\): with \(g\) the density of \(\xi\) and
\(g'/g=-(r-2)\xi/(1-\xi^2)\),
\(E[s_m s_\nu]=-\frac1aE\bigl[-\frac{g'}g\tanh^{-1}\xi\bigr]
=-\frac1a\int g\,\frac{d}{d\xi}\tanh^{-1}\xi\,d\xi
=-\frac1aE(1-\xi^2)^{-1}=-\frac{r-1}{a(r-2)}\), and
\(E[s_as_r]=-\frac1aE\bigl[(1+\xi\frac{g'}g)\frac12\log(1-\xi^2)\bigr]
=\frac1{2a}\int g\,\xi\frac{d}{d\xi}\log(1-\xi^2)\,d\xi
=-\frac1aE\frac{\xi^2}{1-\xi^2}=-\frac1{a(r-2)}\), using
\(\int(\xi g)'h=-\int\xi g h'\) and \(\int(\xi g)'=0\). The products
\(s_ms_a\), \(s_ms_r\), \(s_as_\nu\), \(s_\nu s_r\) are odd, so the other
entries vanish. Positive definiteness: as \(\xi\uparrow1\),
\((1-\xi)s_m\to(r-2)/(2a)\neq0\) and \((1-\xi)s_a\to(r-2)/(2a)\neq0\),
while \((1-\xi)s_\nu\to0\) and \((1-\xi)s_r\to0\). If
\(t_1s_m+t_2s_\nu=0\) almost surely, continuity makes it hold for all
\(|\xi|<1\); multiplying by \(1-\xi\) and letting \(\xi\uparrow1\) gives
\(t_1=0\), and then \(t_2=0\) because \(s_\nu\) is not constant. The same
argument with \(s_a,s_r\) applies to the \((a,r)\) block. So
the two blocks, and \(J\), are positive definite, and (S) holds. The
inverse values are computed as in the proof of
\cref{thm:p15-w02b-p15-057-form-vii}: for \((m,\nu)\),
\(a^2\Delta=\frac{(r-1)(r-2)}{r-4}\cdot\frac12\psi'(\frac r2)
-\frac{(r-1)^2}{(r-2)^2}=\frac{(r-1)H_\nu(r)}{2(r-2)^2(r-4)}\), giving
\(2(r-2)^3/H_\nu\); for \((a,r)\),
\(a^2\Delta=\frac{(r-1)}{2(r-4)}\bigl(\psi'(\frac r2)-\psi'(\frac{r+1}2)\bigr)
-\frac1{(r-2)^2}=\frac{H_r(r)}{2(r-1)(r-2)^2(r-4)}\), giving
\(4(r-1)^2(r-2)^2/H_r\). Positivity of \(\Delta\) gives \(H_\nu,H_r>0\).
Finally \(\psi'(x)=\psi'(x-1)-(x-1)^{-2}\) (termwise from the series of
\cref{thm:p15-w02b-p15-057-trigamma}) gives
\(\psi'(\frac r2)=\psi'(\frac{r-2}2)-\frac4{(r-2)^2}\) and
\(\psi'(\frac{r+1}2)=\psi'(\frac{r-1}2)-\frac4{(r-1)^2}\); hence
\(H_\nu(r)=(r-2)^3\psi'(\frac{r-2}2)-4(r-2)-2(r-1)(r-4)
=(r-2)^3\psi'(\frac{r-2}2)-2(r-3)r=G_\nu(r-4)\), and
\(H_r(r)=(r-1)^2(r-2)^2\bigl(\psi'(\frac{r-2}2)-\psi'(\frac{r-1}2)\bigr)
-4(r-1)^2+4(r-2)^2-2(r-1)(r-4)\), whose last three terms sum to
\(-2(r-2)(r+1)\), so \(H_r(r)=G_r(r-3)\).

\emph{Item 2.} Integrating
\(\frac d{d\xi}\bigl(\xi^k(1-\xi^2)g\bigr)\) over \((-1,1)\), with
\((1-\xi^2)g'/g=-\nu-(r-2)\xi\) for the density of the family and zero
boundary terms, gives \((r+k)\mu_{k+1}'=k\mu_{k-1}'-\nu\mu_k'\). So
\(\mu_1'=-\nu/r\), \(\mu_2'=(r+\nu^2)/(r(r+1))\),
\(\mu_3'=(2\mu_1'-\nu\mu_2')/(r+2)\), and
\(c_2=(r^2-\nu^2)/(r^2(r+1))\), \(c_3=4\nu(r^2-\nu^2)/(r^3(r+1)(r+2))\),
whence the stated \(\gamma_1\); all moments are rational in \((\nu,r)\)
with positive denominators, so \(\gamma_1\) and \(\beta_2\) are
continuously differentiable on \(\{|\nu|<r\}\), odd and even in \(\nu\).
At \(\nu=0\), \(E\xi^{2j}=\frac{(2j-1)(2j-3)\cdots1}{(r+1)(r+3)\cdots(r+2j-1)}\)
from the beta law of \(\xi^2\). These are the Type~VII values of the
proof of \cref{thm:p15-w02b-p15-057-form-vii} with \(r\) replaced by
\(-r\), up to the factor \((-1)^j\), which cancels in every ratio
\(\mu_{2j}/\mu_2^j\); so \(\beta_2(0,r)=3(r+1)/(r+3)\),
\(EV^2/\mu_2^3=6(r^2-r+10)/((r+3)(r+5))\) and
\(EU^2=24r(r+1)^2(r^2+r+18)/((r+3)^3(r+5)(r+7))\), by the same
substitution in the Type~VII formulas. The derivative of \(\Gamma\) at
\((0,r_0)\) is diagonal with \(\partial_\nu\gamma_1
=4\sqrt{r_0+1}/(r_0(r_0+2))\) and \(\partial_r\beta_2=6/(r_0+3)^2\), both
nonzero, and \(\Psi_2(0,b)=(3-3b)/(b-3)\), the inverse of
\(r\mapsto3(r+1)/(r+3)\). The linearizations and the delta step are
those of the Type~VII proof, all moments being finite; the limit
variances are \(\frac{r_0^2(r_0+2)^2}{16(r_0+1)}\cdot
\frac{6(r_0^2-r_0+10)}{(r_0+3)(r_0+5)}\) and
\(\frac{(r_0+3)^4}{36}\cdot\frac{24r_0(r_0+1)^2(r_0^2+r_0+18)}
{(r_0+3)^3(r_0+5)(r_0+7)}\), the stated values.

\emph{Item 3.} The efficiency is defined under (S), which holds for
\(r>4\) by item~1. The expansions follow as in the Type~VII proof from
\(H_\nu(r)=G_\nu(r-4)\), \(H_r(r)=G_r(r-3)\) and the expansions of
\(G_\nu,G_r\) there, together with the expansion in \(1/r\) of the
rational factors; the \(1/r^2\) coefficients are again \(-\frac{144}5\)
and \(-\frac{160}3\).
\end{proof}
\end{annexitem}

\begin{annexitem}{annex:p15-w02b-p15-057.quartic}{thm:p15-w02b-p15-057-quartic}{The quartic-exponent family}
\begin{proof}
\emph{Item 1.} If \(\theta_4<0\), then
\(\theta\cdot T(x)\leq\theta_4x^4/2\) for \(|x|\) large, so
\(e^{\theta\cdot T}\) is integrable. If \(\theta_4>0\), the exponent tends
to \(+\infty\). If \(\theta_4=0\) and \(\theta_3\neq0\), the exponent tends
to \(+\infty\) as \(x\to+\infty\) or as \(x\to-\infty\). If
\(\theta_4=\theta_3=0\), integrability of \(e^{\theta_1x+\theta_2x^2}\)
holds if and only if \(\theta_2<0\), because for \(\theta_2\geq0\) the
integrand is at least \(1\) on a half-line.

\emph{Item 2.} Fix \(\theta\in\Theta^\circ\) and \(\varepsilon>0\) with
\(\varepsilon<|\theta_4|/8\). For \(\theta'\) with
\(|\theta'_j-\theta_j|\leq\varepsilon\) for all \(j\),
\(\theta'\cdot T(x)\leq\theta\cdot T(x)+\varepsilon(|x|+x^2+|x|^3+x^4)\),
and \(|x|^je^{\theta'\cdot T(x)}\) for \(j\leq4\) is bounded by
\(G(x)=(1+x^4)e^{\theta\cdot T(x)+\varepsilon(|x|+x^2+|x|^3+x^4)}\), which
is integrable since the exponent is at most \(\theta_4x^4/2\) for large
\(|x|\). Dominated convergence applied to the difference quotients, by the
mean value theorem, gives \(\partial_jZ(\theta)=\int x^je^{\theta\cdot
T}dx\), so \(\nabla\log Z(\theta)=E_\theta T(X)\). For
\(\theta\neq\theta'\) in \(\Theta^\circ\) and \(0<t<1\), Hölder's
inequality with exponents \(1/t\) and \(1/(1-t)\) gives
\(Z(t\theta+(1-t)\theta')\leq Z(\theta)^tZ(\theta')^{1-t}\), with equality
only if \(e^{\theta\cdot T}\) and \(e^{\theta'\cdot T}\) are proportional
almost everywhere, that is, \((\theta-\theta')\cdot T(x)\) is constant
almost everywhere; a nonzero polynomial of degree at most four without
constant term takes each value at most four times, so this is impossible.
Thus \(\log Z\) is strictly convex on the convex set \(\Theta^\circ\). The
log likelihood is \(\ell_n(\theta)=\theta\cdot\sum_iT(x_i)-n\log
Z(\theta)\), strictly concave and differentiable, with
\(\nabla\ell_n(\theta)=\sum_iT(x_i)-nE_\theta T(X)\). A differentiable
concave function on an open convex set has a global maximum at \(\theta\)
if and only if its gradient vanishes there, and strict concavity makes
such a point unique.

\emph{Item 3.} The integrability statement is item~1 with
\(\theta_4=k_2/\sigma^4\), \(\theta_3=k_1/\sigma^3\),
\(\theta_2=-1/(2\sigma^2)<0\). Next,
\(\phi'(x)=-x/\sigma^2+3k_1x^2/\sigma^3+4k_2x^3/\sigma^4
=-(x/\sigma^2)(1-u)\), and \(1-u=\frac1{1+u}-\frac{u^2}{1+u}\) for
\(u\neq-1\), which gives the displayed identity. If \(|x|\leq L\sigma\),
then \(|u|\leq3|k_1|L+4|k_2|L^2\leq(3L+4L^2)(|k_1|+|k_2|)\leq\frac12\), so
\(|1+u|\geq\frac12\) and
\(\bigl|xu^2/(\sigma^2(1+u))\bigr|\leq2L\sigma u^2/\sigma^2
\leq2L(3L+4L^2)^2(|k_1|+|k_2|)^2/\sigma\).
\end{proof}
\end{annexitem}

\begin{annexitem}{annex:p15-w02b-p15-057.endpoint}{thm:p15-w02b-p15-057-endpoint}{The least observation of a Type III sample}
\begin{proof}
\emph{Item 1.} Let \(G(s)=\Gamma(p+1)^{-1}\int_0^sx^pe^{-x}\,dx\) be the
distribution function of \(\xi=(X-m)/a\). Since \(e^{-s}\leq e^{-x}\leq1\)
on \([0,s]\) and \(\Gamma(p+2)=(p+1)\Gamma(p+1)\),
\[
 \frac{e^{-s}s^{p+1}}{\Gamma(p+2)}\leq G(s)\leq\frac{s^{p+1}}{\Gamma(p+2)} .
\]
The case \(t=0\) is trivial. For \(t>0\) put
\(s_n=tn^{-1/(p+1)}\to0\); then \(nG(s_n)\) lies between
\(e^{-s_n}\lambda\) and \(\lambda\), where
\(\lambda=t^{p+1}/\Gamma(p+2)\), so \(nG(s_n)\to\lambda\). By independence
and continuity of \(G\),
\(P(X_{(1)}-m>as_n)=(1-G(s_n))^n\). For \(0\leq y\leq\frac12\),
\(-y-y^2\leq\log(1-y)\leq-y\); hence, once \(G(s_n)\leq\frac12\),
\(n\log(1-G(s_n))\) lies between \(-nG(s_n)-(nG(s_n))^2/n\) and
\(-nG(s_n)\), and tends to \(-\lambda\).

\emph{Item 2.} Let \(W_n=n^{1/(p+1)}(X_{(1)}-m)/a\). By item~1, \(W_n\)
converges in law, so it is tight, and
\(\sqrt n\,(X_{(1)}-m)=a\,n^{1/2-1/(p+1)}W_n\) with
\(1/2-1/(p+1)<0\) when \(p<1\); the product of a tight sequence and a
deterministic null sequence tends to zero in probability. Let
\(Y_n=\sqrt n\,(T_n-m)\Rightarrow N(0,v)\). Since the normal distribution
function is continuous, \(P(|Y_n|\leq\varepsilon)\to
P(|N(0,v)|\leq\varepsilon)\) for each \(\varepsilon>0\). Then
\[
 P\bigl(|X_{(1)}-m|\geq|T_n-m|\bigr)
 \leq P\bigl(\sqrt n\,|X_{(1)}-m|>\varepsilon\bigr)
 +P(|Y_n|\leq\varepsilon),
\]
whose limit superior is at most \(P(|N(0,v)|\leq\varepsilon)\), which
tends to zero as \(\varepsilon\to0\) because \(v>0\). For the moment
estimate, \(m_\mu=\overline x-a(p+1)\) has
\(\sqrt n\,(m_\mu-m)\Rightarrow N(0,a^2(p+1))\) by \cref{thm:p15-clt-iid},
\(\operatorname{Var}\xi=p+1\) being finite for every \(p>-1\).

\emph{Item 3.} For \(p=1\), \(1/(p+1)=1/2\) and \(\Gamma(3)=2\), so item~1
is the stated limit. For \(p>1\), let \(\kappa=1/2-1/(p+1)>0\) and
\(M>0\). For any \(\delta>0\) and all \(n\) with \(Mn^{-\kappa}\leq\delta\),
\(P(\sqrt n\,(X_{(1)}-m)/a\leq M)=P(W_n\leq Mn^{-\kappa})\leq
P(W_n\leq\delta)\), whose limit is
\(1-\exp(-\delta^{p+1}/\Gamma(p+2))\), and this tends to zero as
\(\delta\to0\).
\end{proof}
\end{annexitem}

\begin{annexitem}{annex:p15-w02b-p15-057.linear-end}{thm:p15-w02b-p15-057-linear-end}{Location faster than \(\sqrt n\) at an end exponent equal to one}
\begin{proof}
Write \(Y\) for \(Y_1\) and fix \(\varepsilon\in(0,\delta)\).

\emph{The criterion.} Both pieces of \(g_\varepsilon\) equal
\(1/\varepsilon\) at \(y=\varepsilon\), so \(g_\varepsilon\) is continuous;
it is nonincreasing with values in \((0,1/\varepsilon]\), and it is
Lipschitz with constant \(\varepsilon^{-2}\), since it is constant on
\((-\infty,\varepsilon]\) and
\(|1/y-1/y'|=|y-y'|/(yy')\leq\varepsilon^{-2}|y-y'|\) for
\(y,y'>\varepsilon\). As \(P(Y>\varepsilon)\geq\int_\varepsilon^\delta
cy\,dy>0\) and \(g_\varepsilon(Y)<1/\varepsilon\) on \(\{Y>\varepsilon\}\),
\(0<\kappa_\varepsilon<1/\varepsilon\). Hence \(\psi_\varepsilon\) is
continuous and nondecreasing, equals \(\kappa_\varepsilon-1/\varepsilon<0\)
on \((-\infty,\varepsilon]\), and tends to \(\kappa_\varepsilon>0\) as
\(y\to\infty\). The integral \(\rho_\varepsilon\) of a continuous
nondecreasing function is continuous and convex, with both one-sided
derivatives equal to \(\psi_\varepsilon\). For \(y\leq0\),
\(\rho_\varepsilon(y)=(1/\varepsilon-\kappa_\varepsilon)|y|\to\infty\) as
\(y\to-\infty\); for \(y\geq y_0=2/\kappa_\varepsilon\),
\(\psi_\varepsilon(y)\geq\kappa_\varepsilon/2\), so
\(\rho_\varepsilon(y)\geq\rho_\varepsilon(y_0)+\kappa_\varepsilon(y-y_0)/2\to\infty\).
Thus \cref{def:p15-convex-location-criterion} applies, and by
\cref{thm:p15-convex-location-minimizers} the midpoint \(T_{n,\varepsilon}\)
is measurable and translation equivariant. Consequently
\(T_{n,\varepsilon}(X_1,\ldots,X_n)-m=T_{n,\varepsilon}(Y_1,\ldots,Y_n)\),
whose law does not depend on \(m\); it suffices to treat \(m=0\).

\emph{The population score.} Put
\(\lambda(t)=E\psi_\varepsilon(Y-t)=\kappa_\varepsilon-Eg_\varepsilon(Y-t)\),
finite since \(g_\varepsilon\) is bounded; \(\lambda(0)=0\). Let \(t>0\). On
\(\{Y>\varepsilon\}\), \(g_\varepsilon(Y)=1/Y\), while
\(g_\varepsilon(Y-t)\) equals \(1/\varepsilon>1/Y\) if \(Y-t\leq\varepsilon\)
and \(1/(Y-t)>1/Y\) otherwise; elsewhere
\(g_\varepsilon(Y-t)\geq g_\varepsilon(Y)\) by monotonicity. Since
\(P(Y>\varepsilon)>0\), \(Eg_\varepsilon(Y-t)>\kappa_\varepsilon\) and
\(\lambda(t)<0\). Let \(t<0\). On \(\{Y>\varepsilon\}\) also
\(Y-t>\varepsilon\), so \(g_\varepsilon(Y-t)=1/(Y-t)<1/Y=g_\varepsilon(Y)\);
elsewhere \(g_\varepsilon(Y-t)\leq g_\varepsilon(Y)\); so \(\lambda(t)>0\).
These are the sign conditions of \cref{thm:p15-convex-location-consistency}
with \(c=0\).

\emph{Derivative and second moment.} For \(t\neq0\),
\((\lambda(t)-\lambda(0))/t=-E[(g_\varepsilon(Y-t)-g_\varepsilon(Y))/t]\).
The integrand is bounded by \(\varepsilon^{-2}\) by the Lipschitz bound. At
every \(Y\neq\varepsilon\), \(g_\varepsilon\) is differentiable at \(Y\)
with derivative \(-Y^{-2}\mathbf 1\{Y>\varepsilon\}\), so the integrand
converges to \(Y^{-2}\mathbf 1\{Y>\varepsilon\}\) as \(t\to0\); and
\(Y\neq\varepsilon\) almost surely, because \(Y\) has a density. By dominated
convergence \(\lambda\) is differentiable at \(0\) with
\(\lambda'(0)=-E[Y^{-2}\mathbf 1\{Y>\varepsilon\}]=-J(\varepsilon)\), and
\(0<J(\varepsilon)\leq\varepsilon^{-2}\). Since \(0<\kappa_\varepsilon\leq
1/\varepsilon\) and \(0<g_\varepsilon\leq1/\varepsilon\),
\(|\psi_\varepsilon|\leq1/\varepsilon\), so
\(B(t)=E\psi_\varepsilon(Y-t)^2\leq\varepsilon^{-2}\) is finite for every
\(t\), and it is continuous at \(0\) by dominated convergence because
\(\psi_\varepsilon\) is continuous. Here
\(B(0)=E(g_\varepsilon(Y)-\kappa_\varepsilon)^2=\operatorname{Var}g_\varepsilon(Y)\).
\Cref{thm:p15-convex-location-normality} with \(c=0\) and
\(a=J(\varepsilon)\) gives
\(\sqrt n\,T_{n,\varepsilon}\Rightarrow N(0,v(\varepsilon))\).

\emph{The bound.} As \(f=0\) on \((-\infty,0]\) and \(f(y)\leq Cy\) on
\((0,\varepsilon)\),
\[
 \operatorname{Var}g_\varepsilon(Y)\leq Eg_\varepsilon(Y)^2
 =\varepsilon^{-2}P(Y\leq\varepsilon)+E[Y^{-2}\mathbf 1\{Y>\varepsilon\}]
 \leq\varepsilon^{-2}\int_0^\varepsilon Cy\,dy+J(\varepsilon)
 =\frac C2+J(\varepsilon),
\]
and \(J(\varepsilon)\geq\int_\varepsilon^\delta cy\,y^{-2}dy
=c\log(\delta/\varepsilon)\). Dividing by \(J(\varepsilon)^2\) gives the
bound on \(v(\varepsilon)\), which tends to \(0\) because
\(J(\varepsilon)\to\infty\).

\emph{Item 2.} By item~1 choose \(\varepsilon_k\in(0,\delta)\) decreasing
to \(0\) with \(v_k=v(\varepsilon_k)\leq k^{-3}\). The law \(N(0,v_k)\)
gives \(\{|z|>1/k\}\) probability at most \(k^2v_k\leq1/k\) by Chebyshev's
inequality, and gives the boundary \(\{|z|=1/k\}\) probability zero (when
\(v_k=0\) it is the point mass at \(0\)). Convergence in law therefore
gives \(N_k\) with \(P(\sqrt n\,|T_{n,\varepsilon_k}(Y_1,\ldots,Y_n)|>1/k)
\leq2/k\) for all \(n\geq N_k\); enlarging the \(N_k\) makes them strictly
increasing. They depend only on \(f\), because the law of
\(T_{n,\varepsilon}(Y_1,\ldots,Y_n)\) does. Each \(T_n\) so defined is
translation equivariant, being one of the \(T_{n,\varepsilon}\). Given
\(\eta>0\), choose \(k_0\) with \(2/k_0<\eta\). For \(n\geq N_{k_0}\), \(n\)
lies in a block \(N_k\leq n<N_{k+1}\) with \(k\geq k_0\), and
\(1/k\leq1/k_0<\eta\), so
\(P_m(\sqrt n\,|T_n-m|>\eta)\leq
P(\sqrt n\,|T_{n,\varepsilon_k}(Y_1,\ldots,Y_n)|>1/k)\leq2/k\leq2/k_0<\eta\).
Hence \(\sqrt n\,(T_n-m)\to0\) in probability.

\emph{Item 3.} For \(\eta>0\),
\(P_m(|T_n-m|\geq|S_n-m|)\leq P_m(\sqrt n|T_n-m|>\eta)
+P_m(\sqrt n|S_n-m|\leq\eta)\). The first term tends to \(0\) by item~2;
the second tends to \(P(|N(0,v)|\leq\eta)\leq2\eta/\sqrt{2\pi v}\), the
boundary having probability zero as \(v>0\). As \(\eta\) is arbitrary,
\(P_m(|T_n-m|<|S_n-m|)\to1\). By
\cref{def:p15-w02b-p15-057-efficiency} the relative efficiency of \(S\) with
respect to \(T_{\cdot,\varepsilon}\) is \(v(\varepsilon)/v\), which tends to
\(0\) by item~1. If \(\int y^2f<\infty\), then
\(\sqrt n(\overline X-\int yf-m)=n^{-1/2}\sum_i(Y_i-EY)\Rightarrow
N(0,\operatorname{Var}Y)\) by \cref{thm:p15-clt-iid}, and
\(\operatorname{Var}Y>0\) because \(Y\) has a density and so is not almost
surely constant.

\emph{Item 4.} The Type~III density of
\cref{sec:p15-w02b-p15-057-type-iii} at \(p=1\) is
\(a^{-1}((x-m)/a)e^{-(x-m)/a}\) on \(x>m\), that is \(f(y)=a^{-2}ye^{-y/a}\)
for \(y>0\); for \(0<y<a\), \(e^{-1}\leq e^{-y/a}\leq1\). For the Type~I
density with \(m_1=1\) put \(y=x+a_1\): then
\(f(y)=Ca_1^{-1}y\bigl((a_1+a_2-y)/a_2\bigr)^{m_2}\) for \(0<y<a_1+a_2\) and
\(f(y)=0\) otherwise. On \(0<y<\delta=(a_1+a_2)/2\) the last factor is
monotone in \(y\) and lies between its values at \(y=0\) and
\(y=\delta\), namely \(((a_1+a_2)/a_2)^{m_2}\) and
\(((a_1+a_2)/(2a_2))^{m_2}\), both positive; this gives \(c\) and \(C\). The
location of the curve is its left end \(-a_1\). If \(m_2=1\), the
observations \(-X_i\) have the Type~I density with the roles of
\((a_1,m_1)\) and \((a_2,m_2)\) exchanged.
\end{proof}
\end{annexitem}

\begin{annexitem}{annex:p15-w02b-p15-057.grouped-mle}{thm:p15-w02b-p15-057-grouped-mle}{Existence, consistency and limit of the grouped estimate}
\begin{proof}
Each \(p_j\) is positive and infinitely differentiable on \(\Theta\), and
\(\sum_jp_j=1\).

\emph{A compactness lemma.} Let \(0<\varepsilon<1\) and let \(S\) be a set
of cells containing a bounded cell \(j_0\) and two cells \(j<j'\) with
\(j'\geq j+2\). Then \(\mathcal C(\varepsilon,S)=\{\theta:p_i(\theta)\geq\varepsilon
\text{ for } i\in S\}\) is a compact subset of \(\Theta\). It is closed in
\(\Theta\) by continuity. Since \(\phi\leq(2\pi)^{-1/2}\),
\(p_{j_0}\leq(b_{j_0+1}-b_{j_0})/(\sigma\sqrt{2\pi})\), so on
\(\mathcal C(\varepsilon,S)\),
\(\sigma\leq\sigma^+=(b_{j_0+1}-b_{j_0})/(\varepsilon\sqrt{2\pi})\). Put
\(q=|\Phi^{-1}(\varepsilon)|\). From \(p_j\geq\varepsilon\),
\(\Phi(z_{j+1})\geq\varepsilon\) and \(z_{j+1}\geq-q\); from
\(p_{j'}\geq\varepsilon\), \(1-\Phi(z_{j'})\geq\varepsilon\) and
\(z_{j'}\leq q\). The endpoints \(b_{j+1}<b_{j'}\) are finite, because
\(1\leq j+1<j'\leq K-1\), and
\(z_{j'}-z_{j+1}=(b_{j'}-b_{j+1})/\sigma\leq2q\). If \(q=0\) the set is
empty; otherwise \(\sigma\geq\sigma^-=(b_{j'}-b_{j+1})/(2q)>0\), and the
two inequalities give \(b_{j'}-q\sigma^+\leq\mu\leq b_{j+1}+q\sigma^+\).
So \(\mathcal C(\varepsilon,S)\) is a closed subset of a compact rectangle inside
\(\Theta\).

\emph{Item 1.} Let \(S\) be the set of occupied cells, satisfying the
condition, fix \(\theta^*\in\Theta\) and put \(c=\ell_n(\theta^*)\leq0\).
Every term \(n_i\log p_i\) is nonpositive, so \(\ell_n(\theta)\geq c\)
implies \(n_i\log p_i(\theta)\geq c\), hence \(p_i(\theta)\geq e^{c/n_i}
\geq e^c\), for each \(i\in S\). The nonempty set \(\{\ell_n\geq c\}\) is
therefore a closed subset of the compact \(\mathcal C(e^c,S)\), and the continuous
\(\ell_n\) attains a maximum on it, which is its maximum on \(\Theta\).
If all \(n\) observations are in one cell \(j\), then \(\ell_n=n\log p_j<0\)
everywhere, because another cell has positive probability, while
\(\ell_n\to0\) along \(\sigma\to0\) with \(\mu\) the centre of a bounded
cell \(j\), or along \(\mu\to\mp\infty\) with \(\sigma=1\) for the tail
cell \(0\) or \(K-1\). If they are in two adjacent cells \(j,j+1\) with
counts \(n_j,n_{j+1}\geq1\), the function \((x,y)\mapsto n_j\log x+
n_{j+1}\log y\) is increasing in each argument, so on \(x+y\leq s<1\) it
is at most \(\ell^*+n\log s\), where
\(\ell^*=n_j\log(n_j/n)+n_{j+1}\log(n_{j+1}/n)\); as
\(p_j+p_{j+1}<1\) everywhere, \(\ell_n<\ell^*\) on \(\Theta\). The endpoint
\(b=b_{j+1}\) is finite. With \(t=\Phi^{-1}(n_j/n)\),
\(\sigma\to0\) and \(\mu=b-\sigma t\), one has \(z_{j+1}=t\) while
\(z_j\to-\infty\) and \(z_{j+2}\to\infty\) (these being \(\mp\infty\)
already when \(j=0\) or \(j+2=K\)), so \(p_j\to n_j/n\),
\(p_{j+1}\to n_{j+1}/n\) and \(\ell_n\to\ell^*\): the supremum is not
attained. Finally every cell has positive probability, so by
\cref{thm:p15-iid-series-strong-law} applied to its indicator it is
occupied for all large \(n\) almost surely; when cells \(0,1,2\) are
occupied, cell \(1\) is bounded and cells \(0,2\) are not adjacent.

\emph{Item 2: identification.} If \(p(\theta)=p(\theta')\), then
\(\Phi(z_1)=p_0\) and \(\Phi(z_2)=p_0+p_1\) agree for \(\theta,\theta'\);
\(\Phi\) being strictly increasing, \(z_1,z_2\) agree, and subtracting,
\((b_2-b_1)/\sigma=(b_2-b_1)/\sigma'\), so \(\sigma=\sigma'\) and then
\(\mu=\mu'\).

\emph{Item 2: consistency.} Put \(\pi_j=p_j(\theta_0)>0\),
\(\hat\pi_j=n_j/n\) and \(\Lambda(\theta)=\sum_j\pi_j\log p_j(\theta)\). By
concavity of \(\log\),
\(\Lambda(\theta)-\Lambda(\theta_0)=\sum_j\pi_j\log(p_j(\theta)/\pi_j)
\leq\log\sum_jp_j(\theta)=0\), with equality only if \(p_j(\theta)/\pi_j\)
is constant in \(j\), that is \(p(\theta)=p(\theta_0)\), that is
\(\theta=\theta_0\). Fix an outcome on which \(\hat\pi_j\to\pi_j\) for all
\(j\), an event of probability one by
\cref{thm:p15-iid-series-strong-law}, and let
\(C=1-\Lambda(\theta_0)\geq1\). For large \(n\),
\(n^{-1}\ell_n(\theta_0)=\sum_j\hat\pi_j\log\pi_j\geq-C\) and
\(\hat\pi_j\geq\pi_j/2\) for all \(j\); a maximum likelihood estimate then
exists by item~1, and any maximizer \(\widehat\theta_n\) satisfies
\(\hat\pi_j\log p_j(\widehat\theta_n)\geq n^{-1}\ell_n(\widehat\theta_n)
\geq-C\), hence \(p_j(\widehat\theta_n)\geq e^{-2C/\pi_j}\), for every
\(j\). Let \(\varepsilon=\min_je^{-2C/\pi_j}\) and
\(\mathcal K=\mathcal C(\varepsilon,\{0,\ldots,K-1\})\), compact by the lemma.
Since \(-\pi_j\log\pi_j\leq1/e<2C\), also \(\theta_0\in\mathcal K\). On
\(\mathcal K\), \(|\log p_j|\leq\log(1/\varepsilon)\), so
\(\sup_{\mathcal K}|n^{-1}\ell_n-\Lambda|\leq\delta_n=\log(1/\varepsilon)
\sum_j|\hat\pi_j-\pi_j|\to0\). Let \(U\) be an open neighbourhood of
\(\theta_0\). If \(\mathcal K\setminus U\) is nonempty, it is compact and
\(\Lambda<\Lambda(\theta_0)\) there, so
\(\eta=\Lambda(\theta_0)-\max_{\mathcal K\setminus U}\Lambda>0\). Now
\(\Lambda(\widehat\theta_n)\geq n^{-1}\ell_n(\widehat\theta_n)-\delta_n
\geq n^{-1}\ell_n(\theta_0)-\delta_n\geq\Lambda(\theta_0)-2\delta_n\), which
exceeds \(\Lambda(\theta_0)-\eta\) once \(2\delta_n<\eta\); then
\(\widehat\theta_n\in U\). This holds for every choice of maximizers.

\emph{Item 3.} Put \(u(z)=\sigma_0^{-1}\bigl(\phi(z),z\phi(z)\bigr)\), with
\(\sigma_0\) the scale component of \(\theta_0\) and
\(u(\pm\infty)=0\). By \cref{def:p15-fg-normal},
\(\nabla p_j(\theta_0)=u(z_j)-u(z_{j+1})\) for every cell, including the
tail cells. If \(t^{\top}Jt=\sum_j(t^{\top}\nabla p_j)^2/p_j=0\), then
\(t^{\top}\nabla p_j=0\) for all \(j\); starting from
\(t^{\top}u(z_1)=-t^{\top}\nabla p_0=0\) and stepping through the cells,
\(t^{\top}u(z_i)=0\) for every finite endpoint. So \(t\) is orthogonal to
\((1,z_1)\) and \((1,z_2)\), which are independent since \(z_1\neq z_2\),
and \(t=0\). Next, \cref{def:p15-w02b-p15-057-multi-regular} holds at
\(\theta_0\) for the model on the finite set of cells with counting
measure: the support is every cell for every \(\theta\);
\(\log p_j\) is twice continuously differentiable; \(\sum_j\nabla p_j=0\)
and \(\sum_j\nabla^2p_j=0\) because \(\sum_jp_j=1\); the score
\(s(j)=\nabla p_j/p_j\) has second-moment matrix
\(\sum_jp_j^{-1}\nabla p_j\nabla p_j^{\top}=J\); and on a closed ball about
\(\theta_0\) inside \(\Theta\) the finitely many continuous functions
\(\|\nabla^2\log p_j\|\) are bounded by a constant \(M\). Maximizers on the
event of existence are interior points of the open set \(\Theta\), hence
roots of \(\nabla\ell_n=0\); that event has probability tending to one by
item~1, and the maximizers converge to \(\theta_0\) almost surely by
item~2, hence in probability. \Cref{thm:p15-w02b-p15-057-multi-regular}
gives the stated limits.
\end{proof}
\end{annexitem}

\begin{annexitem}{annex:p15-w02b-p15-057.grouped-phase}{thm:p15-w02b-p15-057-grouped-phase}{A fixed-phase bound for grouped normal information}
\begin{proof}
Write \(Z=(X-\mu)/\sigma\). In \(Z\)-units the cells are intervals
\([u,v)\) of length \(a\) with midpoint \(c\), and the class of \(X\) is
the class of \(Z\). The location and scale scores are \(Z/\sigma\) and
\((Z^2-1)/\sigma\), with ungrouped information \(1/\sigma^2\) and
\(2/\sigma^2\). By \cref{thm:p15-fg-information},
\(\sigma^2I_\mu(\vartheta)=1-E\operatorname{Var}(Z\mid\text{class})\) and
\(\sigma^2I_\sigma(\vartheta)=2-E\operatorname{Var}(Z^2\mid\text{class})\).

\emph{A comparison lemma.} Let \(q=e^w/\int_u^ve^w\) be a density on an
interval \([u,v]\) of length \(a\) with
\(|w(x)-w(y)|\leq L|x-y|\), and let \(h\) be bounded and measurable. Write
\(V_q(h)\) and \(V_U(h)\) for the variances of \(h\) under \(q\) and under
the uniform law \(U\) on \([u,v]\). For \(x,y\in[u,v]\),
\(q(x)\leq e^{La}q(y)\); averaging in \(y\) gives \(q(x)\leq e^{La}/a\),
and likewise \(q(x)\geq e^{-La}/a\). A variance is the least mean square
deviation from a constant, so
\(V_q(h)\leq E_q(h-E_Uh)^2\leq e^{La}V_U(h)\) and
\(V_q(h)=E_q(h-E_qh)^2\geq e^{-La}E_U(h-E_qh)^2\geq e^{-La}V_U(h)\).
As \(1-e^{-x}\leq e^x-1\), \(|V_q(h)-V_U(h)|\leq(e^{La}-1)V_U(h)\).

\emph{Application.} Given its class \([u,v)\), \(Z\) has density
proportional to \(e^{-z^2/2}\) there, and \(w(z)=-z^2/2\) has
\(|w'|\leq L=\max(|u|,|v|)\). If \(Z\) lies in the class, then
\(L\leq|Z|+a\) and \(|c|\leq|Z|+a/2\). Under \(U\), \(z=c+t\) with \(t\)
uniform on \([-a/2,a/2]\), so \(V_U(z)=a^2/12\) and, as \(Et=Et^3=0\),
\(V_U(z^2)=4c^2Et^2+Et^4-(Et^2)^2=c^2a^2/3+a^4/80-a^4/144
=c^2a^2/3+a^4/180\). Taking expectations over the class of \(Z\),
\[
 \Bigl|E\operatorname{Var}(Z\mid\text{class})-\frac{a^2}{12}\Bigr|
 \leq\frac{a^2}{12}\,E\bigl[e^{a(|Z|+a)}-1\bigr]
 =\frac{a^2}{12}\bigl(2e^{3a^2/2}\Phi(a)-1\bigr),
\]
since \(Ee^{a|Z|}=2\int_0^\infty e^{az}\phi(z)\,dz=2e^{a^2/2}\Phi(a)\).
For the scale, with \(C\) the midpoint of the class of \(Z\),
\[
 \Bigl|E\operatorname{Var}(Z^2\mid\text{class})
 -\frac{a^2}3EC^2-\frac{a^4}{180}\Bigr|
 \leq E\Bigl[\Bigl(\frac{a^2(|Z|+a/2)^2}{3}+\frac{a^4}{180}\Bigr)
 \bigl(e^{a(|Z|+a)}-1\bigr)\Bigr].
\]
Since \(|C-Z|\leq a/2\), Minkowski's inequality gives
\(1-a/2\leq(EC^2)^{1/2}\leq1+a/2\); for \(a\leq2\) squaring yields
\(|EC^2-1|\leq a+a^2/4\). Combining gives the two displayed bounds.

\emph{Periodic remainders.} Each term \((\partial p_j)^2/p_j\) is
continuous in \(\vartheta\) and nonnegative, so \(I_\mu\) and \(I_\sigma\)
are measurable functions of the phase, bounded by the ungrouped values.
Every value \(\sigma^2I_\mu(\vartheta)\) lies within \(\varepsilon_\mu(a)\)
of \(1-a^2/12\), so the phase average does too, and the two differ by at
most \(2\varepsilon_\mu(a)\); likewise for the scale.

\emph{Order.} For \(a\leq1\),
\(2e^{3a^2/2}\Phi(a)-1=(2\Phi(a)-1)+2\Phi(a)(e^{3a^2/2}-1)\leq
a\sqrt{2/\pi}+3a^2e^{3/2}\), so \(\varepsilon_\mu(a)\leq\text{const}\cdot
a^3\). Using \(e^x-1\leq xe^x\), \(a\leq1\) and \(a^4\leq a^2\), the
expectation in \(\varepsilon_\sigma\) is at most
\(a^3E\bigl[\bigl((|Z|+1)^2/3+1/180\bigr)(|Z|+1)e^{|Z|+1}\bigr]\), which is
finite, and the other two terms of \(\varepsilon_\sigma\) are at most a
constant times \(a^3\).
\end{proof}
\end{annexitem}

\begin{annexitem}{annex:p15-w02b-p15-057.periodization}{thm:p15-w02b-p15-057-periodization}{Periodization and its Fourier series}
\begin{proof}
\emph{Continuity, period and coefficients.} Put
\(M_j=\sup_{0\leq u\leq L}|g(u+jL)|\). The series defining \(F\) converges
uniformly on \([0,L]\) by the comparison with \(\sum_jM_j\), so \(F\) is
continuous there; for \(u\in[nL,(n+1)L]\), reindexing the absolutely
convergent series gives \(F(u)=F(u-nL)\), so \(F\) has period \(L\) and is
continuous everywhere. Also
\(\int_{\mathbb R}|g|=\sum_j\int_0^L|g(u+jL)|du\leq L\sum_jM_j<\infty\), so
\(\widehat g\) is defined. Integrating the uniformly convergent series term
by term and using \(e^{-i\omega_kjL}=1\),
\[
 \frac1L\int_0^LF(u)e^{-i\omega_ku}du
 =\frac1L\sum_j\int_{jL}^{(j+1)L}g(x)e^{-i\omega_kx}dx
 =\frac1L\widehat g(\omega_k);
\]
for \(k=0\) this is the average.

\emph{Reconstruction.} If \(\sum_k|\widehat g(\omega_k)|<\infty\), the
series \(H(u)=L^{-1}\sum_k\widehat g(\omega_k)e^{i\omega_ku}\) converges
uniformly to a continuous function of period \(L\), whose coefficients,
computed term by term, are the same \(L^{-1}\widehat g(\omega_k)\). So
\(U(s)=(F-H)(Ls/(2\pi))\) is continuous with period \(2\pi\) and all its
Fourier coefficients vanish. For an integer \(R\geq1\) let
\(K_R(t)=R^{-1}|\sum_{j=0}^{R-1}e^{ijt}|^2=\sum_{|k|<R}(1-|k|/R)e^{ikt}\).
Then \(K_R\geq0\), \((2\pi)^{-1}\int_{-\pi}^{\pi}K_R=1\), and summing the
geometric series, \(K_R(t)=\sin^2(Rt/2)/(R\sin^2(t/2))\leq
1/(R\sin^2(\eta/2))\) for \(\eta\leq|t|\leq\pi\). The convolution
\((2\pi)^{-1}\int_{-\pi}^{\pi}U(s-t)K_R(t)dt\) is the finite combination
\(\sum_{|k|<R}(1-|k|/R)\hat U_ke^{iks}\) of the coefficients \(\hat U_k=0\),
so it vanishes, and
\[
 |U(s)|=\Bigl|\frac1{2\pi}\int_{-\pi}^{\pi}\bigl(U(s)-U(s-t)\bigr)K_R(t)\,dt\Bigr|
 \leq\sup_{|t|<\eta}|U(s)-U(s-t)|+\frac{2\sup|U|}{R\sin^2(\eta/2)}.
\]
Letting \(R\to\infty\) and then \(\eta\to0\), uniform continuity of \(U\)
gives \(U(s)=0\). Hence \(F=H\), and the bound on \(|F-\widehat g(0)/L|\)
follows by the triangle inequality.

\emph{The strip bound.} For \(\omega=0\), \(|\widehat g(0)|\leq\int|g|\leq
M\). Let \(\omega>0\); the case \(\omega<0\) is the same with \(+ic\) in
place of \(-ic\). The function \(z\mapsto g(z)e^{-i\omega z}\) is analytic
on the strip, which contains the closed rectangle with vertices
\(\pm R\) and \(\pm R-ic\). By Cauchy's theorem its integral around that
rectangle is zero. On the vertical sides \(|e^{-i\omega(\pm R-iy)}|=
e^{-\omega y}\leq1\) for \(0\leq y\leq c\), so those two integrals are at
most \(c\sup_{0\leq y\leq c}|g(\pm R-iy)|\), which tends to \(0\) as
\(R\to\infty\). Both horizontal integrals converge absolutely as
\(R\to\infty\), the lower one because
\(\int|g(x-ic)|dx\leq M\). Therefore
\(\widehat g(\omega)=\int_{\mathbb R}g(x-ic)e^{-i\omega(x-ic)}dx\), and since
\(|e^{-i\omega(x-ic)}|=e^{-\omega c}\), \(|\widehat g(\omega)|\leq Me^{-c\omega}\).
\end{proof}
\end{annexitem}

\begin{annexitem}{annex:p15-w02b-p15-057.grouped-exp}{thm:p15-w02b-p15-057-grouped-exp}{Fixed-phase grouped normal information to every order}
\begin{proof}
\emph{The information as a periodization.} Work in units
\(z=(\,\cdot-\mu)/\sigma\), let \(\phi(z)=(2\pi)^{-1/2}e^{-z^2/2}\) on the
whole complex plane, and put \(x_0=(b-\mu)/\sigma\). The class with
endpoints \(b+jh\) and \(b+(j+1)h\) is \([x,x+a)\) in \(z\)-units, with
\(x=x_0+ja\), and has probability
\(P(x)=\int_0^a\phi(x+s)\,ds\). Differentiating
\(\Phi((b_{j+1}-\mu)/\sigma)-\Phi((b_j-\mu)/\sigma)\),
\[
 \sigma\partial_\mu p_j=\phi(x)-\phi(x+a)=\int_0^a(x+s)\phi(x+s)\,ds=:N_\mu(x),
\]
\[
 \sigma\partial_\sigma p_j=x\phi(x)-(x+a)\phi(x+a)
 =-\int_0^a\bigl(1-(x+s)^2\bigr)\phi(x+s)\,ds=:N_\sigma(x),
\]
using \(\phi'(t)=-t\phi(t)\) and \((t\phi(t))'=(1-t^2)\phi(t)\). Hence
\(\sigma^2I_\mu=\sum_jg_\mu(x_0+ja)\), \(\sigma^2I_\sigma=\sum_jg_\sigma(x_0+ja)\)
and \(\sigma^2I_{\mu\sigma}=\sum_jg_{\mu\sigma}(x_0+ja)\), with
\(g_\mu=N_\mu^2/P\), \(g_\sigma=N_\sigma^2/P\) and
\(g_{\mu\sigma}=N_\mu N_\sigma/P\). As \(\vartheta\) runs over
\([0,2\pi)\), \(x_0\) runs over an interval of length \(a\), so the phase
averages are the averages of these periodizations over a period, namely
\(a^{-1}\widehat g(0)=a^{-1}\int g\), by
\cref{thm:p15-w02b-p15-057-periodization}.

\emph{The averages are \(J_\mu\), \(J_\sigma\) and \(0\).} In the
phase-randomized experiment the observation \(X+U\), with \(U\) uniform on
\((-h/2,h/2)\) and independent of \(X\), has density
\(q(y)=h^{-1}P(z-a/2)\) with \(z=(y-\mu)/\sigma\); differentiating as above,
\(\sigma\partial_\mu q=h^{-1}N_\mu(z-a/2)\) and
\(\sigma\partial_\sigma q=h^{-1}N_\sigma(z-a/2)\). So its information,
times \(\sigma^2\), is \(\int h^{-1}N^2/P\,dy=a^{-1}\int N^2/P\,dx\) for
each parameter: the phase averages above are \(J_\mu\) and \(J_\sigma\).
Under \(x\mapsto-x-a\), \(P\) and \(N_\sigma\) are unchanged and
\(N_\mu\) changes sign, so \(g_{\mu\sigma}\) is odd about \(-a/2\) and its
integral, which converges by the bounds below, is \(0\).

\emph{Continuation into a strip.} Let \(z=x+iy\) with \(|y|\leq c\) and
\(ca=\tau<\pi/2\). Then
\(\phi(x+s+iy)=\phi(x+s)e^{y^2/2}e^{-iy(x+s)}\), so
\(|\phi(x+s+iy)|=e^{y^2/2}\phi(x+s)\) and
\[
 \operatorname{Re}\bigl(e^{iyx}P(z)\bigr)
 =e^{y^2/2}\int_0^a\phi(x+s)\cos(ys)\,ds\geq e^{y^2/2}\cos\tau\,P(x)>0,
\]
since \(|ys|\leq\tau\). Thus \(|P(z)|\geq e^{y^2/2}\cos\tau\,P(x)\), and
\(P\), \(N_\mu\), \(N_\sigma\) being entire, the three functions \(g\) are
analytic on \(|\operatorname{Im}z|<\pi/(2a)\). Further
\(|N_\mu(z)|\leq e^{y^2/2}\int_0^a(|x+s|+c)\phi(x+s)ds\), and by the
Cauchy--Schwarz inequality the square of this integral is at most
\(P(x)\int_0^a(|x+s|+c)^2\phi(x+s)ds\). Therefore
\[
 |g_\mu(z)|\leq\frac{e^{c^2/2}}{\cos\tau}\int_0^a(|x+s|+c)^2\phi(x+s)\,ds,
\]
and in the same way, from \(|1-(x+s+iy)^2|\leq1+(|x+s|+c)^2\),
\[
 |g_\sigma(z)|\leq\frac{e^{c^2/2}}{\cos\tau}\int_0^a\bigl(1+(|x+s|+c)^2\bigr)^2
 \phi(x+s)\,ds,
\]
while \(|g_{\mu\sigma}|\leq(|N_\mu|^2+|N_\sigma|^2)/(2|P|)\) is at most half
the sum of the two right-hand sides. These bounds tend to \(0\) as
\(|x|\to\infty\) uniformly in \(|y|\leq c\). Integrating in \(x\) and
substituting \(t=x+s\) (Tonelli),
\(\int|g_\mu(x+iy)|dx\leq M_\mu=ae^{c^2/2}E(|Z|+c)^2/\cos\tau\) and
\(\int|g_\sigma(x+iy)|dx\leq M_\sigma=ae^{c^2/2}E(1+(|Z|+c)^2)^2/\cos\tau\).
With \(E|Z|=\beta\), \(EZ^2=1\), \(E|Z|^3=2\beta\) and \(EZ^4=3\),
\(E(|Z|+c)^2=1+2\beta c+c^2\) and
\(E(1+(|Z|+c)^2)^2=1+2E(|Z|+c)^2+E(|Z|+c)^4=6+12\beta c+8c^2+4\beta c^3+c^4\).
At \(y=0\) the same bounds give
\(\sup_{0\leq u\leq a}|g(u+ja)|\leq a\sup_{t\in[ja,(j+2)a]}(1+(|t|+c)^2)^2\phi(t)\),
which is summable over \(j\).

\emph{Conclusion.} By \cref{thm:p15-w02b-p15-057-periodization} with
\(L=a\), \(\omega_k=2\pi k/a\) and \(|\widehat g(\omega_k)|\leq
Me^{-2\pi c|k|/a}=M\eta^{|k|}\),
\[
 |\sigma^2I_\mu(\vartheta)-J_\mu|\leq\frac2a\sum_{k\geq1}M_\mu\eta^k
 =\frac{2M_\mu}a\,\frac\eta{1-\eta}=B_\mu,
\]
and likewise for \(\sigma\); for the cross entry the average is \(0\) and
\(M\leq(M_\mu+M_\sigma)/2\). With \(\tau=\pi/3\),
\(e^{c^2/2}\eta=e^{\pi^2/(18a^2)-2\pi^2/(3a^2)}=e^{-11\pi^2/(18a^2)}\), and for
\(a\leq1\), \(1/(1-\eta)\leq1/(1-e^{-2\pi^2/3})\) while the remaining factors
are polynomials in \(c=\pi/(3a)\). Since
\(a^{-N'-10}e^{-11\pi^2/(18a^2)}\to0\) as \(a\to0\) for every fixed
\(N'\), \(B_\mu,B_\sigma\leq a^{10}\) for all small \(a\). Adding this to the remainder of
\cref{thm:p15-fg-expansion} gives its expansions of \(J_\mu\) and
\(J_\sigma\) for \(\sigma^2I_\mu(\vartheta)\) and
\(\sigma^2I_\sigma(\vartheta)\) with an \(O(a^{10})\) remainder uniform in
phase. For the inverses, when \(J_\mu\geq1/2\) and \(B_\mu\leq1/4\) (all
small \(a\)), \(\sigma^2I_\mu(\vartheta)\geq1/4\) and
\(|1/(\sigma^2I_\mu)-1/J_\mu|=|\sigma^2I_\mu-J_\mu|/(\sigma^2I_\mu J_\mu)\leq
8B_\mu\); the same argument with \(J_\sigma\geq1\) and \(B_\sigma\leq1/2\)
gives \(|2/(\sigma^2I_\sigma)-2/J_\sigma|\leq4B_\sigma\). Both are
\(O(a^{10})\), which transfers the expansions of \(J_\mu^{-1}\) and
\(2J_\sigma^{-1}\).
\end{proof}
\end{annexitem}

\begin{annexitem}{annex:p15-w02b-p15-057.chisq-sampling}{thm:p15-w02b-p15-057-chisq-sampling}{Twice Fisher's \(L\) and Pearson's \(\chi^2\) under sampling with fitted cells}
\begin{proof}
Each \(n_s\) has the binomial law with index \(N\) and probability
\(\pi_s\), so \(E(n_s-N\pi_s)^2=N\pi_s(1-\pi_s)\) and, by Chebyshev's
inequality, \((n_s-N\pi_s)/\sqrt N\) is tight. Put \(x_s=n_s-m_s\). Then
\[
 \frac{x_s}{\sqrt N}=\frac{n_s-N\pi_s}{\sqrt N}-\sqrt N(\hat p_{N,s}-\pi_s)
\]
is tight, being a difference of tight sequences, and
\(\hat p_{N,s}-\pi_s=N^{-1/2}\cdot\sqrt N(\hat p_{N,s}-\pi_s)\to0\) in
probability. Let \(G_N\) be the event that \(\hat p_{N,s}>\pi_s/2\) for
every \(s\); its probability tends to one, and on it every \(m_s\geq
N\pi_s/2>0\), so \(D\) and \(Q\) are defined. On \(G_N\),
\[
 Q=\sum_s\frac{x_s^2}{m_s}\leq\sum_s\frac2{\pi_s}\Bigl(\frac{x_s}{\sqrt N}
 \Bigr)^2,\qquad
 \rho=\max_s\frac{|x_s|}{m_s}\leq\max_s\frac2{\pi_s\sqrt N}\cdot
 \frac{|x_s|}{\sqrt N},
\]
so \(Q\) is tight and \(\rho\to0\) in probability. Because
\(\sum_sn_s=\sum_sm_s=N\), \cref{thm:p15-fg-deviance} applies on
\(G_N\cap\{\rho\leq\frac12\}\), an event of probability tending to one,
and gives \(|2D-Q|\leq\rho Q/(3(1-\rho)^2)\leq\frac43\rho Q\). Given
\(\varepsilon,\gamma>0\) choose \(M\) with \(P(Q>M)<\varepsilon\) for all
large \(N\); then \(P(|2D-Q|>\gamma)\leq\varepsilon+P(\rho>3\gamma/(4M))
+P(G_N^c)+P(\rho>\tfrac12)\), whose last three terms tend to zero. So
\(2D-Q\to0\) in probability. With \(m_s=N\pi_s\) take
\(\hat p_{N,s}=\pi_s\), for which the premises hold trivially.
\end{proof}
\end{annexitem}

\begin{annexitem}{annex:p15-w02b-p15-057.lindeberg-vector}{thm:p15-w02b-p15-057-lindeberg-vector}{A vector central limit statement for smooth functions}
\begin{proof}
\emph{Gaussian comparison sequence.} On a product space let
\(G_1,G_2,\ldots\) be independent random vectors with independent standard
normal coordinates, independent of the \(\xi_i\). Each coordinate of
\(W_N=N^{-1/2}\sum_{i\leq N}G_i\) is a sum of \(N\) independent
\(\operatorname{normal}(0,1/N)\) variables, hence standard normal by
\cref{thm:p15-normal-sums}, and distinct coordinates are functions of
disjoint independent families, hence independent. So \(W_N\) has the law of
\(G\).

\emph{Constants.} Let \(\|h\|\), \(H_2\), \(H_3\) bound \(|h|\) and the
absolute values of all second and third partial derivatives. The first
derivatives are bounded too: for a unit vector \(e\) and \(t>0\), Taylor's
formula gives \(|h(v+te)-h(v)-t\,\partial_eh(v)|\leq rH_2t^2/2\), so
\(|\partial_eh(v)|\leq2\|h\|/t+rH_2t/2\). For \(w\in\mathbb R^r\), Taylor's
formula with remainder gives
\(h(R+w)=h(R)+\nabla h(R)^{\top}w+\tfrac12w^{\top}\nabla^2h(R)w+e(R,w)\) with
\(|e(R,w)|\leq H_3(\sum_l|w_l|)^3/6\leq H_3r^{3/2}|w|^3/6\).

\emph{Replacement.} For \(1\leq k\leq N\) put
\(R_k=N^{-1/2}(\sum_{i<k}G_i+\sum_{k<i\leq N}\xi_i)\), independent of
\((\xi_k,G_k)\). Then \(R_1+\xi_1/\sqrt N=V_N\),
\(R_k+G_k/\sqrt N=R_{k+1}+\xi_{k+1}/\sqrt N\) for \(k<N\), and
\(R_N+G_N/\sqrt N=W_N\), so
\[
 Eh(V_N)-Eh(W_N)=\sum_{k=1}^N\bigl(Eh(R_k+\xi_k/\sqrt N)-Eh(R_k+G_k/\sqrt N)\bigr).
\]
Expand both terms about \(R_k\). By independence and boundedness of
\(\nabla h\) and \(\nabla^2h\),
\(E[\nabla h(R_k)^{\top}\xi_k]=E[\nabla h(R_k)]^{\top}E\xi_k=0\), likewise for
\(G_k\), and
\(E[\xi_k^{\top}\nabla^2h(R_k)\xi_k]=E\operatorname{tr}(\nabla^2h(R_k))
=E[G_k^{\top}\nabla^2h(R_k)G_k]\), because both vectors have second-moment
matrix \(I_r\). So each summand is at most
\(H_3r^{3/2}(E|\xi_1|^3+E|G_1|^3)/(6N^{3/2})\) in absolute value, with
\(E|\xi_1|^3\leq b^3\) and \(E|G_1|^3\leq(E|G_1|^4)^{3/4}=(r^2+2r)^{3/4}\).
The sum is \(O(N^{-1/2})\), so \(Eh(V_N)-Eh(G)=Eh(V_N)-Eh(W_N)\to0\).

\emph{The squared length.} Let \(S(t)=0\) for \(t\leq0\),
\(S(t)=35t^4-84t^5+70t^6-20t^7\) for \(0\leq t\leq1\) and \(S(t)=1\) for
\(t\geq1\). Since \(S'(t)=140t^3(1-t)^3\) on \([0,1]\), \(S\) is
nondecreasing and three times continuously differentiable on \(\mathbb R\).
For real \(x\) and \(\eta>0\) put
\(h^+(v)=1-S((|v|^2-x)/\eta)\) and \(h^-(v)=1-S((|v|^2-x+\eta)/\eta)\).
They are three times continuously differentiable, take values in
\([0,1]\), and their partial derivatives vanish outside the bounded set
\(\{|v|^2\leq x+\eta\}\), so they are bounded; and
\(\mathbf 1\{|v|^2\leq x-\eta\}\leq h^-\leq\mathbf 1\{|v|^2\leq x\}\leq
h^+\leq\mathbf 1\{|v|^2\leq x+\eta\}\). Hence
\[
 F_r(x-\eta)\leq Eh^-(G)=\lim Eh^-(V_N)\leq\liminf P(|V_N|^2\leq x)
 \leq\limsup P(|V_N|^2\leq x)\leq Eh^+(G)\leq F_r(x+\eta),
\]
where \(|G|^2\) has law \(\chi^2_r\) by
\cref{thm:p15-student-beta-direction}. That law has a density
(\cref{def:p15-gamma-density}), so \(F_r\) is continuous, and letting
\(\eta\to0\) gives \(P(|V_N|^2\leq x)\to F_r(x)\).
\end{proof}
\end{annexitem}

\begin{annexitem}{annex:p15-w02b-p15-057.chisq-law}{thm:p15-w02b-p15-057-chisq-law}{The limiting law of \(\chi^2\) and \(2L\) with fitted cells}
\begin{proof}
\emph{Geometry.} Let \(u=(\sqrt{\pi_s})_s\), a unit vector. Since
\(\sum_sp_s\equiv1\), \(\sum_s\nabla p_s(\theta_0)=0\), that is
\(A^{\top}u=0\): the columns of \(A\) lie in \(u^\perp\), which has dimension
\(K-1\), so \(q\leq K-1\) and \(r\geq0\). Let \(P_A=A(A^{\top}A)^{-1}A^{\top}\)
(\(P_A=0\) if \(q=0\)), the orthogonal projection onto the column space of
\(A\), and \(M=I-P_A\). Then \(Mu=u\). The subspace
\(H=\{v:A^{\top}v=0,\ u^{\top}v=0\}\) has dimension \(K-q-1=r\), because
\(u\) is a nonzero vector orthogonal to the \(q\)-dimensional column space.

\emph{The counts.} Write \(Z_N=N^{-1/2}\sum_{i\leq N}\zeta_i\) with
\(\zeta_{i,s}=(\mathbf 1\{\text{observation }i\text{ in cell }s\}-\pi_s)/\sqrt{\pi_s}\).
The \(\zeta_i\) are independent and identically distributed, bounded by
\(\max_s\pi_s^{-1/2}\) in each coordinate, with \(E\zeta_i=0\) and
\(E\zeta_i\zeta_i^{\top}=I-uu^{\top}\), since
\(E[\mathbf 1_s\mathbf 1_t]=\pi_s\delta_{st}\); and
\(u^{\top}\zeta_i=\sum_s(\mathbf 1_s-\pi_s)=0\). In particular
\(EZ_{N,s}^2=1-\pi_s\leq1\), so each \(Z_{N,s}\) is tight by Chebyshev's
inequality.

\emph{Linearization.} From the premise and tightness of \(Z_N\),
\(\sqrt N(\widehat\theta_N-\theta_0)\) is tight, so
\(\widehat\theta_N\to\theta_0\) in probability. Let \(B\) be a closed ball
about \(\theta_0\) on which \(p\) is continuously differentiable. By
continuity, \(p_s(\widehat\theta_N)\to\pi_s\) in probability, so the event
that every \(p_s(\widehat\theta_N)>0\) and \(\widehat\theta_N\in B\) has
probability tending to one, and \(\pi_s/p_s(\widehat\theta_N)\to1\) in
probability. On that event the mean value theorem gives points
\(\tilde\theta_s\) on the segment from \(\theta_0\) to \(\widehat\theta_N\)
with \(p_s(\widehat\theta_N)-\pi_s=\nabla p_s(\tilde\theta_s)^{\top}
(\widehat\theta_N-\theta_0)\), so
\[
 \frac{\sqrt N(p_s(\widehat\theta_N)-\pi_s)}{\sqrt{\pi_s}}
 =A_s^{\top}\sqrt N(\widehat\theta_N-\theta_0)+\epsilon_s,\qquad
 |\epsilon_s|\leq\frac{|\nabla p_s(\tilde\theta_s)-\nabla p_s(\theta_0)|}
 {\sqrt{\pi_s}}\sqrt N|\widehat\theta_N-\theta_0|,
\]
where \(A_s^{\top}\) is row \(s\) of \(A\); \(\epsilon_s\to0\) in probability
by continuity of \(\nabla p_s\) and tightness. With the premise, the vector
\(e_N\) with \(e_{N,s}=(n_s-Np_s(\widehat\theta_N))/\sqrt{N\pi_s}\) satisfies
\(e_N=Z_N-A(A^{\top}A)^{-1}A^{\top}Z_N+o_p(1)=MZ_N+o_p(1)\).

\emph{The statistic.} On the event above,
\(Q=\sum_se_{N,s}^2\,\pi_s/p_s(\widehat\theta_N)\). Since \(MZ_N\) is tight,
\(e_N=MZ_N+o_p(1)\) and \(\pi_s/p_s(\widehat\theta_N)\to1\),
\(Q-|MZ_N|^2\to0\) in probability. Each \(M\zeta_i\) lies in \(H\):
\(A^{\top}M\zeta_i=0\), and \(u^{\top}M\zeta_i=(Mu)^{\top}\zeta_i=
u^{\top}\zeta_i=0\). If \(r=0\), then \(H=\{0\}\), \(MZ_N=0\) and
\(Q\to0\) in probability. If \(r\geq1\), let \(e_1,\ldots,e_r\) be an
orthonormal basis of \(H\) and \(\xi_i=(e_l^{\top}\zeta_i)_{l\leq r}\).
These are independent, identically distributed and bounded, with mean
\(0\) and second-moment matrix \((e_l^{\top}(I-uu^{\top})e_k)_{lk}=I_r\).
As \(MZ_N\in H\) and \(Me_l=e_l\),
\(|MZ_N|^2=\sum_l(e_l^{\top}MZ_N)^2=\sum_l(e_l^{\top}Z_N)^2=|V_N|^2\) with
\(V_N=N^{-1/2}\sum_i\xi_i\). By
\cref{thm:p15-w02b-p15-057-lindeberg-vector},
\(P(|V_N|^2\leq x)\to F_r(x)\). For \(\eta>0\),
\(P(Q\leq x)\leq P(|V_N|^2\leq x+\eta)+P(|Q-|V_N|^2|>\eta)\) and
\(P(Q\leq x)\geq P(|V_N|^2\leq x-\eta)-P(|Q-|V_N|^2|>\eta)\); letting
\(N\to\infty\) and then \(\eta\to0\), continuity of \(F_r\) gives
\(P(Q\leq x)\to F_r(x)\).

\emph{Twice \(L\).} The fitted probabilities \(\hat p_{N,s}=p_s(\widehat
\theta_N)\) sum to one, and \(\sqrt N(\hat p_{N,s}-\pi_s)\) is tight by the
linearization. \Cref{thm:p15-w02b-p15-057-chisq-sampling} gives
\(2D-Q\to0\) in probability, so \(2D\) has the same limit, by the same
sandwich.

\emph{Item 2.} With \(m_s=N\pi_s\), \(e_N=Z_N\) exactly and \(Q=|Z_N|^2\);
the argument above with \(M=I\) and \(H=u^\perp\) gives \(r=K-1\), and
\cref{thm:p15-w02b-p15-057-chisq-sampling} covers \(2D\).

\emph{Item 3.} The proof of item~3 of
\cref{thm:p15-w02b-p15-057-grouped-mle} verifies
\cref{def:p15-w02b-p15-057-multi-regular} for the grouped model, shows that
maximizers exist with probability tending to one, are then roots of the
likelihood equations, and converge to \(\theta_0\) in probability. By
\cref{thm:p15-w02b-p15-057-multi-regular} and item~2 of
\cref{thm:p15-w02b-p15-057-multi},
\(\sqrt N(\widehat\theta_N-\theta_0)-J^{-1}N^{-1/2}\sum_{i\leq N}s(X_i)\to0\)
in probability, where \(s(j)=\nabla p_j(\theta_0)/\pi_j\) for an
observation in cell \(j\) and \(J=\sum_j\pi_j^{-1}\nabla p_j\nabla p_j^{\top}
=A^{\top}A\), positive definite by that theorem, so that \(A\) has rank
\(q=2\). Since \(\sum_j\nabla p_j(\theta_0)=0\),
\(N^{-1/2}\sum_is(X_i)=\sum_jn_j\nabla p_j/(\sqrt N\pi_j)=
\sum_j(n_j-N\pi_j)\nabla p_j/(\sqrt N\pi_j)=A^{\top}Z_N\). This is the
premise, and \(r=K-1-2\).
\end{proof}
\end{annexitem}

\begin{annexitem}{annex:p15-w02b-p15-057.poisson}{thm:p15-w02b-p15-057-poisson}{The Poisson series}
\begin{proof}
\emph{Moments.} Shifting the summation index,
\(E_mX=\sum_{x\geq1}e^{-m}m^x/(x-1)!=m\) and
\(E_mX(X-1)=\sum_{x\geq2}e^{-m}m^x/(x-2)!=m^2\), so
\(\operatorname{Var}_mX=m^2+m-m^2=m\). All moments are finite, since
\(x^k\,m^x/x!\) is summable for each \(k\).

\emph{Item 1.} With \(t=\sum_iX_i\),
\(\ell_n(m)=-nm+t\log m-\sum_i\log X_i!\). If \(t>0\),
\(\ell_n'(m)=-n+t/m\) and \(\ell_n''(m)=-t/m^2<0\), so \(\ell_n\) is
strictly concave on \((0,\infty)\) with the single critical point
\(m=t/n=\overline X\), which is its unique maximizer; the equation
\(\ell_n'=0\) is \(\sum_i(X_i/m-1)=0\). If \(t=0\),
\(\ell_n(m)=-nm\) is strictly decreasing, so its supremum over
\((0,\infty)\) is approached as \(m\downarrow0\) and not attained.

\emph{Item 2.} With respect to counting measure on
\(\{0,1,2,\ldots\}^n\) the joint density is
\(e^{-nm}m^{t}\cdot\prod_i(X_i!)^{-1}\), a function of \(m\) and \(t\)
times a function of the sample alone, so
\cref{thm:p15-w02b-p15-057-factorization} gives sufficiency of \(t\).

\emph{Item 3.} (R1): every \(f_m(x)=e^{-m}m^x/x!\) is positive on
\(A=\{0,1,2,\ldots\}\). (R2): \(\log f_m(x)=-m+x\log m-\log x!\) is smooth
in \(m>0\), with \(s_m(x)=x/m-1\) and \(\partial_m^2\log f_m(x)=-x/m^2\).
(R3): \(\partial_mf_m=f_m\,(x/m-1)\) and
\(\partial_m^2f_m=f_m\bigl((x/m-1)^2-x/m^2\bigr)\) are absolutely summable
over \(x\) because the first two moments are finite, and their sums are
\(E_mX/m-1=0\) and \(\operatorname{Var}_mX/m^2-E_mX/m^2=1/m-1/m=0\).
(R4): \(I(m)=E_m(X/m-1)^2=\operatorname{Var}_mX/m^2=1/m\in(0,\infty)\).
(R5): for \(|m'-m|<m/2\), \(|\partial_{m'}^2\log f_{m'}(x)|=x/m'^2\leq
4x/m^2=M(x)\), and \(E_mM=4/m<\infty\). So the family is regular at \(m\).
By the moments, \(E_m\overline X=m\) and
\(\operatorname{Var}_m\overline X=m/n=1/(nI(m))\), so the finite-sample
information efficiency \(1/(nI(m)\operatorname{Var}_m\overline X)\) is one.
The limit law is \cref{thm:p15-clt-iid} with \(\sigma^2=m\).
\end{proof}
\end{annexitem}

\begin{annexitem}{annex:p15-w02b-p15-057.dilution}{thm:p15-w02b-p15-057-dilution}{Maximum likelihood and the sterile count in a dilution series}
\begin{proof}
\emph{Derivatives.} Fix \(x\) and write \(m=m_x=e^\theta a^{-x}\), so that
\(dm/d\theta=m\), \(\log p_x=-m\) and \(d\log p_x/d\theta=\log p_x\). Then
\(dp_x/d\theta=p_x\log p_x=-dq_x/d\theta\), which is Fisher's
\(\partial p/\partial\log n\), and
\[
 \frac{d^2\log p_x}{d\theta^2}=-m<0,\qquad
 \frac{d\log q_x}{d\theta}=-\frac{p_x\log p_x}{q_x}=\frac{m}{e^m-1},\qquad
 \frac{d^2\log q_x}{d\theta^2}
 =m\,\frac{e^m-1-me^m}{(e^m-1)^2}<0,
\]
the last because \(k(m)=e^m-1-me^m\) has \(k(0)=0\) and
\(k'(m)=-me^m<0\) for \(m>0\).

\emph{Replicates.} Number the plates so that replicate \(r\in\{1,\ldots,s\}\)
contains one plate at each \(x\in D\), and let \(B_{r,x}\) indicate that
this plate is sterile. The vectors \(B_r=(B_{r,x})_{x\in D}\) are
independent and identically distributed, with density
\(f_\theta(b)=\prod_xp_x^{b_x}q_x^{1-b_x}\) with respect to counting measure
on \(\{0,1\}^D\), and
\(\ell_s(\theta)=\sum_r\log f_\theta(B_r)=\sum_x\bigl(S_x\log p_x
+(s-S_x)\log q_x\bigr)\).

\emph{Item 1 and regularity.} (R1): \(0<p_x<1\), so \(f_\theta(b)>0\) for
every \(b\) and \(\theta\). (R2): \(\log f_\theta(b)\) is smooth in
\(\theta\). (R3): the sample space is finite, so \(\sum_b\partial_\theta
f_\theta(b)=\partial_\theta\sum_bf_\theta(b)=0\), and likewise for the
second derivative. (R4): the score is
\(\sum_x\bigl(b_x\log p_x-(1-b_x)p_x\log p_x/q_x\bigr)
=\sum_x(b_x-p_x)\log p_x/q_x\), a sum of independent centred terms, so
\(I(\theta)=\sum_xp_xq_x(\log p_x)^2/q_x^2=\sum_xw(p_x)=\mathcal W_D\);
each term is positive since \(\log p_x=-m_x\neq0\). The same computation
for one plate gives \(w(p_x)\), and information adds over the
\(s|D|\) independent plates, giving \(s\mathcal W_D\). (R5):
\(\partial_\theta^2\log f_\theta(b)\) is continuous in \(\theta\) for each
of the finitely many \(b\), so it is bounded on \(|\theta-\theta_0|\leq1\)
by a constant \(M\).

\emph{Item 2.} By the derivatives, \(\ell_s\) is a combination with
nonnegative coefficients \(S_x\), \(s-S_x\), summing to \(s|D|>0\), of
strictly concave functions, hence strictly concave. As
\(\theta\to\infty\), \(\log p_x\to-\infty\) and \(\log q_x\to0\); as
\(\theta\to-\infty\), \(\log p_x\to0\) and \(\log q_x\to-\infty\). So
\(\ell_s\to-\infty\) at both ends exactly when some \(S_x>0\) and some
\(S_x<s\); then the continuous strictly concave \(\ell_s\) has a unique
maximizer, the unique zero of \(\ell_s'\), and \(\ell_s'\) is Fisher's
expression by the derivative formulas. If every plate is sterile,
\(\ell_s=\sum_xS_x\log p_x\) is strictly decreasing; if every plate is
fertile, \(\ell_s\) is strictly increasing; in neither case is there a
maximizer. By the strong law, \(S_x/s\to p_x(\theta_0)\in(0,1)\) almost
surely for each \(x\), so almost surely the first case holds for all large
\(s\). For consistency, put
\[
 \lambda(\theta)=\sum_x\Bigl(p_x(\theta_0)\frac{d\log p_x}{d\theta}(\theta)
 +q_x(\theta_0)\frac{d\log q_x}{d\theta}(\theta)\Bigr),
\]
the derivative of the strictly concave function
\(\sum_x(p_x(\theta_0)\log p_x+q_x(\theta_0)\log q_x)\), hence strictly
decreasing, with \(\lambda(\theta_0)=\sum_x(dp_x/d\theta+dq_x/d\theta)=0\)
at \(\theta_0\). Since \(\ell_s'(\theta)/s\) is the same expression with
\(S_x/s\) and \(1-S_x/s\) in place of \(p_x(\theta_0)\), \(q_x(\theta_0)\),
almost surely \(\ell_s'(\theta)/s\to\lambda(\theta)\) at each \(\theta\).
For \(\varepsilon>0\), \(\lambda(\theta_0-\varepsilon)>0>
\lambda(\theta_0+\varepsilon)\), so almost surely, for large \(s\),
\(\ell_s'(\theta_0-\varepsilon)>0>\ell_s'(\theta_0+\varepsilon)\), and the
unique zero of the decreasing function \(\ell_s'\) lies in
\((\theta_0-\varepsilon,\theta_0+\varepsilon)\). Intersecting over
\(\varepsilon=1/k\) gives \(\widehat\theta\to\theta_0\) almost surely.
Define \(\widehat\theta=0\) when no maximizer exists; then
\(\widehat\theta\) is a function of the counts, hence measurable,
\(\widehat\theta\to\theta_0\) in probability and
\(P(\ell_s'(\widehat\theta)=0)\to1\).
\cref{thm:p15-w02b-p15-057-asymptotic}, with the regularity just verified
and \(s\) in place of \(n\), gives
\(\sqrt s(\widehat\theta-\theta_0)\Rightarrow N(0,1/\mathcal W_D)\).

\emph{Item 3.} Let \(\mu(\theta)=\sum_xp_x(\theta)\). It is continuously
differentiable with \(\mu'(\theta)=\sum_xp_x\log p_x=-\mathcal
D_D(\theta)<0\), tends to \(|D|\) as \(\theta\to-\infty\) and to \(0\) as
\(\theta\to\infty\), so it is a strictly decreasing bijection of
\(\mathbb R\) onto \((0,|D|)\) with continuous inverse. Hence the equation
\(s\mu(\widetilde\theta)=T\) has exactly one solution when \(0<T<s|D|\)
and none otherwise; put \(\widetilde\theta=0\) in the second case. The
replicate counts \(Y_r=\sum_xB_{r,x}\) are independent and identically
distributed with mean \(\mu(\theta_0)\) and variance
\(\sum_xp_xq_x=\mathcal V_D\) at \(\theta_0\), since the plates within a
replicate are independent, and \(T=\sum_rY_r\). By the strong law
\(T/s\to\mu(\theta_0)\in(0,|D|)\) almost surely, so the solution exists for
all large \(s\) and \(\widetilde\theta=\mu^{-1}(T/s)\to\theta_0\) almost
surely. By the mean value theorem, on the event that it exists,
\(T/s-\mu(\theta_0)=\mu'(\xi_s)(\widetilde\theta-\theta_0)\) for some
\(\xi_s\) between \(\theta_0\) and \(\widetilde\theta\); so
\(\mu'(\xi_s)\to\mu'(\theta_0)=-\mathcal D_D\) almost surely. By
\cref{thm:p15-clt-iid}, \(\sqrt s(T/s-\mu(\theta_0))\Rightarrow
N(0,\mathcal V_D)\), and by \cref{thm:p15-clt-perturbation}, dividing by
\(\mu'(\xi_s)\) (and absorbing the event of nonexistence, whose
probability tends to zero),
\(\sqrt s(\widetilde\theta-\theta_0)\Rightarrow
N(0,\mathcal V_D/\mathcal D_D^2)\).

\emph{Item 4.} The relative efficiency is the ratio of limit variances,
\((1/\mathcal W_D)/(\mathcal V_D/\mathcal D_D^2)\). By the Cauchy--Schwarz
inequality,
\[
 \mathcal D_D^2=\Bigl(\sum_x\sqrt{p_xq_x}\cdot\sqrt{p_x/q_x}\,
 (-\log p_x)\Bigr)^2\leq\sum_xp_xq_x\sum_x\frac{p_x}{q_x}(\log p_x)^2
 =\mathcal V_D\mathcal W_D,
\]
with equality exactly when the vectors \((\sqrt{p_xq_x})_x\) and
\((\sqrt{p_x/q_x}\,m_x)_x\) are proportional, that is, when
\(m_x/q_x=m_x/(1-e^{-m_x})\) is the same for every \(x\in D\). The function
\(m\mapsto m/(1-e^{-m})\) is strictly increasing on \((0,\infty)\) (its
derivative has numerator \(1-e^{-m}-me^{-m}=e^{-m}(e^m-1-m)>0\)), and the
\(m_x\) are distinct for distinct \(x\), so equality holds if and only if
\(D\) has one element.
\end{proof}
\end{annexitem}

\begin{annexitem}{annex:p15-w02b-p15-057.dilution-average}{thm:p15-w02b-p15-057-dilution-average}{The infinite series and the figure 87.71 per cent}
\begin{proof}
\emph{Terms as functions of \(m\).} With \(m=m_x\), the terms are
\(\omega(m)=m^2e^{-m}/(1-e^{-m})\), \(\delta(m)=me^{-m}\) and
\(\nu(m)=e^{-m}-e^{-2m}\). Since \(m\leq e^m-1\), \(me^{-m}\leq1-e^{-m}\),
so \(\omega(m)\leq m\); also \(\delta(m)\leq m\) and, since
\(1-e^{-m}\leq m\), \(\nu(m)\leq m\). For \(m\geq1\),
\(1-e^{-m}\geq1-e^{-1}>\tfrac12\), so \(\omega(m)\leq2m^2e^{-m}\), and
\(\delta(m),\nu(m)\leq2m^2e^{-m}\) as well; and \(m^3e^{-m}\leq27e^{-3}<2\)
(its maximum is at \(m=3\)), so all three are at most \(4/m\) for
\(m\geq1\).

\emph{Convergence and periodicity.} The values \(m_x=e^\theta a^{-x}\) form
a two-sided geometric sequence with ratio \(a\). Those with \(m_x\leq1\)
are at most \(1,a^{-1},a^{-2},\ldots\) in decreasing order, so their terms
sum to at most \(a/(a-1)\); those with \(m_x>1\) are at least
\(1,a,a^2,\ldots\) in increasing order, so their terms sum to at most
\(4a/(a-1)\). Hence the three series converge for every \(\theta\).
Replacing \(\theta\) by \(\theta+\log a\) replaces \(m_x\) by \(m_{x-1}\),
which permutes the terms, so each series is periodic with period
\(\log a\).

\emph{Averages.} For a measurable \(g\geq0\) on \((0,\infty)\), put
\(L=\log a\) and \(h(u)=g(e^u)\). Then \(g(m_x)=h(\theta-xL)\), and by
Tonelli's theorem and the substitution \(u=\theta-xL\),
\[
 \frac1L\int_0^L\sum_{x\in\mathbb Z}h(\theta-xL)\,d\theta
 =\frac1L\sum_{x\in\mathbb Z}\int_{-xL}^{-xL+L}h(u)\,du
 =\frac1L\int_{\mathbb R}h(u)\,du=\frac1L\int_0^\infty g(m)\,\frac{dm}m .
\]
For \(\delta\): \(\int_0^\infty e^{-m}\,dm=1\). For \(\nu\), by Tonelli,
\(\int_0^\infty(e^{-m}-e^{-2m})m^{-1}\,dm
=\int_0^\infty\int_1^2e^{-um}\,du\,dm=\int_1^2u^{-1}\,du=\log2\). For
\(\omega\): for \(m>0\), \(e^{-m}/(1-e^{-m})=\sum_{k\geq1}e^{-km}\), so by
monotone convergence
\(\int_0^\infty me^{-m}(1-e^{-m})^{-1}\,dm=\sum_{k\geq1}\int_0^\infty
me^{-km}\,dm=\sum_{k\geq1}k^{-2}\).

\emph{The value \(\sum_kk^{-2}=\pi^2/6\).} For \(0<t<\pi/2\),
\(\sin t<t<\tan t\), so \(\cot^2t<t^{-2}<1+\cot^2t\). Fix \(N\geq1\) and put
\(t_k=k\pi/(2N+1)\), \(k=1,\ldots,N\). By de Moivre's formula,
\(\sin((2N+1)t)\) is the imaginary part of
\((\cos t+i\sin t)^{2N+1}=\sin^{2N+1}t\,(\cot t+i)^{2N+1}\), so
\[
 \sin((2N+1)t)=\sin^{2N+1}t\sum_{j=0}^{N}(-1)^j\binom{2N+1}{2j+1}
 \cot^{2N-2j}t .
\]
At \(t=t_k\) the left side vanishes and \(\sin t_k\neq0\), so the \(N\)
numbers \(\cot^2t_k\), which are distinct because \(\cot\) is strictly
decreasing and positive on \((0,\pi/2)\), are the \(N\) roots of the
polynomial \(\sum_{j=0}^N(-1)^j\binom{2N+1}{2j+1}y^{N-j}\) of degree \(N\).
Their sum is
\(\binom{2N+1}3/\binom{2N+1}1=N(2N-1)/3\). Summing the two inequalities at
\(t_k\),
\[
 \frac{N(2N-1)}3<\frac{(2N+1)^2}{\pi^2}\sum_{k=1}^Nk^{-2}
 <N+\frac{N(2N-1)}3 .
\]
Multiplying by \(\pi^2/(2N+1)^2\) and letting \(N\to\infty\), both bounds
tend to \(\pi^2/6\).

So \(\overline{\mathcal W}=\pi^2/(6\log a)\),
\(\overline{\mathcal D}=1/\log a\) and \(\overline{\mathcal V}=\log2/\log a\).

\emph{The case \(a=2\).} Here \(2m_x=m_{x-1}\), so
\(\nu(m_x)=e^{-m_x}-e^{-m_{x-1}}\), and the partial sum over
\(-N\leq x\leq N\) telescopes to \(e^{-m_N}-e^{-m_{-N-1}}\), which tends
to \(1-0=1\) as \(N\to\infty\).

\emph{Limits and the ratio.} For finite \(D\) increasing to \(\mathbb Z\),
the partial sums \(\mathcal W_D,\mathcal D_D,\mathcal V_D\) of convergent
series of positive terms tend to \(\mathcal W,\mathcal D,\mathcal V>0\), so
\(E_D(\theta)\to\mathcal D^2/(\mathcal V\mathcal W)\). Finally
\(\overline{\mathcal D}^{\,2}/(\overline{\mathcal V}\,\overline{\mathcal W})
=(\log a)^{-2}\big/\bigl(\log2\cdot\pi^2/(6(\log a)^2)\bigr)
=6/(\pi^2\log2)\), free of \(a\); with \(\pi^2=9.8696\ldots\) and
\(\log2=0.69315\ldots\) it is \(0.87705\ldots\).
\end{proof}
\end{annexitem}

\begin{annexitem}{annex:p15-w02b-p15-057.dilution-oscillation}{thm:p15-w02b-p15-057-dilution-oscillation}{The efficiency of the counting method at a given density}
\begin{proof}
\emph{Periodizations.} With \(m=e^u\), the terms of the three series are
\(h_{\mathcal D}(u)=e^ue^{-e^u}\), \(h_{\mathcal W}(u)=e^{2u}/(e^{e^u}-1)\)
and \(h_{\mathcal V}(u)=e^{-e^u}-e^{-2e^u}\), the functions
\(\delta,\omega,\nu\) of the proof of
\cref{thm:p15-w02b-p15-057-dilution-average} evaluated at \(e^u\); and
\(\mathcal D(\theta)=\sum_xh_{\mathcal D}(\theta-x\Lambda)
=\sum_jh_{\mathcal D}(\theta+j\Lambda)\), and likewise for \(\mathcal W\) and
\(\mathcal V\). That proof shows each term is at most \(m\), and at most
\(4/m\) when \(m\geq1\), hence at most \(\min(m,4/m)\). For
\(0\leq u\leq\Lambda\), \(e^{u+j\Lambda}\in[a^j,a^{j+1}]\), so
\(\sup_{0\leq u\leq\Lambda}h(u+j\Lambda)\leq\min(a^{j+1},4a^{-j})\), which
is summable over \(j\in\mathbb Z\). The functions are continuous, so
\cref{thm:p15-w02b-p15-057-periodization} applies with \(L=\Lambda\) once
the transforms are found and shown summable.

\emph{Transforms.} Substitute \(t=e^u\), \(du=dt/t\). For \(\mathcal D\),
\(\widehat h_{\mathcal D}(\omega)=\int_0^\infty t^{-i\omega}e^{-t}dt=
\Gamma(1-i\omega)\). For \(\mathcal W\),
\(\widehat h_{\mathcal W}(\omega)=\int_0^\infty t^{1-i\omega}(e^t-1)^{-1}dt
=\sum_{j\geq1}\int_0^\infty t^{1-i\omega}e^{-jt}dt\), the interchange being
justified by \(\sum_j\int_0^\infty te^{-jt}dt=\sum_jj^{-2}<\infty\); the
substitution \(t=s/j\) gives \(j^{-(2-i\omega)}\Gamma(2-i\omega)\), so
\(\widehat h_{\mathcal W}(\omega)=\Gamma(2-i\omega)\zeta(2-i\omega)\). For
\(\mathcal V\), \(e^{-t}-e^{-2t}=\int_1^2te^{-st}ds\), and by Fubini's
theorem (the absolute integrand has integral \(\log2\))
\(\widehat h_{\mathcal V}(\omega)=\int_1^2\int_0^\infty t^{-i\omega}e^{-st}dt\,ds
=\Gamma(1-i\omega)\int_1^2s^{-1+i\omega}ds\), which is
\(\Gamma(1-i\omega)(2^{i\omega}-1)/(i\omega)\) for \(\omega\neq0\) and
\(\log2\) for \(\omega=0\). Integration by parts gives
\(\Gamma(2-i\omega)=(1-i\omega)\Gamma(1-i\omega)\), the boundary terms
vanishing; \(|\zeta(2-i\omega)|\leq\sum_jj^{-2}=\zeta(2)\); and
\(\Gamma(1+i\omega)\) is the complex conjugate of \(\Gamma(1-i\omega)\), so
\(|\widehat h(\omega_{-k})|=|\widehat h(\omega_k)|\) in all three cases.

\emph{A bound for \(\Gamma(1-i\omega)\).} The function
\(h_{\mathcal D}(z)=e^z\exp(-e^z)\) is entire. For \(z=x+iy\) with
\(|y|\leq c<\pi/2\), \(|h_{\mathcal D}(z)|=e^x\exp(-e^x\cos y)\), so
\(\int|h_{\mathcal D}(x+iy)|dx=\int_0^\infty e^{-t\cos y}dt\leq1/\cos c\),
and \(|h_{\mathcal D}(x+iy)|\leq e^x\exp(-e^x\cos c)\to0\) as
\(|x|\to\infty\) uniformly in \(|y|\leq c\). The strip bound of
\cref{thm:p15-w02b-p15-057-periodization} gives
\(|\Gamma(1-i\omega)|\leq e^{-c|\omega|}/\cos c\) for every
\(c\in(0,\pi/2)\). With \(c=\pi/4\) this shows that
\(\sum_k(1+|\omega_k|)|\Gamma(1-i\omega_k)|<\infty\), so all three
transform sequences are summable and item~1 follows from
\cref{thm:p15-w02b-p15-057-periodization}, which states
\(\Lambda\mathcal D(\theta)=\sum_k\widehat h_{\mathcal D}(\omega_k)
e^{i\omega_k\theta}\) and the same for \(\mathcal W\) and \(\mathcal V\).
For \(|\omega|\geq1\) take \(c=\pi/2-1/|\omega|\), which lies in
\((0,\pi/2)\): \(\cos c=\sin(1/|\omega|)\geq\sin(1)/|\omega|\) by concavity
of \(\sin\) on \([0,1]\), and
\(e^{-c|\omega|}=e\,e^{-\pi|\omega|/2}\), giving the bound of item~3.

\emph{Item 2.} The constant terms are \(\Gamma(1)=1\),
\(\Gamma(2)\zeta(2)=\zeta(2)\) and \(\log2\). The remaining terms of
\(\Lambda\mathcal D\) have moduli \(\gamma_k\), those of \(\Lambda\mathcal W\)
at most \((1+\omega_k^2)^{1/2}\gamma_k\zeta(2)\), and those of
\(\Lambda\mathcal V\) equal to \(|2^{i\omega_k}-1|\gamma_k/|\omega_k|\); the
two indices \(\pm k\) contribute equally. If \(\log2/\Lambda=J\) is an
integer, \(2^{i\omega_k}=e^{2\pi ikJ}=1\) for every \(k\), so every
nonconstant term of \(\Lambda\mathcal V\) vanishes.

\emph{Item 3.} For \(k\geq1\), \(\omega_k=k\omega_1\geq1\) and
\(e^{-\pi\omega_k/2}=e^{-\pi^2k/\Lambda}=\rho^k\), so
\(\gamma_k\leq b\omega_1k\rho^k\) and \((1+\omega_k^2)^{1/2}\leq\sqrt2\,
\omega_1k\). Using \(\sum_{k\geq1}k\rho^k=\rho/(1-\rho)^2\),
\(\sum_{k\geq1}k^2\rho^k=\rho(1+\rho)/(1-\rho)^3\) and, with
\(|2^{i\omega_k}-1|\leq2\), \(\sum_{k\geq1}\rho^k=\rho/(1-\rho)\), the three
displayed bounds follow.

\emph{Item 4.} The efficiency is
\(\mathcal D^2/(\mathcal V\mathcal W)=(\Lambda\mathcal D)^2/
((\Lambda\mathcal V)(\Lambda\mathcal W))\). Under the stated conditions item~2
places the three factors in the positive intervals
\([1-\delta_{\mathcal D},1+\delta_{\mathcal D}]\),
\([\log2-\delta_{\mathcal V},\log2+\delta_{\mathcal V}]\) and
\([\zeta(2)-\delta_{\mathcal W},\zeta(2)+\delta_{\mathcal W}]\), which gives
both bounds except the \(1\); that one holds because \(E_D\leq1\) for every
finite \(D\) by \cref{thm:p15-w02b-p15-057-dilution} and \(E_D\) tends to
the efficiency by \cref{thm:p15-w02b-p15-057-dilution-average}.

\emph{Item 5.} For \(a=2\), \(\Lambda=0.6931472\), \(\log2/\Lambda=1\),
\(\omega_1=9.064720\), \(\pi^2/\Lambda=14.23883\),
\(\rho=6.5487\times10^{-7}\) and \(b=3.230393\); item~3 gives
\(\delta_{\mathcal D}\leq3.8353\times10^{-5}\) and
\(\delta_{\mathcal W}\leq8.0875\times10^{-4}\). Inserting the larger values
\(3.9\times10^{-5}\) and \(8.1\times10^{-4}\) in item~4 (whose bounds are
monotone in the \(\delta\)) gives the bounds \(0.876553\) and \(0.877554\), rounded outward. For \(a=3\),
\(\Lambda=1.0986123\), \(\omega_1=5.719202\), \(\pi^2/\Lambda=8.983701\) and
\(\rho=1.25438\times10^{-4}\); item~3 gives
\(\delta_{\mathcal D}\leq4.6362\times10^{-3}\),
\(\delta_{\mathcal W}\leq6.1697\times10^{-2}\) and
\(\delta_{\mathcal V}\leq1.6211\times10^{-3}\), and item~4 with
\(4.7\times10^{-3}\), \(6.2\times10^{-2}\) and \(1.7\times10^{-3}\) gives
the bounds \(0.83522\) and \(0.92226\), rounded outward.
\end{proof}
\end{annexitem}

\begin{annexitem}{annex:p15-w02b-p15-057.three-statements}{thm:p15-w02b-p15-057-three-statements}{The curvature interval is neither a coverage guarantee nor a posterior probability}
\begin{proof}
\emph{The maximizer.} For \(0<K<n\),
\(\ell_n'(p)=K/p-(n-K)/(1-p)\) and
\(\ell_n''(p)=-K/p^2-(n-K)/(1-p)^2<0\) on \((0,1)\), with
\(\ell_n\to-\infty\) at both ends, so \(\hat p=K/n\) is the unique
maximizer. For \(K=0\), \(\ell_n(p)=n\log(1-p)\) is strictly decreasing, so
\(\hat p=0\); \(K=n\) is symmetric.

\emph{Item 1.} For \(0<K<n\),
\(-\ell_n''(\hat p)=n^2/K+n^2/(n-K)=n^3/(K(n-K))=n/(\hat p(1-\hat p))\),
so \(v_O=\hat p(1-\hat p)/n\). The score of one observation is
\(X/p-(1-X)/(1-p)=(X-p)/(p(1-p))\), with variance \(1/(p(1-p))=I(p)\), so
\(1/(nI(\hat p))=\hat p(1-\hat p)/n=v_E\). For \(K=0\),
\(\ell_n''(p)=-n/(1-p)^2\), which is \(-n\) at \(p=0\), so \(v_O=1/n\),
while \(v_E=0\); \(K=n\) is symmetric.

\emph{Item 2.} \(W_E(0)=\{0\}\) and \(W_E(n)=\{1\}\). At \(p=0\),
\(K=0\) almost surely and \(0\in W_E(0)\), so the coverage is \(1\);
likewise at \(p=1\). For \(0<p<1\), \(p\in W_E(K)\) forces
\(0<K<n\), so the coverage is at most
\(P_p(0<K<n)=1-(1-p)^n-p^n\), which tends to \(0\) as \(p\downarrow0\).
Coverages are nonnegative, so the infimum is \(0\).

\emph{Item 3.} Fix \(p\in(0,1)\) and let \(G_n=\{0<K<n\}\), so
\(P_p(G_n)=1-(1-p)^n-p^n\to1\). On \(G_n\), \(W_E(K)=W_O(K)\) by item~1,
and since \(p\in[0,1]\) the intersection with \([0,1]\) does not matter:
\(p\in W(K)\) if and only if \(|S_n|\leq z\), where
\(S_n=\sqrt n(\hat p-p)/\sqrt{\hat p(1-\hat p)}\) on \(G_n\) and
\(S_n=0\) off \(G_n\). By \cref{thm:p15-clt-iid},
\(T_n=\sqrt n(\hat p-p)/\sqrt{p(1-p)}\Rightarrow N(0,1)\). By
\cref{thm:p15-iid-series-strong-law}, \(\hat p\to p\) almost surely, so
\(A_n=\sqrt{\hat p(1-\hat p)/(p(1-p))}\to1\) in probability, and
\(T_n/A_n\Rightarrow N(0,1)\) by \cref{thm:p15-clt-perturbation}; \(S_n\)
differs from \(T_n/A_n\) only off \(G_n\), so \(S_n-T_n/A_n\to0\) in
probability and \(S_n\Rightarrow N(0,1)\). Let \(F_n\) be the distribution
function of \(S_n\). As \(\Phi\) is continuous, \(F_n\to\Phi\) at every
point (\cref{def:p15-weak-law-convergence}). For \(\gamma>0\),
\(F_n(-z-\gamma)\leq P(S_n<-z)\leq F_n(-z)\), so every limit point of
\(P(S_n<-z)\) lies in \([\Phi(-z-\gamma),\Phi(-z)]\), and letting
\(\gamma\downarrow0\), \(P(S_n<-z)\to\Phi(-z)\). Hence
\(P(|S_n|\leq z)=F_n(z)-P(S_n<-z)\to\Phi(z)-\Phi(-z)=2\Phi(z)-1\), and both
coverages differ from \(P(|S_n|\leq z)\) by at most \(P_p(G_n^c)\to0\).

\emph{Item 4.} For \(n=4\): \(W_E(0)=\{0\}\), \(W_E(4)=\{1\}\),
\(W_E(1)=[0,\frac14+\frac{\sqrt3}8z]\) (the lower end
\(\frac14-\frac{\sqrt3}8z<0\) because \(z>\frac2{\sqrt3}\)),
\(W_E(2)=[\frac12-\frac z4,\frac12+\frac z4]\),
\(W_E(3)=[\frac34-\frac{\sqrt3}8z,1]\), and \(W_O(k)=W_E(k)\) for
\(k=1,2,3\), while \(W_O(0)=[0,\frac z2]\) and \(W_O(4)=[1-\frac z2,1]\).
With \(\frac54\leq z<\sqrt3\): at \(p=\frac58\), \(W(1)\) excludes \(p\)
since \(\frac{\sqrt3}8z<\frac38\); \(W(2)\) contains it since
\(\frac z4\geq\frac5{16}>\frac18\); \(W(3)\) contains it since
\(\frac{\sqrt3}8z\geq\frac{5\sqrt3}{32}>\frac18\); \(W_O(0)\) contains it
since \(\frac z2\geq\frac58\); \(W_O(4)\) contains it since
\(1-\frac z2\leq\frac38\). With
\(P_{5/8}(K=k)=\binom4k5^k3^{4-k}/4096\), equal to \(81,540,1350,1500,625\)
over \(4096\) for \(k=0,\ldots,4\), the coverage of \(W_E\) is
\((1350+1500)/4096=1425/2048\) and that of \(W_O\) is
\((81+1350+1500+625)/4096=889/1024\). At \(p=\frac12\), \(W(1)\) and
\(W(3)\) contain \(p\) since \(\frac{\sqrt3}8z\geq\frac14\), \(W(2)\)
contains it, \(W_E(0)\) and \(W_E(4)\) do not, and \(W_O(0)\), \(W_O(4)\)
do since \(\frac z2\geq\frac12\); so the coverages are \(1-\frac2{16}=\frac78\)
and \(1\).

\emph{Item 5.} The posterior law under the uniform prior is
\(\operatorname{beta}(k+1,n-k+1)\) by
\cref{thm:p15-confidence-credibility-boundary}. It has a density, so the
one-point sets \(W_E(0)\) and \(W_E(n)\) have posterior content \(0\).
For integer shapes, \(P(\operatorname{beta}(a,b)\leq u)\) is the
probability that a binomial variable with index \(a+b-1\) and probability
\(u\) is at least \(a\). With \(u_1=\frac14+\frac{\sqrt3}8z_{0.95}
=0.606121\) and \(a_2=\frac12-\frac{z_{0.95}}4=0.088787\) to six decimals,
the content of \(W_E(1)\) is
\(1-(1-u_1)^5-5u_1(1-u_1)^4=0.9176\), that of \(W_E(2)\) is
\(1-2\bigl(10a_2^3(1-a_2)^2+5a_2^4(1-a_2)+a_2^5\bigr)=0.9878\), and that of
\(W_E(3)\) equals that of \(W_E(1)\) by the symmetry \(p\mapsto1-p\), which
maps \(K=k\) to \(K=4-k\). The content of \(W_O(0)=[0,z_{0.95}/2]\) under
\(\operatorname{beta}(1,5)\) is \(1-(1-z_{0.95}/2)^5=0.9998\), and
\(W_O(4)\) is symmetric. The value \(z_{0.95}=1.6449\) to four decimals
lies in \([\frac54,\sqrt3)\).

\emph{Item 6.} For \(K=2\), \(\hat p=\frac12\) and
\(L(p)/L(\hat p)=\bigl(p(1-p)\bigr)^2/(1/4)^2=\bigl(1-(2p-1)^2\bigr)^2\).
For \(p\in[0,1]\), \(1-(2p-1)^2\in[0,1]\), and \(1-\frac{z^2}4>0\) when
\(z<2\), so the ratio is at least \((1-z^2/4)^2\) if and only if
\((2p-1)^2\leq z^2/4\), that is \(|p-\frac12|\leq\frac z4\). This interval
lies inside \([0,1]\) when \(z<2\) and is \(W_E(2)=W_O(2)\).
\end{proof}
\end{annexitem}

\begin{annexitem}{annex:p15-w02b-p15-057.sandwich}{thm:p15-w02b-p15-057-sandwich}{A root of the likelihood equation under a misspecified family}
\begin{proof}
\emph{Setting.} Since \(Q(A)=1\), almost surely every \(X_i\in A\), and
then by (M2) \(\ell_n\) is twice continuously differentiable on \(N\). The
functions \(s_\theta\) and \(\partial_\theta^2\log f_\theta\) are
measurable in \(x\) as pointwise limits of difference quotients of
measurable functions. Fix \(\eta_0>0\) with \(\eta_0<\delta\) and
\([\theta^*-\eta_0,\theta^*+\eta_0]\subseteq N\).

\emph{Uniform control of the curvature.} For \(0<\eta\leq\eta_0\) put
\[
 h_\eta(x)=\sup_{|\theta-\theta^*|\leq\eta}
 \bigl|\partial_\theta^2\log f_\theta(x)-\partial_\theta^2\log
 f_{\theta^*}(x)\bigr|\qquad(x\in A),
\]
measurable as a supremum over the rationals of the interval and its two
end points, by continuity. By (M4), \(0\leq h_\eta\leq2M\), and by (M2),
\(h_\eta(x)\to0\) as \(\eta\downarrow0\) for each \(x\in A\). Dominated
convergence gives \(E_Qh_\eta\to0\).

\emph{Expansion and item 2.} Let \(\Delta_n=\widehat\theta_n-\theta^*\).
For \(0<\eta\leq\eta_0\) let \(\Omega_n(\eta)\) be the event that all
\(X_i\in A\), \(|\Delta_n|<\eta\) and \(\ell_n'(\widehat\theta_n)=0\); by
hypothesis \(Q(\Omega_n(\eta))\to1\) for each fixed \(\eta\). Put
\(B_n=\int_0^1\ell_n''(\theta^*+t\Delta_n)\,dt\) on \(\Omega_n(\eta_0)\)
and \(B_n=\ell_n''(\theta^*)\) elsewhere, measurable as a limit of
Riemann sums of a continuous integrand. On \(\Omega_n(\eta_0)\) the
fundamental theorem of calculus gives
\(0=\ell_n'(\widehat\theta_n)=\ell_n'(\theta^*)+\Delta_nB_n\). On
\(\Omega_n(\eta)\),
\[
 \Bigl|\frac{B_n}{n}-\frac{\ell_n''(\theta^*)}{n}\Bigr|
 \leq\frac1n\sum_{i\leq n}h_\eta(X_i),\qquad
 \Bigl|\frac{\ell_n''(\widehat\theta_n)}{n}-\frac{\ell_n''(\theta^*)}{n}\Bigr|
 \leq\frac1n\sum_{i\leq n}h_\eta(X_i),
\]
and the right side converges almost surely to \(E_Qh_\eta\) by
\cref{thm:p15-iid-series-strong-law}. By the same strong law applied to the
integrable variables \(\partial_\theta^2\log f_{\theta^*}(X_i)\),
\(\ell_n''(\theta^*)/n\to-J\) almost surely. Given \(\varepsilon>0\),
choose \(\eta\) with \(E_Qh_\eta<\varepsilon\); then
\(Q(|B_n/n+J|>2\varepsilon)\to0\) and
\(Q(|\ell_n''(\widehat\theta_n)/n+J|>2\varepsilon)\to0\). Hence
\(-B_n/n\to J\) and \(-\ell_n''(\widehat\theta_n)/n\to J\) in probability.
Write \(a_n=-\ell_n''(\widehat\theta_n)/n\). On \(\{|a_n-J|\leq J/2\}\),
whose probability tends to one, \(a_n\geq J/2\) and
\(|1/a_n-1/J|=|a_n-J|/(a_nJ)\leq2|a_n-J|/J^2\); so
\(n\cdot\bigl(-1/\ell_n''(\widehat\theta_n)\bigr)=1/a_n\to1/J\) in
probability. Finally \(1/J=B/J^2\) if and only if \(J=B\), as \(J>0\).

\emph{Item 1.} By (M3), the variables \(s_{\theta^*}(X_i)\) are
independent identically distributed with mean zero and variance
\(B\in(0,\infty)\), so \(\ell_n'(\theta^*)/\sqrt n\Rightarrow N(0,B)\) by
\cref{thm:p15-clt-iid}. Put \(Y_n=(\ell_n'(\theta^*)/\sqrt n)/(-B_n/n)\),
with the value zero on \(\{B_n=0\}\). Since \(-B_n/n\to J\neq0\) in
probability, \cref{thm:p15-clt-perturbation} gives
\(Y_n\Rightarrow N(0,B)/J=N(0,B/J^2)\). On
\(\Omega_n(\eta_0)\cap\{B_n\neq0\}\), \(\sqrt n\,\Delta_n=Y_n\), and the
probability of this event tends to one. Hence
\(\sqrt n\,\Delta_n-Y_n\to0\) in probability, and the same lemma gives
\(\sqrt n\,\Delta_n\Rightarrow N(0,B/J^2)\).

\emph{Item 3.} For \(x\in A\) and \(|\theta-\theta^*|\leq\eta_0\), the
fundamental theorem of calculus and (M4) give
\(|s_\theta(x)-s_{\theta^*}(x)|\leq M(x)|\theta-\theta^*|\), hence
\[
 |s_\theta(x)^2-s_{\theta^*}(x)^2|
 =|s_\theta(x)-s_{\theta^*}(x)|\,|s_\theta(x)+s_{\theta^*}(x)|
 \leq M(x)|\theta-\theta^*|\bigl(2|s_{\theta^*}(x)|+M(x)|\theta-\theta^*|\bigr).
\]
On \(\Omega_n(\eta_0)\), with the Cauchy--Schwarz inequality for the
averages,
\[
 \Bigl|\widehat B_n-\frac1n\sum_{i\leq n}s_{\theta^*}(X_i)^2\Bigr|
 \leq2|\Delta_n|\Bigl(\frac1n\sum_{i\leq n}M(X_i)^2\Bigr)^{1/2}
 \Bigl(\frac1n\sum_{i\leq n}s_{\theta^*}(X_i)^2\Bigr)^{1/2}
 +\Delta_n^2\,\frac1n\sum_{i\leq n}M(X_i)^2 .
\]
By \cref{thm:p15-iid-series-strong-law}, the two averages converge almost
surely to \(E_QM(X)^2<\infty\) and to \(B\). Since \(\Delta_n\to0\) in
probability, the right side tends to zero in probability, and so
\(\widehat B_n\to B\) in probability. By item~2,
\((\ell_n''(\widehat\theta_n)/n)^2=a_n^2\to J^2>0\) in probability, and
on \(\{a_n\geq J/2\}\),
\[
 \Bigl|\frac{\widehat B_n}{a_n^2}-\frac B{J^2}\Bigr|
 \leq\frac{|\widehat B_n-B|}{a_n^2}+B\,\frac{|J^2-a_n^2|}{a_n^2J^2}
 \leq\frac{4|\widehat B_n-B|}{J^2}+\frac{4B\,|J^2-a_n^2|}{J^4},
\]
which tends to zero in probability.

\emph{Item 4.} If \(Q=P_{\theta^*}\) and the model is regular at
\(\theta^*\), then (R1) supplies (M1) and (M2) is (R2). The first step of
the proof of \cref{thm:p15-w02b-p15-057-asymptotic} shows from (R3) that
\(E_{\theta^*}s_{\theta^*}=0\) and
\(E_{\theta^*}\partial_\theta^2\log f_{\theta^*}=-I(\theta^*)\); with
(R4) and (R5) this gives (M3) and (M4) with \(B=J=I(\theta^*)\). Then
\(B/J^2=I(\theta^*)^{-1}\) and \(1/J=I(\theta^*)^{-1}\), the conclusions
of that theorem.
\end{proof}
\end{annexitem}

\begin{annexitem}{annex:p15-w02b-p15-057.poisson-misspecified}{thm:p15-w02b-p15-057-poisson-misspecified}{Counts fitted by a Poisson family}
\begin{proof}
\emph{Item 1.} For an integer \(x\geq1\),
\(0\leq\log x!\leq x\log x\leq x(x-1)\leq x^2\), using \(\log x\leq x-1\);
and \(\log0!=0\). Since \(E_QX^2=\sigma^2+\lambda^{*2}<\infty\),
\(E_Q\log X!\in[0,\infty)\). For \(\lambda>0\),
\(\log f_\lambda(x)=-\lambda+x\log\lambda-\log x!\) and
\(|\log f_\lambda(X)|\leq\lambda+X|\log\lambda|+\log X!\), which is
integrable; so \(\Lambda_Q(\lambda)=-\lambda+\lambda^*\log\lambda-E_Q\log
X!\) is finite. Its derivative \(-1+\lambda^*/\lambda\) is positive on
\((0,\lambda^*)\) and negative on \((\lambda^*,\infty)\), so
\(\Lambda_Q\) is strictly increasing then strictly decreasing, and
\(\lambda^*\) is its unique maximizer.

\emph{Item 2.} \(\ell_n(\lambda)=-n\lambda+n\overline X_n\log\lambda
-\sum_i\log X_i!\). If \(\overline X_n>0\), then
\(\ell_n'(\lambda)=-n+n\overline X_n/\lambda\) is positive below
\(\overline X_n\) and negative above it, so \(\overline X_n\) is the unique
maximizer on \((0,\infty)\). If \(\overline X_n=0\), then
\(\ell_n(\lambda)=-n\lambda\) is strictly decreasing on \((0,\infty)\) and
has no maximizer there. Because \(\lambda^*>0\), \(Q(\{0\})<1\), so
\(Q(\overline X_n>0)=1-Q(\{0\})^n\to1\). The variables \(X_i\) are
integrable, and \(\overline X_n\to\lambda^*\) almost surely by
\cref{thm:p15-iid-series-strong-law}.

\emph{Item 3.} The convergence \(\sqrt n(\overline X_n-\lambda^*)
\Rightarrow N(0,\sigma^2)\) is \cref{thm:p15-clt-iid}. For the
hypotheses: \(A=\{0,1,2,\ldots\}\) and \(f_\lambda(x)>0\) for all
\(\lambda>0\), which is (M1); \(\log f_\lambda(x)\) is infinitely
differentiable in \(\lambda\), with \(s_\lambda(x)=-1+x/\lambda\) and
\(\partial_\lambda^2\log f_\lambda(x)=-x/\lambda^2\), which is (M2).
\(E_Q|s_{\lambda^*}(X)|\leq1+E_QX/\lambda^*=2\),
\(E_Qs_{\lambda^*}(X)=-1+\lambda^*/\lambda^*=0\), and
\(B=E_Q(X-\lambda^*)^2/\lambda^{*2}=\sigma^2/\lambda^{*2}\in(0,\infty)\),
which is (M3). With \(\delta=\lambda^*/2\) and \(N=(0,\infty)\), for
\(|\lambda-\lambda^*|<\delta\) we have \(\lambda>\lambda^*/2\) and
\(|\partial_\lambda^2\log f_\lambda(x)|=x/\lambda^2\leq4x/\lambda^{*2}
=M(x)\), with \(E_QM(X)=4/\lambda^*\) and
\(E_QM(X)^2=16E_QX^2/\lambda^{*4}\), both finite; and
\(J=E_QX/\lambda^{*2}=1/\lambda^*>0\). This is (M4), and
\(B/J^2=\sigma^2\). Define \(\widehat\lambda_n=\overline X_n\) on
\(\{\overline X_n>0\}\) and \(\widehat\lambda_n=1\) elsewhere; by item~2,
\(\widehat\lambda_n\to\lambda^*\) almost surely and
\(Q(\ell_n'(\widehat\lambda_n)=0)\geq Q(\overline X_n>0)\to1\), so
\cref{thm:p15-w02b-p15-057-sandwich} applies, and its item~1 agrees with
the direct central limit statement.

\emph{Item 4.} On \(\{\overline X_n>0\}\),
\(\ell_n''(\lambda)=-n\overline X_n/\lambda^2\), so
\(-\ell_n''(\overline X_n)=n/\overline X_n\) and
\(-1/\ell_n''(\overline X_n)=\overline X_n/n\); \(n\) times this is
\(\overline X_n\to\lambda^*\) almost surely. Also
\[
 \widehat B_n=\frac1n\sum_{i\leq n}\Bigl(-1+\frac{X_i}{\overline X_n}\Bigr)^2
 =\frac{S_n^2}{\overline X_n^2},\qquad
 \Bigl(\frac{\ell_n''(\overline X_n)}n\Bigr)^2=\frac1{\overline X_n^2},
\]
so the corrected estimate is \(S_n^2\). As
\(S_n^2=n^{-1}\sum_iX_i^2-\overline X_n^2\) and \(X_i^2\) is integrable,
the strong law gives \(S_n^2\to E_QX^2-\lambda^{*2}=\sigma^2\) almost
surely.

\emph{Item 5.} Let \(G_n=\{\overline X_n>0\}\). Off \(G_n\) the curvature
interval is \(\{0\}\), which does not contain \(\lambda^*\). Put
\(U_n=\sqrt n(\overline X_n-\lambda^*)/\sigma\), so
\(U_n\Rightarrow N(0,1)\) by \cref{thm:p15-clt-iid}, and
\(A_n=\sqrt{\overline X_n}/\sigma\), which tends almost surely, hence in
probability, to \(c=\sqrt{\lambda^*}/\sigma>0\). By
\cref{thm:p15-clt-perturbation}, \(T_n=U_n/A_n\) (with the value zero on
\(\{A_n=0\}=G_n^c\)) satisfies \(T_n\Rightarrow N(0,1)/c\), the law
\(N(0,\sigma^2/\lambda^*)\), whose distribution function \(F\) is
continuous with \(F(t)=\Phi(t\sqrt{\lambda^*}/\sigma)\). On \(G_n\),
\(\lambda^*\) lies in the curvature interval if and only if
\(|T_n|\leq z\); so the coverage differs from \(Q(|T_n|\leq z)\) by at
most \(Q(G_n^c)\to0\). Let \(F_n\) be the distribution function of
\(T_n\). By \cref{def:p15-weak-law-convergence}, \(F_n(t)\to F(t)\) for
every \(t\). For \(\gamma>0\),
\(F_n(-z-\gamma)\leq Q(T_n<-z)\leq F_n(-z)\), so every limit point of
\(Q(T_n<-z)\) lies in \([F(-z-\gamma),F(-z)]\); letting
\(\gamma\downarrow0\) and using continuity of \(F\),
\(Q(T_n<-z)\to F(-z)\). Hence
\(Q(|T_n|\leq z)=F_n(z)-Q(T_n<-z)\to F(z)-F(-z)
=2\Phi(z\sqrt{\lambda^*}/\sigma)-1\). For the second interval let
\(H_n=\{S_n>0\}\); since \(S_n^2\to\sigma^2>0\) almost surely,
\(Q(H_n^c)\to0\). Put \(A_n'=S_n/\sigma\) on \(H_n\) and \(A_n'=1\)
elsewhere, so \(A_n'\to1\) in probability and \(T_n'=U_n/A_n'\Rightarrow
N(0,1)\) by \cref{thm:p15-clt-perturbation}. On \(H_n\), \(\lambda^*\)
lies in \(\overline X_n\pm zS_n/\sqrt n\) if and only if \(|T_n'|\leq z\),
and the argument just given, with \(\Phi\) in place of \(F\), shows
\(Q(|T_n'|\leq z)\to2\Phi(z)-1\); the coverage differs from this by at
most \(Q(H_n^c)\to0\).

\emph{Item 6.} The exponential function is continuous, so
\(e^{-\overline X_n}\to e^{-\lambda^*}\) almost surely by item~2.

\emph{The geometric example.} For \(|x|<1\),
\(\sum_{k\geq0}kx^k=x/(1-x)^2\) and \(\sum_{k\geq0}k^2x^k=x(1+x)/(1-x)^3\),
by differentiating the geometric series termwise inside its radius of
convergence. At \(x=\frac12\) these are \(2\) and \(6\), so under
\(Q(\{k\})=2^{-k-1}\), \(E_QX=\frac12\cdot2=1\) and
\(E_QX^2=\frac12\cdot6=3\); hence \(\lambda^*=1\) and \(\sigma^2=2\), and
\(Q(\{0\})=\frac12\). Item~5 gives the limits \(2\Phi(z/\sqrt2)-1\) and
\(2\Phi(z)-1\); at \(z=1.96\) these are \(0.8342\) and \(0.9500\) to four
decimals, and \(e^{-1}=0.3679\) to four decimals.
\end{proof}
\end{annexitem}
